% Special Continuous Distributions — Pure Mathematics with Statistics Upper Sixth — Cours signé OnBuch+
% Official MINESEC syllabus. Compiler avec : tectonic PMS14-special-continuous-distributions.tex
\newcommand{\DOCMATIERE}{Pure Mathematics with Statistics}
\newcommand{\DOCNIVEAU}{Upper Sixth}
\newcommand{\DOCMODULE}{Upper Sixth — Pure mathematics with statistics}
\newcommand{\DOCLECON}{14}
\newcommand{\DOCTITRE}{Special Continuous Distributions}
% OnBuch+ — préambule « cours signés » ENRICHI (compilé avec tectonic / XeLaTeX)
% Système visuel « Pop » d'OnBuch+ : fond crème (jamais blanc), bordures encre
% épaisses, ombres « dures » décalées, titres Archivo Black en capitales.
% Version MATHÉMATIQUES : graphes (pgfplots/tikz), boîtes pédagogiques
% (définition, propriété, méthode, attention, le savais-tu, exercice résolu,
% exercice, TP, bilan), page de garde et signature OnBuch+ ancrée en bas.
%
% Macros attendues dans main.tex AVANT \input{preamble} :
%   \DOCMATIERE  \DOCNIVEAU  \DOCMODULE  \DOCLECON (numéro)  \DOCTITRE
\documentclass[11pt]{article}

\usepackage{fontspec}
\usepackage[english]{babel}
\usepackage[a4paper,margin=2cm,top=3.4cm,bottom=2.8cm,headheight=1.4cm,footskip=1.3cm]{geometry}
\usepackage{amsmath,amssymb,amsfonts}
\usepackage{mathtools} % \xmapsto, \coloneqq... souvent utilisés par les modèles
\usepackage{cancel} % simplifications barrées (bilans de forces)
% \abs{z} (module, valeur absolue) et \norm{u} (norme), utilisés par les modèles
\providecommand{\abs}{}\let\abs\relax\DeclarePairedDelimiter{\abs}{\lvert}{\rvert}
\providecommand{\norm}{}\let\norm\relax\DeclarePairedDelimiter{\norm}{\lVert}{\rVert}
\usepackage[table]{xcolor} % [table] : fournit aussi \rowcolors (alternance de lignes)
\usepackage{pagecolor}
\usepackage{tikz}
\usetikzlibrary{arrows.meta,positioning,calc,shapes.geometric,shapes.misc,shapes.arrows,decorations.pathmorphing,decorations.pathreplacing,decorations.markings,patterns,shadows,mindmap,trees,fit,backgrounds,matrix,intersections,angles,quotes}
\usepackage{pgfplots}
\usepackage[version=4]{mhchem} % équations nucléaires \ce{^{235}_{92}U}
\usepackage[european,siunitx]{circuitikz} % schémas électriques
\pgfplotsset{compat=1.18}
\usepgfplotslibrary{fillbetween,groupplots} % aires entre courbes (calcul intégral)
\usepackage{enumitem}
\usepackage{fancyhdr}
% [most] charge toutes les bibliothèques tcolorbox (listings, minted, raster,
% documentation...) et gonfle inutilement la mémoire TeX sur un document long
% (~300 boîtes) : on ne charge que ce qui est réellement utilisé.
\usepackage[skins,breakable]{tcolorbox}
\usepackage{graphicx}
\usepackage{colortbl}
\usepackage{booktabs}
\usepackage{tabularx}
\usepackage{diagbox} % case d'en-tête diagonale des tableaux à double entrée
% Colonnes extensibles centrée / gauche / droite (C, L, R), souvent utilisées par les modèles
\newcolumntype{C}{>{\centering\arraybackslash}X}
\newcolumntype{L}{>{\raggedright\arraybackslash}X}
\newcolumntype{R}{>{\raggedleft\arraybackslash}X}
\usepackage{array}
\usepackage{multicol}
\usepackage{siunitx}
% parse-numbers=false : accepte aussi \SI{2,5 \times 10^{-3}}{...}
\sisetup{locale=UK,output-decimal-marker={.},group-minimum-digits=4,parse-numbers=false}
\DeclareSIUnit\ohm{\ensuremath{\symup{\Omega}}} % U+2126 absent de la police
\DeclareSIUnit\division{div} % réglages d'oscilloscope (ms/div, V/div)
\DeclareSIUnit\tr{tr} % tours (vitesse de rotation, ex. \SI{1500}{\tr\per\minute})
\usepackage{qrcode}
% \degree : commande fréquemment utilisée par les modèles pour un angle ou une
% température en dehors de \SI{}{} (non fournie par les paquets chargés).
\providecommand{\degree}{^{\circ}}
% \qed (fin de démonstration), utilisé par les modèles sans amsthm
\providecommand{\qed}{\hfill$\square$}
% \barre : double barre des valeurs interdites dans un tableau de variations (tkz-tab)
\providecommand{\barre}{\ensuremath{\|}}
% environnement proof (amsthm non chargé) utilisé par les modèles
\ifdefined\proof\else\newenvironment{proof}[1][Proof]{\par\noindent\textit{#1.}\ }{\qed\par}\fi
\usepackage{lastpage}
\usepackage{hyperref}
\hypersetup{hidelinks}

% ============================================================
% Palette « Pop » OnBuch+ — mêmes hex que la classe P (plus_ui.dart)
% ============================================================
\definecolor{poporange}{HTML}{F59321}
\definecolor{popdark}{HTML}{E07A0C}
\definecolor{popcream}{HTML}{FFF6E8}
\definecolor{popink}{HTML}{1C1714}
\definecolor{poppurple}{HTML}{7A5AE0}
\definecolor{popgreen}{HTML}{2E9E6A}
\definecolor{poppink}{HTML}{E0457B}
\definecolor{popblue}{HTML}{2F7DE1}
\definecolor{popgold}{HTML}{C98A12}
\definecolor{popmuted}{HTML}{8A7F76}
\definecolor{popline}{HTML}{EADFD2}
\definecolor{popgray}{HTML}{8A7F76}
\colorlet{popgrey}{popgray}
% Alias tolérants (le modèle nomme parfois les couleurs différemment)
\colorlet{popwhite}{white}
\colorlet{popblack}{popink}
\colorlet{popred}{poppink}
\colorlet{popyellow}{popgold}
% Teintes claires pour fonds de tableaux / zones
\colorlet{popblueL}{popblue!10!white}
\colorlet{poporangeL}{poporange!14!white}
\colorlet{popgreenL}{popgreen!12!white}
\colorlet{poppurpleL}{poppurple!12!white}
\colorlet{poppinkL}{poppink!12!white}
\colorlet{popgoldL}{popgold!14!white}

\pagecolor{popcream}

% ============================================================
% Typographie
% ============================================================
\defaultfontfeatures{Ligatures=TeX}
\setmainfont{PlusJakartaSans-Medium.ttf}[Path=../../../fonts/,BoldFont=PlusJakartaSans-Bold.ttf]
% Maths en sans-serif (Fira Math) avec chiffres et lettres droites de Plus
% Jakarta, pour que les formules se fondent dans le texte.
\usepackage[math-style=ISO,bold-style=ISO]{unicode-math}
\setmathfont{FiraMath-Regular.otf}[Path=../../../fonts/]
\setmathfont{PlusJakartaSans-Medium.ttf}[Path=../../../fonts/,range={up/num,up/{Latin,latin},\mathup}]
% (icomma retiré : décimales avec point en anglais)
% Caractères mathématiques Unicode absents des polices de texte (Plus Jakarta,
% Archivo Black) : rendus par leur équivalent mathématique, en texte comme en
% formule — sinon ils s'affichent en carré vide (ex. « Arithmétique dans ℤ »).
\usepackage{newunicodechar}
\newunicodechar{ℝ}{\ensuremath{\mathbb{R}}}
\newunicodechar{ℤ}{\ensuremath{\mathbb{Z}}}
\newunicodechar{ℂ}{\ensuremath{\mathbb{C}}}
\newunicodechar{ℕ}{\ensuremath{\mathbb{N}}}
\newunicodechar{ℚ}{\ensuremath{\mathbb{Q}}}
\newunicodechar{𝔻}{\ensuremath{\mathbb{D}}}
\newunicodechar{α}{\ensuremath{\alpha}}
\newunicodechar{β}{\ensuremath{\beta}}
\newunicodechar{γ}{\ensuremath{\gamma}}
\newunicodechar{δ}{\ensuremath{\delta}}
\newunicodechar{ε}{\ensuremath{\varepsilon}}
\newunicodechar{θ}{\ensuremath{\theta}}
\newunicodechar{λ}{\ensuremath{\lambda}}
\newunicodechar{μ}{\ensuremath{\mu}}
\newunicodechar{σ}{\ensuremath{\sigma}}
\newunicodechar{τ}{\ensuremath{\tau}}
\newunicodechar{φ}{\ensuremath{\varphi}}
\newunicodechar{ψ}{\ensuremath{\psi}}
\newunicodechar{ω}{\ensuremath{\omega}}
\newunicodechar{Σ}{\ensuremath{\Sigma}}
% majuscules produites par \MakeUppercase dans les titres (α-aminés -> Α)
\newunicodechar{Α}{\ensuremath{\alpha}}
\newunicodechar{Β}{\ensuremath{\beta}}
\newunicodechar{Γ}{\ensuremath{\Gamma}}
\newunicodechar{Δ}{\ensuremath{\Delta}}
\newunicodechar{Ε}{\ensuremath{\varepsilon}}
\newunicodechar{Θ}{\ensuremath{\Theta}}
\newunicodechar{Λ}{\ensuremath{\Lambda}}
\newunicodechar{Μ}{\ensuremath{\mu}}
\newunicodechar{Τ}{\ensuremath{\tau}}
\newunicodechar{Φ}{\ensuremath{\Phi}}
\newunicodechar{Ψ}{\ensuremath{\Psi}}
\newunicodechar{Ω}{\ensuremath{\Omega}}
\newunicodechar{↦}{\ensuremath{\mapsto}}
\newunicodechar{∈}{\ensuremath{\in}}
\newunicodechar{∉}{\ensuremath{\notin}}
\newunicodechar{∧}{\ensuremath{\wedge}}
\newunicodechar{⇔}{\ensuremath{\Leftrightarrow}}
\newunicodechar{⟺}{\ensuremath{\Longleftrightarrow}}
\newunicodechar{⇒}{\ensuremath{\Rightarrow}}
\newunicodechar{⟹}{\ensuremath{\Longrightarrow}}
\newunicodechar{∘}{\ensuremath{\circ}}
\newunicodechar{∪}{\ensuremath{\cup}}
\newunicodechar{∩}{\ensuremath{\cap}}
\newunicodechar{≡}{\ensuremath{\equiv}}
\newunicodechar{⊂}{\ensuremath{\subset}}
\newunicodechar{⊥}{\ensuremath{\perp}}
\newunicodechar{∃}{\ensuremath{\exists}}
\newunicodechar{∀}{\ensuremath{\forall}}
\newunicodechar{∛}{\ensuremath{\sqrt[3]{\,}}}
\newunicodechar{∜}{\ensuremath{\sqrt[4]{\,}}}
\newunicodechar{⃗}{\ensuremath{{}^{\to}}}
\newunicodechar{ⁿ}{\ifmmode^{n}\else\textsuperscript{\ensuremath{n}}\fi}
\newunicodechar{ˣ}{\ifmmode^{x}\else\textsuperscript{\ensuremath{x}}\fi}
\newunicodechar{ᵇ}{\ifmmode^{b}\else\textsuperscript{\ensuremath{b}}\fi}
\newunicodechar{ʳ}{\ifmmode^{r}\else\textsuperscript{\ensuremath{r}}\fi}
\newunicodechar{ᵘ}{\ifmmode^{u}\else\textsuperscript{\ensuremath{u}}\fi}
\newunicodechar{ʸ}{\ifmmode^{y}\else\textsuperscript{\ensuremath{y}}\fi}
\newunicodechar{ᵃ}{\ifmmode^{a}\else\textsuperscript{\ensuremath{a}}\fi}
\newunicodechar{ᵉ}{\ifmmode^{e}\else\textsuperscript{\ensuremath{e}}\fi}
\newunicodechar{ᵏ}{\ifmmode^{k}\else\textsuperscript{\ensuremath{k}}\fi}
\newunicodechar{ᵖ}{\ifmmode^{p}\else\textsuperscript{\ensuremath{p}}\fi}
\newunicodechar{ᴺ}{\ifmmode^{N}\else\textsuperscript{\ensuremath{N}}\fi}
\newunicodechar{ᶜ}{\ifmmode^{c}\else\textsuperscript{\ensuremath{c}}\fi}
\newunicodechar{ʰ}{\ifmmode^{h}\else\textsuperscript{\ensuremath{h}}\fi}
\newunicodechar{ᵠ}{\ifmmode^{q}\else\textsuperscript{\ensuremath{q}}\fi}
\newunicodechar{ₙ}{\ifmmode_{n}\else\textsubscript{\ensuremath{n}}\fi}
\newunicodechar{ₐ}{\ifmmode_{a}\else\textsubscript{\ensuremath{a}}\fi}
\newunicodechar{ᵢ}{\ifmmode_{i}\else\textsubscript{\ensuremath{i}}\fi}
\newunicodechar{ᵣ}{\ifmmode_{r}\else\textsubscript{\ensuremath{r}}\fi}
\newunicodechar{ₖ}{\ifmmode_{k}\else\textsubscript{\ensuremath{k}}\fi}
\newunicodechar{ᵦ}{\ifmmode_{\beta}\else\textsubscript{\ensuremath{\beta}}\fi}
\newunicodechar{ₓ}{\ifmmode_{x}\else\textsubscript{\ensuremath{x}}\fi}
\newfontfamily\popdisplay{ArchivoBlack-Regular.ttf}[Path=../../../fonts/]
\newfontfamily\poplogo{SpaceGrotesk-Bold.ttf}[Path=../../../fonts/]
\newfontfamily\popsemi{PlusJakartaSans-SemiBold.ttf}[Path=../../../fonts/]
\newfontfamily\popxbold{PlusJakartaSans-ExtraBold.ttf}[Path=../../../fonts/]
\newfontfamily\popxboldls{PlusJakartaSans-ExtraBold.ttf}[Path=../../../fonts/,LetterSpace=3.0]

\setlist[enumerate]{leftmargin=1.7em,itemsep=5pt,topsep=4pt}
\setlist[itemize]{leftmargin=1.7em,itemsep=4pt,topsep=4pt,label={\color{poporange}\textbullet}}
\setlength{\parindent}{0pt}
\setlength{\parskip}{5pt}
\renewcommand{\arraystretch}{1.25}

% pgfplots : style Pop par défaut
\pgfplotsset{
  every axis/.append style={axis line style={popink,line width=1pt},tick style={popink},
    grid style={popline,line width=0.5pt},label style={font=\small},tick label style={font=\footnotesize}},
  popcurve/.style={poporange,line width=1.8pt,no markers,smooth},
}
% Même style disponible dans un \draw TikZ simple (hors pgfplots)
\tikzset{popcurve/.style={poporange,line width=1.8pt}}
\tikzset{popgrid/.style={popline,very thin}}
\colorlet{popcurve}{poporange} % le nom du style est parfois utilisé comme couleur

% ============================================================
% Pill / squiggle
% ============================================================
\newcommand{\poppill}[3]{%
  \tikz[baseline=-0.55ex]\node[draw=popink,line width=1.2pt,rounded corners=2.6pt,%
    fill=#2,inner xsep=6.5pt,inner ysep=3pt]{%
    \popxboldls\fontsize{7}{7}\selectfont\color{#3}\MakeUppercase{#1}};%
}
\newcommand{\popsquiggle}[1]{%
  \tikz[baseline=-0.4ex]{
    \draw[#1,line width=2.6pt,line cap=round]
      plot [smooth] coordinates {(0,0.09) (0.34,0.01) (0.68,0.21) (1.02,0.01) (1.36,0.10)};
  }%
}
% Mot-clé surligné (marqueur orange sous le texte)
\newcommand{\cle}[1]{{\popxbold\color{popdark}#1}}

% ============================================================
% En-tête / pied
% ============================================================
\pagestyle{fancy}
\fancyhf{}
\renewcommand{\headrulewidth}{0pt}
\renewcommand{\footrulewidth}{0pt}
\lhead{\raisebox{-0.15ex}{\poplogo\fontsize{13}{13}\selectfont\color{popink}OnBuch\color{poporange}+}}
\rhead{\poppill{\DOCMATIERE\ \textbullet\ \DOCNIVEAU\ \textbullet\ Chapter \DOCLECON}{white}{popink}}
\lfoot{\popxbold\fontsize{7.6}{7.6}\selectfont\color{popmuted}\MakeUppercase{Signed course OnBuch+}\ \textbullet\ {\popsemi\DOCTITRE}}
\rfoot{\tikz[baseline=-0.6ex]\node[rounded corners=8.5pt,fill=popink,minimum height=17pt,minimum width=17pt,inner xsep=5pt,inner sep=0]{\popxbold\fontsize{8}{8}\selectfont\color{popcream}\thepage\,/\,\pageref*{LastPage}};}

% ============================================================
% Titres
% ============================================================
\newcommand{\chaptitre}[2]{%
  \begin{tcolorbox}[enhanced,colback=poporange,colframe=popink,boxrule=2.6pt,arc=11pt,
      drop shadow=popink,left=15pt,right=15pt,top=12pt,bottom=11pt,before skip=3pt,after skip=18pt]
    \poppill{#2}{popink}{popcream}\par\vspace{9pt}
    {\raggedright\hyphenpenalty=10000\exhyphenpenalty=10000\popdisplay\fontsize{22}{24}\selectfont\color{popcream}\MakeUppercase{#1}\par}
  \end{tcolorbox}
}
% Section numérotée : pastille orange avec le numéro + titre Archivo
\newcounter{popsec}
\newcommand{\coursec}[1]{%
  \stepcounter{popsec}\setcounter{popsub}{0}%
  \par\vspace{16pt}\noindent
  \begin{minipage}{\linewidth}
  \tikz[baseline=-0.7ex]\node[draw=popink,line width=1.6pt,fill=poporange,rounded corners=4pt,minimum size=22pt,inner sep=0]{\popdisplay\fontsize{12}{12}\selectfont\color{popcream}\thepopsec};%
  \hspace{8pt}{\popdisplay\fontsize{15}{17}\selectfont\color{popink}\MakeUppercase{#1}}\par
  \vspace{4pt}\noindent{\color{poporange}\rule{36pt}{3.6pt}}
  \end{minipage}\par\nopagebreak\vspace{8pt}
}
\newcounter{popsub}
\newcommand{\courssub}[1]{%
  \stepcounter{popsub}%
  \par\vspace{9pt}\noindent{\popxbold\fontsize{11.5}{13}\selectfont\color{popink}{\color{poporange}\thepopsec.\thepopsub}\hspace{6pt}#1}\par\nopagebreak\vspace{4pt}
}

% ============================================================
% Boîtes pédagogiques — même recette : fond blanc, bordure épaisse colorée,
% ombre dure encre, badge plein. Titre optionnel [#1].
% ============================================================
\tcbset{popbase/.style={breakable,enhanced,colback=white,boxrule=2.3pt,arc=9pt,
  shadow={3pt}{-3pt}{0pt}{popink},left=14pt,right=14pt,top=11pt,bottom=11pt,before skip=11pt,after skip=15pt,
  fontupper=\popsemi\fontsize{10.6}{14.5}\selectfont\color{popink},
  fontlower=\popsemi\fontsize{10.6}{14.5}\selectfont\color{popink}}}
% Coche dessinée (le glyphe ✓ n'existe pas dans les polices du document)
\newcommand{\popcheck}[1]{\tikz[baseline=-0.5ex]\draw[#1,line width=1.6pt,line cap=round,line join=round](-0.13,0.0)--(-0.04,-0.09)--(0.13,0.1);}
\providecommand{\coche}{\popcheck{popgreen}}
\providecommand{\resultat}[1]{\par\smallskip\noindent\fcolorbox{popgreen}{popgreenL}{\parbox{\dimexpr\linewidth-2\fboxsep-2\fboxrule\relax}{\textbf{Result.} #1}}\par\smallskip}
\newcommand{\popboxhead}[3]{\poppill{#1}{#2}{white}\ifx\relax#3\relax\else\hspace{6pt}{\popxbold\fontsize{10.6}{12}\selectfont\color{popink}#3}\fi\par\vspace{7pt}}

\newtcolorbox{aretenir}[1][]{popbase,colframe=popgreen,before upper={\popboxhead{Key points}{popgreen}{#1}}}
\newtcolorbox{exemplebox}[1][]{popbase,colframe=poporange,before upper={\popboxhead{Exemple}{poporange}{#1}}}
\newtcolorbox{definition}[1][]{popbase,colframe=popblue,before upper={\popboxhead{Definition}{popblue}{#1}}}
\newtcolorbox{propriete}[1][]{popbase,colframe=poppurple,before upper={\popboxhead{Property}{poppurple}{#1}}}
\newtcolorbox{methode}[1][]{popbase,colframe=popgold,before upper={\popboxhead{Method}{popgold}{#1}}}
\newtcolorbox{attention}[1][]{popbase,colframe=poppink,before upper={\popboxhead{Warning}{poppink}{#1}}}
\newtcolorbox{experience}[1][]{popbase,colframe=popblue,colback=popblueL,before upper={\popboxhead{Experiment}{popblue}{#1}}}
% « Le savais-tu ? » : carte violette pleine (contexte, culture, Cameroun)
\newtcolorbox{savaistu}[1][]{popbase,colback=poppurple,colframe=popink,
  fontupper=\popsemi\fontsize{10.4}{14.2}\selectfont\color{white},
  before upper={\poppill{Did you know?}{white}{poppurple}\ifx\relax#1\relax\else\hspace{6pt}{\popxbold\color{white}#1}\fi\par\vspace{7pt}}}
% Exercice résolu : énoncé en haut, \tcblower puis la solution sur fond crème
\newcounter{popexo}\newcounter{popexr}
\newtcolorbox{exoresolu}[1][]{popbase,colframe=poporange,colbacklower=popcream,
  segmentation style={popink,line width=1.2pt,dashed},
  before upper={\stepcounter{popexr}\popboxhead{Worked example \thepopexr}{poporange}{#1}},
  before lower={\poppill{Solution}{popink}{popcream}\par\vspace{6pt}}}
% Exercice (non résolu en place) — la correction est dans la partie Corrigés
\newtcolorbox{exercice}[1][]{popbase,colframe=popink,
  before upper={\stepcounter{popexo}\popboxhead{Exercise \thepopexo}{popink}{#1}}}
% Corrigé
\newtcolorbox{corrige}[1][]{popbase,colframe=popgreen,colback=popgreenL,
  before upper={\popboxhead{Solution}{popgreen}{#1}}}
% Objectifs / prérequis / bilan
\newtcolorbox{objectifs}{popbase,colframe=popink,colback=poporangeL,
  before upper={\popboxhead{Chapter objectives}{popink}{}}}
\newtcolorbox{prerequis}{popbase,colframe=popink,colback=popblueL,
  before upper={\popboxhead{Prerequisites}{popblue}{}}}
\newtcolorbox{bilan}{popbase,colframe=popink,colback=white,boxrule=2.8pt,
  before upper={\popboxhead{Fiche bilan}{poporange}{L'essentiel à connaître pour l'examen}}}
% Figure encadrée avec légende
\newtcolorbox{popfigure}[1][]{enhanced,breakable=false,colback=white,colframe=popline,boxrule=1.4pt,arc=8pt,
  left=10pt,right=10pt,top=10pt,bottom=8pt,before skip=10pt,after skip=12pt,halign=center,
  lower separated=false,halign lower=center,
  fontlower=\popsemi\fontsize{9}{11.5}\selectfont\color{popmuted},
  #1}
\newcommand{\legende}[1]{\par\vspace{6pt}{\popsemi\fontsize{9}{11.5}\selectfont\color{popmuted}#1}}

% ============================================================
% Signature OnBuch+
% ============================================================
\newcommand{\poplogoinline}[1]{{\poplogo\fontsize{#1}{#1}\selectfont\color{popink}OnBuch\color{poporange}+}}
\newcommand{\signatureonbuch}{%
  \begin{tcolorbox}[enhanced,colback=popink,colframe=popink,boxrule=2.2pt,arc=10pt,
      drop shadow=poporange,left=14pt,right=14pt,top=10pt,bottom=10pt,before skip=0pt,after skip=0pt]
    \tikz[baseline=-0.7ex]\node[circle,fill=popgreen,draw=popcream,line width=1.3pt,minimum size=22pt,inner sep=0]{\popcheck{white}};%
    \hspace{9pt}\begin{minipage}[c]{0.62\linewidth}
      {\popxbold\fontsize{12}{13}\selectfont\color{popcream}Signed\ \textbullet\ {\poplogo OnBuch\color{poporange}+}}\par\vspace{2pt}
      {\raggedright\popsemi\fontsize{8}{10.5}\selectfont\color{popline}\MakeUppercase{OnBuch+ teaching team}\\\MakeUppercase{Follows the official MINESEC syllabus}\par}
    \end{minipage}\hfill
    {\poplogo\fontsize{22}{22}\selectfont\color{popcream}OnBuch\color{poporange}+}
  \end{tcolorbox}%
}

% ============================================================
% Page de garde
% #1 titre · #2 matière · #3 niveau · #4 module · #5 n° leçon · #6 durée ·
% #7 description / objectifs (texte ou itemize)
% ============================================================
\newcommand{\pagegarde}[7]{%
  \thispagestyle{empty}
  \newgeometry{a4paper,margin=1.7cm,top=1.6cm,bottom=1.4cm}%
  \begin{tikzpicture}[remember picture,overlay]
    \fill[poporange!18] ([xshift=-3.2cm,yshift=-3.6cm]current page.north east) circle (4.6cm);
    \fill[poppurple!14] ([xshift=2.4cm,yshift=7.6cm]current page.south west) circle (3.8cm);
  \end{tikzpicture}%
  {\poplogo\fontsize{30}{30}\selectfont\color{popink}OnBuch\color{poporange}+}\hfill
  \raisebox{0.6ex}{\poppill{Signed course \textbullet\ Premium}{popink}{popcream}}\par
  \vspace{1pt}\popsquiggle{poppurple}\par
  \vspace{8pt}
  \begin{tcolorbox}[enhanced,colback=poporange,colframe=popink,boxrule=2.8pt,arc=13pt,
      drop shadow=popink,left=18pt,right=18pt,top=12pt,bottom=13pt,after skip=10pt]
    \poppill{#2 \textbullet\ #3}{popink}{popcream}\hspace{5pt}\poppill{#4}{white}{popink}\par\vspace{8pt}
    {\popdisplay\fontsize{14}{14}\selectfont\color{popink}CHAPTER #5}\par\vspace{4pt}
    {\raggedright\hyphenpenalty=10000\exhyphenpenalty=10000\popdisplay\fontsize{25}{27}\selectfont\color{popcream}\MakeUppercase{#1}\par}
    \vspace{8pt}
    {\popxbold\fontsize{9.5}{9.5}\selectfont\color{popink}\MakeUppercase{Suggested time: #6}}
  \end{tcolorbox}
  \begin{tcolorbox}[enhanced,colback=white,colframe=popink,boxrule=2.2pt,arc=10pt,
      drop shadow=popink,left=16pt,right=16pt,top=10pt,bottom=10pt,after skip=8pt]
    \poppill{In this chapter}{poppurple}{white}\par\vspace{5pt}
    \popsemi\fontsize{9.3}{12.6}\selectfont\color{popink}#7
  \end{tcolorbox}
  \noindent\begin{minipage}[t]{0.487\linewidth}
    \begin{tcolorbox}[enhanced,colback=popgreen,colframe=popink,boxrule=2.2pt,arc=10pt,
        drop shadow=popink,left=11pt,right=11pt,top=10pt,bottom=10pt,height=3.1cm,valign=center]
      \tcbox[colback=white,colframe=popink,boxrule=1.3pt,arc=5pt,boxsep=0pt,nobeforeafter,
        left=3pt,right=3pt,top=3pt,bottom=3pt,valign=center]{\qrcode[height=1.9cm]{https://play.google.com/store/apps/details?id=com.luvvix.onbuch&hl=en}}%
      \hspace{8pt}\begin{minipage}[c]{0.5\linewidth}
        {\popdisplay\fontsize{11}{12.5}\selectfont\color{white}\MakeUppercase{Download OnBuch}}\par\vspace{4pt}
        {\raggedright\popsemi\fontsize{8.4}{11}\selectfont\color{white}Free \textbullet\ Google Play\\AI tutor, past papers, results\par}
      \end{minipage}
    \end{tcolorbox}
  \end{minipage}\hfill
  \begin{minipage}[t]{0.487\linewidth}
    \begin{tcolorbox}[enhanced,colback=poppurple,colframe=popink,boxrule=2.2pt,arc=10pt,
        drop shadow=popink,left=11pt,right=11pt,top=10pt,bottom=10pt,height=3.1cm,valign=center]
      \tikz[baseline=-0.7ex]\node[circle,fill=white,draw=popink,line width=1.3pt,minimum size=1.6cm,inner sep=0]{\popdisplay\fontsize{24}{24}\selectfont\color{poppurple}+};%
      \hspace{8pt}\begin{minipage}[c]{0.6\linewidth}
        {\popdisplay\fontsize{11}{12.5}\selectfont\color{white}\MakeUppercase{Go OnBuch+}}\par\vspace{4pt}
        {\raggedright\popsemi\fontsize{8.4}{11}\selectfont\color{white}All signed courses, unlimited AI tutor, PDF sheets — OnBuch+ tab in the app\par}
      \end{minipage}
    \end{tcolorbox}
  \end{minipage}
  \vspace{8pt}
  \popcontenu
  \vfill
  \signatureonbuch
  \restoregeometry
  \newpage
}

% Rangée « Ce cours contient » (page de garde)
\newcommand{\poptile}[3]{% #1 couleur #2 gros texte #3 légende
  \begin{tcolorbox}[enhanced,colback=#1,colframe=popink,boxrule=2pt,arc=9pt,shadow={2.5pt}{-2.5pt}{0pt}{popink},
      left=6pt,right=6pt,top=9pt,bottom=9pt,halign=center,valign=center,height=1.75cm,width=\linewidth,nobeforeafter]
    {\popdisplay\fontsize{12.5}{13}\selectfont\color{white}\MakeUppercase{#2}}\par\vspace{3pt}
    {\centering\popsemi\fontsize{7.8}{9.5}\selectfont\color{white}#3\par}
  \end{tcolorbox}}
\newcommand{\popcontenu}{%
  \par\noindent{\popxboldls\fontsize{7.5}{7.5}\selectfont\color{popmuted}THIS SIGNED COURSE CONTAINS}\par\vspace{7pt}
  \noindent\begin{minipage}[t]{0.235\linewidth}\poptile{popblue}{Course}{complete, illustrated, step by step}\end{minipage}\hfill
  \begin{minipage}[t]{0.235\linewidth}\poptile{poporange}{Methods}{and worked examples}\end{minipage}\hfill
  \begin{minipage}[t]{0.235\linewidth}\poptile{poppink}{Exercises}{exam style + solutions}\end{minipage}\hfill
  \begin{minipage}[t]{0.235\linewidth}\poptile{popgreen}{Summary sheet}{the essentials to remember}\end{minipage}\par}

% Sommaire
\newtcolorbox{sommaire}{enhanced,colback=white,colframe=popink,boxrule=2.2pt,arc=10pt,
  drop shadow=popink,left=18pt,right=18pt,top=13pt,bottom=13pt,before skip=6pt,after skip=20pt,
  before upper={\poppill{Contents}{poppurple}{white}\par\vspace{9pt}\popsemi\fontsize{10.6}{14.5}\selectfont\color{popink}}}

% Signature de fin de cours — poussée tout en bas de la dernière page
\newcommand{\findecours}{%
  \par\vfill\nopagebreak
  \begin{center}\popsquiggle{poporange}\end{center}\vspace{4pt}
  \signatureonbuch
}
\AtBeginDocument{\def\llbracket{\mathopen{[\mkern-3.5mu[}}\def\rrbracket{\mathclose{]\mkern-3.5mu]}}} % intervalles d'entiers (glyphe ⟦ absent de la police)
% Glyphes absents de Fira Math : redéfinis à partir de glyphes disponibles
\AtBeginDocument{%
  \def\setminus{\mathbin{\backslash}}\let\smallsetminus\setminus
  \def\vdots{\mathord{\vbox{\baselineskip3.2pt\lineskiplimit0pt\kern4pt\hbox{.}\hbox{.}\hbox{.}}}}%
  \def\ddots{\mathinner{\mkern1mu\raise6.5pt\hbox{.}\mkern2mu\raise3.6pt\hbox{.}\mkern2mu\raise0.7pt\hbox{.}\mkern1mu}}%
  \def\triangle{\mathord{\scalebox{0.9}{$\symup{\Delta}$}}}%
  \def\nabla{\mathord{\raisebox{\depth}{\scalebox{1}[-1]{$\symup{\Delta}$}}}}%
  \def\checkmark{\popcheck{popdark}}%
  \def\star{\mathbin{\ast}}\def\bigstar{\mathord{\ast}}%
  \def\diamond{\mathbin{\scalebox{0.6}{$\Diamond$}}}%
  \def\bigcup{\mathop{\mathchoice{\raisebox{-0.3em}{\scalebox{1.7}{$\cup$}}}{\scalebox{1.25}{$\cup$}}{\cup}{\cup}}\displaylimits}%
  \def\bigcap{\mathop{\mathchoice{\raisebox{-0.3em}{\scalebox{1.7}{$\cap$}}}{\scalebox{1.25}{$\cap$}}{\cap}{\cap}}\displaylimits}%
  \def\lesseqqgtr{\mathrel{\lessgtr}}\def\bigoplus{\mathop{\oplus}\displaylimits}%
}
\providecommand{\ssi}{if and only if} % abréviation fréquente
\AtBeginDocument{\providecommand{\wideparen}{\overparen}} % arc de cercle

% Fonctions usuelles que les modèles utilisent sans que amsmath les fournisse
\providecommand{\arsinh}{\operatorname{arsinh}}\providecommand{\arcosh}{\operatorname{arcosh}}
\providecommand{\artanh}{\operatorname{artanh}}\providecommand{\arsech}{\operatorname{arsech}}
\providecommand{\arcsch}{\operatorname{arcsch}}\providecommand{\arcoth}{\operatorname{arcoth}}
\providecommand{\arcsinh}{\operatorname{arsinh}}\providecommand{\arccosh}{\operatorname{arcosh}}
\providecommand{\arctanh}{\operatorname{artanh}}\providecommand{\sech}{\operatorname{sech}}
\providecommand{\csch}{\operatorname{csch}}\providecommand{\arcsec}{\operatorname{arcsec}}
\providecommand{\arccsc}{\operatorname{arccsc}}\providecommand{\arccot}{\operatorname{arccot}}
\providecommand{\sgn}{\operatorname{sgn}}\providecommand{\lcm}{\operatorname{lcm}}
\providecommand{\Var}{\operatorname{Var}}\providecommand{\Cov}{\operatorname{Cov}}

%% FIGFIX-BEGIN v4 — correctifs globaux des figures (docs/onbuch-generation/tools/figfix)
\usepackage{adjustbox}\usepackage{etoolbox}
% toute figure TikZ/pgfplots plus large que la ligne est réduite à la largeur disponible
\BeforeBeginEnvironment{tikzpicture}{\begin{adjustbox}{max width=\linewidth}}
\AfterEndEnvironment{tikzpicture}{\end{adjustbox}}
% légendes pgfplots lisibles par-dessus les courbes
\pgfplotsset{every axis/.append style={legend style={fill=white,fill opacity=0.92,text opacity=1,draw=black!15,inner sep=3pt}}}
% étiquettes d'arêtes (syntaxe edge["..."]) avec fond blanc
\tikzset{every edge quotes/.append style={fill=white,inner sep=1.2pt}}
%% FIGFIX-END
\begin{document}

\pagegarde{Special Continuous Distributions}{Pure Mathematics with Statistics}{Upper Sixth}{Upper Sixth — Pure mathematics with statistics}{14}{18 periods}{This chapter studies the three special continuous probability distributions required at GCE Advanced Level: the uniform (rectangular), the exponential and the Normal distributions. Students learn their probability density functions, expectations, variances, the use of standard Normal tables, continuity corrections, Normal approximations to the binomial and Poisson distributions, and combinations of independent Normal variables.
\vspace{4pt}
\begin{multicols}{2}\small
\begin{itemize}
\item By the end of this chapter I can state and use the probability density function, expectation and variance of a uniform (rectangular) distribution.
\item By the end of this chapter I can solve applied problems modelled by the uniform distribution.
\item By the end of this chapter I can derive and use the expectation and variance of an exponential distribution and apply it to waiting-time problems.
\item By the end of this chapter I can explain the memoryless property of the exponential distribution and its link with the Poisson distribution.
\item By the end of this chapter I can describe the Normal curve, standardise a Normal variable and read the standard Normal tables Φ(z).
\item By the end of this chapter I can use the symmetry properties of Φ(z) and find central 95\% and 99\% intervals.
\end{itemize}\end{multicols}}

\begin{sommaire}
\begin{multicols}{2}
\begin{enumerate}
\item The Uniform (Rectangular) Distribution
\item The Exponential Distribution
\item The Normal Distribution and Standardisation
\item Central Intervals and Working with Any Normal Distribution
\item Determining Unknown μ and/or σ
\item Continuity Correction and Normal Approximations
\item Combinations of Independent Normal Variables
\item Integration, Assessment and Remediation
\item Integration activity
\item Methods and worked examples
\item Exercises
\item Solutions to the exercises
\item Summary sheet
\end{enumerate}
\end{multicols}
\end{sommaire}

\begin{objectifs}
\begin{itemize}
\item By the end of this chapter I can state and use the probability density function, expectation and variance of a uniform (rectangular) distribution.
\item By the end of this chapter I can solve applied problems modelled by the uniform distribution.
\item By the end of this chapter I can derive and use the expectation and variance of an exponential distribution and apply it to waiting-time problems.
\item By the end of this chapter I can explain the memoryless property of the exponential distribution and its link with the Poisson distribution.
\item By the end of this chapter I can describe the Normal curve, standardise a Normal variable and read the standard Normal tables Φ(z).
\item By the end of this chapter I can use the symmetry properties of Φ(z) and find central 95\% and 99\% intervals.
\item By the end of this chapter I can de-standardise and determine unknown values of μ and/or σ from given probabilities.
\item By the end of this chapter I can apply continuity corrections and use the Normal approximation to the binomial and Poisson distributions.
\item By the end of this chapter I can find the distribution of sums and differences of two independent Normal variables.
\end{itemize}
\end{objectifs}

\begin{prerequis}
\begin{itemize}
\item Continuous random variables: probability density functions, P(a < X < b) as an area, cumulative distribution function (Lower Sixth / PMS on continuous random variables)
\item Expectation E(X) and variance Var(X) of continuous random variables, including E(aX + b) and Var(aX + b)
\item The Poisson distribution: conditions of validity, mean = variance = λ, calculation of P(X = r)
\item The binomial distribution B(n, p): mean np, variance npq, conditions of use
\item Basic integration of polynomials and exponential functions, and use of limits at infinity
\end{itemize}
\end{prerequis}

\newpage

\coursec{The Uniform (Rectangular) Distribution}

At the Mfoundi market in Yaoundé, a vendor notices that customers arrive at her stall roughly every 4 minutes, while the waiting time at the SONEL cash office across the road seems completely unpredictable — sometimes 1 minute, sometimes 40. At the GCE marking centre, examiners observe that candidates' marks in Mathematics cluster around a central value with fewer and fewer scripts at the extremes. Why do waiting times, lifetimes of phone batteries, and examination marks follow such different yet perfectly predictable patterns? This chapter reveals the three special continuous distributions — \cle{uniform}, \cle{exponential} and \cle{Normal} — that model these everyday Cameroonian realities and that dominate the GCE Advanced Level statistics paper.

We begin with the simplest of the three: the uniform distribution, in which \emph{every} value in an interval is equally likely.

\courssub{Introducing the Idea: A Concrete Situation}

\begin{experience}[The bus that comes "any time within 20 minutes"]
A student at the Carrefour Warda bus stop in Yaoundé is told that the next bus will arrive \emph{at a completely random instant} within the next 20 minutes: no moment is more likely than any other. Let $X$ be the waiting time, in minutes. Then $X$ takes values in $[0, 20]$ and:
\begin{itemize}
\item the probability of waiting between 5 and 10 minutes should be the same as waiting between 12 and 17 minutes (both intervals have length 5);
\item the probability of waiting more than 15 minutes should be $\tfrac{5}{20} = \tfrac14$.
\end{itemize}
This "equal chance everywhere on an interval" is exactly the spirit of the uniform distribution. Probabilities depend only on the \emph{length} of the interval considered, not on its position.
\end{experience}

\courssub{Definition and Probability Density Function}

\begin{definition}[Continuous uniform (rectangular) distribution]
A continuous random variable $X$ follows the \cle{uniform distribution} on the interval $[a, b]$ (with $a < b$), written $X \sim U(a, b)$, if its probability density function is
\[
f(x) = \begin{cases} \dfrac{1}{b - a} & \text{for } a \le x \le b, \\[6pt] 0 & \text{otherwise.} \end{cases}
\]
Because the graph of $f$ is a rectangle, this is also called the \cle{rectangular distribution}.
\end{definition}

Let us check that $f$ is a genuine pdf. It is clearly non-negative, and
\[
\int_{-\infty}^{+\infty} f(x)\,dx = \int_a^b \frac{1}{b-a}\,dx = \frac{b - a}{b - a} = 1.
\]
The total area under the curve is the area of a rectangle of width $b - a$ and height $\tfrac{1}{b-a}$, which is indeed $1$.

\begin{propriete}[Probabilities as areas of rectangles]
If $X \sim U(a, b)$ and $[c, d] \subseteq [a, b]$, then
\[
P(c < X < d) = \int_c^d \frac{1}{b-a}\,dx = \frac{d - c}{b - a}.
\]
The probability is the ratio of the length of the sub-interval to the length of the whole interval. (Proof not examinable; it follows directly from the definition of a pdf.)
\end{propriete}

\begin{attention}[Endpoints do not matter]
For a continuous variable, $P(X = c) = 0$, so $P(c < X < d)$, $P(c \le X \le d)$, $P(c \le X < d)$ are all equal. Never waste time in the exam worrying about strict versus non-strict inequalities for continuous distributions.
\end{attention}

\begin{popfigure}
\begin{center}
\begin{tikzpicture}
\begin{axis}[
  width=11cm, height=5.5cm,
  axis lines=middle,
  xlabel={$x$}, ylabel={$f(x)$},
  xmin=-0.5, xmax=5.5, ymin=0, ymax=0.3,
  xtick={0,1,3,4}, xticklabels={$a$,$c$,$d$,$b$},
  ytick={0.25}, yticklabels={$\frac{1}{b-a}$},
  clip=false
]
\addplot[fill=popblueL, draw=none] coordinates {(1,0) (3,0) (3,0.25) (1,0.25)} \closedcycle;
\addplot[very thick, popblue] coordinates {(0,0.25) (4,0.25)};
\addplot[very thick, popblue] coordinates {(-0.5,0) (0,0)};
\addplot[very thick, popblue] coordinates {(4,0) (5.5,0)};
\draw[dashed, popmuted] (axis cs:0,0) -- (axis cs:0,0.25);
\draw[dashed, popmuted] (axis cs:4,0) -- (axis cs:4,0.25);
\draw[dashed, popmuted] (axis cs:1,0) -- (axis cs:1,0.25);
\draw[dashed, popmuted] (axis cs:3,0) -- (axis cs:3,0.25);
\node[popdark] at (axis cs:2,0.14) {$P(c<X<d)$};
\end{axis}
\end{tikzpicture}
\end{center}
\legende{Pdf of $X \sim U(a,b)$: the probability $P(c<X<d)$ is the area of the shaded rectangle, $\frac{d-c}{b-a}$.}
\end{popfigure}

\courssub{Expectation and Variance}

\begin{propriete}[Mean and variance of $U(a,b)$ — proof examinable]
If $X \sim U(a, b)$, then
\[
E(X) = \frac{a + b}{2} \qquad \text{and} \qquad \operatorname{Var}(X) = \frac{(b - a)^2}{12}.
\]
\end{propriete}

\textbf{Derivation of $E(X)$ by integration.}
\[
E(X) = \int_a^b x \cdot \frac{1}{b-a}\,dx = \frac{1}{b-a}\left[\frac{x^2}{2}\right]_a^b = \frac{b^2 - a^2}{2(b-a)} = \frac{(b-a)(b+a)}{2(b-a)} = \frac{a+b}{2}.
\]

\textbf{Derivation of $E(X)$ by symmetry.} The pdf is symmetric about the midpoint of $[a,b]$, namely $\tfrac{a+b}{2}$. For any symmetric continuous distribution, the mean is the axis of symmetry, so $E(X) = \tfrac{a+b}{2}$ immediately. This argument is a useful check in the exam.

\textbf{Derivation of $\operatorname{Var}(X)$.} First compute $E(X^2)$:
\[
E(X^2) = \int_a^b x^2 \cdot \frac{1}{b-a}\,dx = \frac{1}{b-a}\left[\frac{x^3}{3}\right]_a^b = \frac{b^3 - a^3}{3(b-a)} = \frac{a^2 + ab + b^2}{3},
\]
using $b^3 - a^3 = (b-a)(a^2 + ab + b^2)$. Then
\begin{align*}
\operatorname{Var}(X) &= E(X^2) - [E(X)]^2 = \frac{a^2 + ab + b^2}{3} - \frac{(a+b)^2}{4} \\
&= \frac{4(a^2 + ab + b^2) - 3(a^2 + 2ab + b^2)}{12} = \frac{a^2 - 2ab + b^2}{12} = \frac{(b-a)^2}{12}.
\end{align*}

\begin{attention}[The most common GCE error]
The variance is $\operatorname{Var}(X) = \dfrac{(b-a)^2}{12}$, \textbf{not} $\dfrac{b-a}{12}$. The squaring is forgotten by a large number of candidates every year. Remember: a variance has the units of $X$ \emph{squared}, so $(b-a)$ must appear squared.
\end{attention}

\begin{aretenir}[Key facts about $X \sim U(a,b)$]
\begin{itemize}
\item Pdf: $f(x) = \dfrac{1}{b-a}$ on $[a,b]$, $0$ elsewhere.
\item $P(c < X < d) = \dfrac{d-c}{b-a}$ for $[c,d] \subseteq [a,b]$.
\item $E(X) = \dfrac{a+b}{2}$, \quad $\operatorname{Var}(X) = \dfrac{(b-a)^2}{12}$, \quad $\sigma = \dfrac{b-a}{\sqrt{12}} = \dfrac{b-a}{2\sqrt{3}}$.
\item Cdf: $F(x) = \dfrac{x-a}{b-a}$ for $a \le x \le b$.
\end{itemize}
\end{aretenir}

\courssub{The Cumulative Distribution Function}

\begin{propriete}[Cdf of the uniform distribution]
For $X \sim U(a,b)$,
\[
F(x) = P(X \le x) = \begin{cases} 0 & \text{if } x < a, \\[4pt] \dfrac{x - a}{b - a} & \text{if } a \le x \le b, \\[6pt] 1 & \text{if } x > b. \end{cases}
\]
\end{propriete}

\textbf{Justification.} For $a \le x \le b$, $F(x) = \displaystyle\int_a^x \frac{1}{b-a}\,dt = \frac{x-a}{b-a}$; outside $[a,b]$ the cdf takes its limiting values $0$ and $1$.

The cdf is convenient for probabilities of the form $P(X \le x)$ or $P(X > x) = 1 - F(x)$, and for finding percentiles: the median $m$ satisfies $F(m) = \tfrac12$, giving $m = \tfrac{a+b}{2}$, which coincides with the mean — as expected for a symmetric distribution.

\begin{popfigure}
\begin{center}
\begin{tikzpicture}[scale=1.1]
\draw[->, thick] (-0.6,0) -- (6,0) node[below left] {$x$};
\draw[->, thick] (0,-0.3) -- (0,1.3) node[below left] {$F(x)$};
\draw[very thick, poporange] (-0.6,0) -- (1.5,0);
\draw[very thick, poporange] (1.5,0) -- (5.5,1);
\draw[very thick, poporange] (5.5,1) -- (6,1);
\draw[dashed, popmuted] (1.5,0) node[below] {$a$} -- (1.5,0.05);
\draw[dashed, popmuted] (5.5,0) node[below] {$b$} -- (5.5,1);
\draw[dashed, popmuted] (0,1) node[left] {$1$} -- (5.5,1);
\fill[poporange] (1.5,0) circle (1.5pt);
\fill[poporange] (5.5,1) circle (1.5pt);
\end{tikzpicture}
\end{center}
\legende{Cdf of $U(a,b)$: zero before $a$, a straight ramp of slope $\frac{1}{b-a}$ on $[a,b]$, then $1$ after $b$.}
\end{popfigure}

\courssub{Method and Worked Examples}

\begin{methode}[Computing probabilities with a uniform distribution]
\begin{enumerate}
\item \textbf{Sketch the pdf}: a rectangle over $[a,b]$ of height $\tfrac{1}{b-a}$.
\item \textbf{Locate} the interval asked for on the sketch.
\item \textbf{Compute the area} of the corresponding rectangle: probability $= \dfrac{\text{length of sub-interval}}{b-a}$.
\item For $E(X)$ and $\operatorname{Var}(X)$, quote the standard formulae directly.
\end{enumerate}
\end{methode}

\begin{exemplebox}[Exam tip]
Always sketch the pdf first. With a uniform distribution the picture literally \emph{is} the calculation: probabilities are areas of rectangles, so a correct sketch gives the answer immediately and lets you spot absurd results (e.g.\ a probability greater than 1).
\end{exemplebox}

\begin{exoresolu}[Waiting for the bus]
The waiting time $X$ (in minutes) for a bus at a stop is uniformly distributed on $[0, 20]$.
\begin{enumerate}
\item Find the probability that a passenger waits between 6 and 14 minutes.
\item Find the probability of waiting more than 15 minutes.
\item Find the mean and standard deviation of the waiting time.
\item Given that the passenger has already waited 10 minutes, find the probability that the total wait exceeds 15 minutes.
\end{enumerate}
\tcblower
\textbf{Solution.} Here $a = 0$, $b = 20$, so $f(x) = \tfrac{1}{20}$ on $[0,20]$.
\begin{enumerate}
\item $P(6 < X < 14) = \dfrac{14 - 6}{20} = \dfrac{8}{20} = 0.4$.
\item $P(X > 15) = \dfrac{20 - 15}{20} = \dfrac{5}{20} = 0.25$.
\item $E(X) = \dfrac{0 + 20}{2} = 10$ minutes; $\operatorname{Var}(X) = \dfrac{20^2}{12} = \dfrac{100}{3}$, so $\sigma = \sqrt{\tfrac{100}{3}} \approx 5.77$ minutes.
\item Conditional probability:
\[
P(X > 15 \mid X > 10) = \frac{P(X > 15)}{P(X > 10)} = \frac{5/20}{10/20} = \frac{1}{2}.
\]
Given 10 minutes have elapsed, the remaining wait is uniform on $[10, 20]$, and half of that interval lies beyond 15 — the picture confirms the answer.
\end{enumerate}
\end{exoresolu}

\begin{exoresolu}[Finding $b$ from the variance]
A random variable $X$ is uniformly distributed on $[2, b]$ with $b > 2$, and $\operatorname{Var}(X) = 12$. Find $b$, and hence write down $E(X)$ and $P(X > 8)$.
\tcblower
\textbf{Solution.} We use $\operatorname{Var}(X) = \dfrac{(b-a)^2}{12}$ with $a = 2$:
\[
\frac{(b-2)^2}{12} = 12 \implies (b-2)^2 = 144 \implies b - 2 = 12 \implies b = 14
\]
(taking the positive root since $b > 2$). Then
\[
E(X) = \frac{2 + 14}{2} = 8, \qquad P(X > 8) = \frac{14 - 8}{14 - 2} = \frac{6}{12} = \frac{1}{2}.
\]
Note that $P(X > 8) = \tfrac12$ is obvious from the sketch, since $8$ is the midpoint of $[2, 14]$.
\end{exoresolu}

\courssub{Applications of the Rectangular Distribution}

The uniform distribution appears whenever "all values in a range are equally likely":

\begin{itemize}
\item \textbf{Random waiting.} A bus, a shared taxi or a customer known to arrive "at any time within the next $T$ minutes" gives $X \sim U(0, T)$.
\item \textbf{Rounding errors.} When a measurement is rounded to the nearest integer (or nearest 100 FCFA in a market transaction), the rounding error is modelled by $U(-0.5, 0.5)$, with mean $0$ and variance $\tfrac{1}{12}$. This is why the number $12$ appears so often in error analysis.
\item \textbf{Random number generation.} The random numbers produced by a calculator or computer between $0$ and $1$ follow $U(0,1)$; this is the foundation of simulation methods.
\end{itemize}

\begin{exoresolu}[Rounding errors]
The mass of a bag of garri is recorded to the nearest gram. The rounding error $E$ (in grams) is uniformly distributed on $[-0.5, 0.5]$. Find the probability that the error exceeds $0.3$ g in absolute value, and the standard deviation of $E$.
\tcblower
\textbf{Solution.} We need $P(|E| > 0.3) = P(E > 0.3) + P(E < -0.3)$. By symmetry of the rectangle:
\[
P(|E| > 0.3) = 2 \times \frac{0.5 - 0.3}{1} = 2 \times 0.2 = 0.4.
\]
The standard deviation is
\[
\sigma = \frac{b - a}{\sqrt{12}} = \frac{1}{\sqrt{12}} = \frac{1}{2\sqrt{3}} \approx 0.289 \text{ g}.
\]
\end{exoresolu}

\begin{savaistu}[Why $\frac{1}{12}$?]
The variance $\tfrac{(b-a)^2}{12}$ of the uniform distribution is one of the most quoted results in engineering and metrology. When instruments round readings, the total uncertainty of a sum of many rounded measurements is computed using exactly this $\tfrac{1}{12}$ factor. The same formula underlies the "quantisation noise" theory used in digital music and mobile phone signal processing — technology used every day in Douala and Bamenda!
\end{savaistu}

\begin{exercice}
The time $T$ (in minutes) taken by a photocopying machine at a cybercafé in Buea to warm up is uniformly distributed between $1$ and $5$ minutes.
\begin{enumerate}
\item Write down the pdf of $T$ and sketch it.
\item Find $P(T < 2.5)$ and $P(2 < T < 4.5)$.
\item Calculate $E(T)$ and $\operatorname{Var}(T)$.
\item Find the value $t$ such that $P(T < t) = 0.75$.
\end{enumerate}
\end{exercice}

\begin{corrige}[Exercise 1]
\begin{enumerate}
\item $f(t) = \tfrac{1}{4}$ for $1 \le t \le 5$, $0$ otherwise: a rectangle of height $0.25$ over $[1,5]$.
\item $P(T < 2.5) = \dfrac{2.5 - 1}{4} = 0.375$; \quad $P(2 < T < 4.5) = \dfrac{4.5 - 2}{4} = 0.625$.
\item $E(T) = \dfrac{1+5}{2} = 3$ minutes; \quad $\operatorname{Var}(T) = \dfrac{(5-1)^2}{12} = \dfrac{16}{12} = \dfrac{4}{3}$ min$^2$.
\item $F(t) = \dfrac{t-1}{4} = 0.75 \implies t - 1 = 3 \implies t = 4$ minutes.
\end{enumerate}
\end{corrige}

\begin{exercice}
$X \sim U(a, 10)$ and $P(X < 6) = 0.5$. Find $a$, $E(X)$ and $\operatorname{Var}(X)$.
\end{exercice}

\begin{corrige}[Exercise 2]
$P(X < 6) = \dfrac{6 - a}{10 - a} = 0.5 \implies 2(6 - a) = 10 - a \implies 12 - 2a = 10 - a \implies a = 2$.
Then $E(X) = \dfrac{2 + 10}{2} = 6$ (consistent with $P(X<6) = 0.5$) and $\operatorname{Var}(X) = \dfrac{(10-2)^2}{12} = \dfrac{64}{12} = \dfrac{16}{3}$.
\end{corrige}

\begin{attention}[Final checklist for the GCE]
\begin{itemize}
\item Check that the interval given in the question really is $[a,b]$ before applying formulae.
\item Square $(b-a)$ in the variance formula.
\item For conditional probabilities, restrict the interval first, then use lengths.
\item Sketch the rectangle: it costs ten seconds and saves marks.
\end{itemize}
\end{attention}

\coursec{The Exponential Distribution}

In the previous section we studied the uniform distribution, in which every value in an interval $[a,b]$ is \emph{equally likely}. That model suits situations with complete symmetry. But many real waiting times are not symmetric at all: the time you wait for a taxi in Douala, the lifetime of a phone battery, the gap between two customers arriving at an MTN service counter. Such times are often short, occasionally long, and never negative. The distribution that models them is the \cle{exponential distribution}, one of the most important special continuous distributions in the GCE syllabus.

\courssub{Definition and probability density function}

\begin{experience}[A motivating observation]
Customers arrive at a bank counter in Yaound\'e at an average rate of $\lambda = 2$ customers per minute. Let $X$ be the waiting time (in minutes) between two consecutive arrivals. Observation shows that very short gaps are the most frequent, and long gaps become rarer and rarer. A histogram of such gaps has the shape of a rapidly decreasing curve starting at $x = 0$. The curve $f(x) = \lambda e^{-\lambda x}$ fits this behaviour exactly.
\end{experience}

\begin{definition}[Exponential distribution]
A continuous random variable $X$ follows the \cle{exponential distribution} with parameter $\lambda > 0$, written $X \sim \mathrm{Exp}(\lambda)$, if its probability density function is
\[
f(x) = \begin{cases} \lambda e^{-\lambda x} & \text{if } x \geq 0, \\ 0 & \text{if } x < 0. \end{cases}
\]
The parameter $\lambda$ is a \cle{rate}: the average number of events per unit time.
\end{definition}

Let us check that $f$ is a genuine pdf. Clearly $f(x) \geq 0$ for all $x$, and
\[
\int_{-\infty}^{+\infty} f(x)\,dx = \int_{0}^{+\infty} \lambda e^{-\lambda x}\,dx = \left[-e^{-\lambda x}\right]_{0}^{+\infty} = 0 - (-1) = 1.
\]

\begin{popfigure}
\begin{center}
\begin{tikzpicture}
\begin{axis}[
  width=11cm, height=6.5cm,
  axis lines=left,
  xlabel={$x$}, ylabel={$f(x)$},
  xmin=0, xmax=6, ymin=0, ymax=2.2,
  domain=0:6, samples=100,
  legend style={at={(0.97,0.97)}, anchor=north east, font=\small}
]
\addplot[popblue, very thick] {0.5*exp(-0.5*x)};
\addlegendentry{$\lambda = 0.5$}
\addplot[poporange, very thick] {exp(-x)};
\addlegendentry{$\lambda = 1$}
\addplot[poppurple, very thick] {2*exp(-2*x)};
\addlegendentry{$\lambda = 2$}
\end{axis}
\end{tikzpicture}
\end{center}
\legende{The pdf $f(x) = \lambda e^{-\lambda x}$ for $\lambda = 0.5,\ 1,\ 2$. Larger $\lambda$ means faster decay: events occur more often, so long waiting times become rarer.}
\end{popfigure}

\courssub{The cumulative distribution function and survival probability}

For $x \geq 0$,
\[
F(x) = P(X \leq x) = \int_{0}^{x} \lambda e^{-\lambda t}\,dt = \left[-e^{-\lambda t}\right]_{0}^{x} = 1 - e^{-\lambda x}.
\]

\begin{propriete}[CDF and survival probability]
If $X \sim \mathrm{Exp}(\lambda)$, then for all $x \geq 0$:
\[
F(x) = P(X \leq x) = 1 - e^{-\lambda x}, \qquad P(X > x) = e^{-\lambda x}.
\]
The quantity $P(X > x)$ is called the \cle{survival probability}: the probability that the event has \emph{not yet} occurred by time $x$.
\end{propriete}

\begin{exemplebox}[First calculation]
The lifetime $X$ (in years) of a certain electronic component is exponential with $\lambda = 0.25$ per year. Then
\[
P(X > 4) = e^{-0.25 \times 4} = e^{-1} \approx 0.3679,
\]
and $P(X \leq 4) = 1 - e^{-1} \approx 0.6321$.
\end{exemplebox}

\courssub{Expectation and variance of the exponential distribution}

These derivations \textbf{are examinable}: you must be able to carry out the integration by parts carefully, with limits.

\textbf{Expectation.} By definition,
\[
E(X) = \int_{0}^{+\infty} x\,\lambda e^{-\lambda x}\,dx.
\]
Integrate by parts with $u = x$ (so $du = dx$) and $dv = \lambda e^{-\lambda x}dx$ (so $v = -e^{-\lambda x}$):
\begin{align*}
E(X) &= \left[-x e^{-\lambda x}\right]_{0}^{+\infty} + \int_{0}^{+\infty} e^{-\lambda x}\,dx \\
&= (0 - 0) + \left[-\frac{1}{\lambda}e^{-\lambda x}\right]_{0}^{+\infty} = \frac{1}{\lambda},
\end{align*}
since $x e^{-\lambda x} \to 0$ as $x \to +\infty$ (exponential decay beats linear growth).

\textbf{Variance.} First compute $E(X^2) = \displaystyle\int_{0}^{+\infty} x^2 \lambda e^{-\lambda x}\,dx$. Integrate by parts with $u = x^2$, $dv = \lambda e^{-\lambda x}dx$:
\begin{align*}
E(X^2) &= \left[-x^2 e^{-\lambda x}\right]_{0}^{+\infty} + \int_{0}^{+\infty} 2x\,e^{-\lambda x}\,dx \\
&= 0 + \frac{2}{\lambda}\int_{0}^{+\infty} x\,\lambda e^{-\lambda x}\,dx = \frac{2}{\lambda}\,E(X) = \frac{2}{\lambda^2}.
\end{align*}
Therefore
\[
\mathrm{Var}(X) = E(X^2) - [E(X)]^2 = \frac{2}{\lambda^2} - \frac{1}{\lambda^2} = \frac{1}{\lambda^2}.
\]

\begin{propriete}[Mean and variance of $\mathrm{Exp}(\lambda)$ — proof examinable]
If $X \sim \mathrm{Exp}(\lambda)$, then
\[
E(X) = \frac{1}{\lambda}, \qquad \mathrm{Var}(X) = \frac{1}{\lambda^2}, \qquad \sigma = \frac{1}{\lambda}.
\]
Interpretation: $1/\lambda$ is the \cle{mean waiting time} between events. Remarkably, the standard deviation equals the mean.
\end{propriete}

\begin{attention}[$\lambda$ is a rate — convert units first!]
$\lambda$ must be expressed in events \emph{per unit of the time used for $x$}. If calls arrive at a rate of $3$ per hour and $X$ is measured in \emph{minutes}, then $\lambda = \frac{3}{60} = 0.05$ per minute, and $E(X) = \frac{1}{0.05} = 20$ minutes. Mixing hours and minutes is the most frequent error in GCE scripts. Always state your unit and convert before substituting.
\end{attention}

\courssub{The memoryless (no-ageing) property}

\begin{experience}[A surprising fact]
A bulb has already burned for $1000$ hours without failing. Does this make failure in the next hour \emph{more} likely? For an exponential lifetime, the answer is \textbf{no}: the bulb is ``as good as new''. The distribution \emph{forgets} the past.
\end{experience}

\begin{propriete}[Memoryless property — proof examinable]
If $X \sim \mathrm{Exp}(\lambda)$, then for all $s, t \geq 0$:
\[
P(X > s + t \mid X > s) = P(X > t).
\]
\textbf{Proof.} By the definition of conditional probability, and since $\{X > s+t\} \subset \{X > s\}$:
\[
P(X > s+t \mid X > s) = \frac{P(X > s+t)}{P(X > s)} = \frac{e^{-\lambda(s+t)}}{e^{-\lambda s}} = e^{-\lambda t} = P(X > t). \qquad \blacksquare
\]
\end{propriete}

Interpretation: given that the component has survived $s$ units of time, the probability that it survives a \emph{further} $t$ units is the same as the probability that a brand-new component survives $t$ units. The exponential distribution is the \emph{only} continuous distribution with this property.

\begin{exemplebox}[Using memorylessness]
With $X \sim \mathrm{Exp}(0.25)$ (lifetimes in years), the probability that a component already 4 years old lasts at least 2 more years is
\[
P(X > 6 \mid X > 4) = P(X > 2) = e^{-0.25 \times 2} = e^{-0.5} \approx 0.6065.
\]
\end{exemplebox}

\courssub{The link with the Poisson distribution}

Suppose events occur randomly in time at a constant average rate $\lambda$ per unit time, independently of one another (a \cle{Poisson process}). Then two random variables describe the same situation from two different points of view:

\begin{itemize}
\item the \textbf{number} $N$ of events occurring in a time interval of length $t$ is discrete: $N \sim \mathrm{Po}(\lambda t)$;
\item the \textbf{waiting time} $T$ between two consecutive events is continuous: $T \sim \mathrm{Exp}(\lambda)$.
\end{itemize}

\begin{propriete}[Poisson counts $\leftrightarrow$ exponential waiting times]
In a Poisson process of rate $\lambda$, the waiting time $T$ until the next event satisfies $T \sim \mathrm{Exp}(\lambda)$. Indeed, ``$T > t$'' means ``no event occurs in $[0,t]$'', and since the count in $[0,t]$ is $\mathrm{Po}(\lambda t)$:
\[
P(T > t) = P(N = 0) = \frac{e^{-\lambda t}(\lambda t)^0}{0!} = e^{-\lambda t},
\]
which is exactly the survival probability of $\mathrm{Exp}(\lambda)$.
\end{propriete}

\begin{popfigure}
\begin{center}
\begin{tikzpicture}[scale=1.0]
\draw[->, thick] (0,0) -- (10.6,0) node[right] {time};
\foreach \x/\lab in {1.3/{$E_1$}, 2.1/{$E_2$}, 4.6/{$E_3$}, 5.4/{$E_4$}, 8.9/{$E_5$}} {
  \draw[popblue, very thick] (\x,0) -- (\x,0.35);
  \node[popblue, above, font=\small] at (\x,0.35) {\lab};
}
\draw[<->, poporange, thick] (1.3,-0.5) -- (2.1,-0.5) node[midway, below, font=\small] {$T_1$};
\draw[<->, poporange, thick] (2.1,-0.5) -- (4.6,-0.5) node[midway, below, font=\small] {$T_2$};
\draw[<->, poporange, thick] (4.6,-0.5) -- (5.4,-0.5) node[midway, below, font=\small] {$T_3$};
\draw[<->, poporange, thick] (5.4,-0.5) -- (8.9,-0.5) node[midway, below, font=\small] {$T_4$};
\node[font=\small, popdark] at (5.3,-1.2) {Each gap $T_i \sim \mathrm{Exp}(\lambda)$; the count of events in any interval of length $t$ is $\mathrm{Po}(\lambda t)$.};
\end{tikzpicture}
\end{center}
\legende{A Poisson process on the time axis: events $E_1, E_2, \dots$ occur at rate $\lambda$; the gaps between consecutive events are independent exponential waiting times.}
\end{popfigure}

\begin{methode}[Poisson or exponential? Choosing the right model]
\begin{enumerate}
\item Read the question and identify the random variable.
\item If it \textbf{counts} events in a fixed interval (``how many customers in 10 minutes''), use $N \sim \mathrm{Po}(\lambda t)$ with $t$ the length of the interval.
\item If it measures a \textbf{waiting time, duration or lifetime} (``how long until the next call''), use $T \sim \mathrm{Exp}(\lambda)$: $P(T > t) = e^{-\lambda t}$, $P(T \leq t) = 1 - e^{-\lambda t}$.
\item Convert $\lambda$ to the same time unit as the variable before substituting.
\item Translate boundary phrases: ``at least'' $\to$ $P(T \geq t) = e^{-\lambda t}$; ``no event in the interval'' $\to$ $P(N = 0) = e^{-\lambda t}$.
\end{enumerate}
\end{methode}

\begin{exoresolu}[GCE-style problem mixing Poisson and exponential]
Calls arrive at a call centre in Bamenda at an average rate of $4$ per hour, according to a Poisson process.\\
(a) Find the probability that exactly $3$ calls arrive in one hour.\\
(b) Find the probability that the operator waits more than $30$ minutes for the next call.\\
(c) Given that no call has arrived for $20$ minutes, find the probability that the operator waits at least $10$ further minutes.\\
(d) Find the mean and standard deviation of the waiting time between calls.
\tcblower
Work in \textbf{hours}: $\lambda = 4$ per hour.\\
\textbf{(a)} The count in one hour is $N \sim \mathrm{Po}(4)$:
\[
P(N = 3) = \frac{e^{-4}\,4^3}{3!} = \frac{64e^{-4}}{6} \approx 0.1954.
\]
\textbf{(b)} Let $T$ be the waiting time in hours, $T \sim \mathrm{Exp}(4)$. Since $30$ min $= 0.5$ h:
\[
P(T > 0.5) = e^{-4 \times 0.5} = e^{-2} \approx 0.1353.
\]
\textbf{(c)} By the memoryless property:
\[
P\!\left(T > \tfrac{30}{60} \,\middle|\, T > \tfrac{20}{60}\right) = P\!\left(T > \tfrac{10}{60}\right) = e^{-4/6} = e^{-2/3} \approx 0.5134.
\]
\textbf{(d)} $E(T) = \dfrac{1}{\lambda} = \dfrac{1}{4}$ h $= 15$ minutes, and $\sigma = \dfrac{1}{\lambda} = 15$ minutes as well.
\end{exoresolu}

\begin{attention}[Common mistakes]
\begin{itemize}
\item Using $P(X > t) = 1 - e^{-\lambda t}$: this is $F(t) = P(X \leq t)$. The \emph{survival} probability is $e^{-\lambda t}$.
\item Forgetting that for a continuous variable $P(X = t) = 0$, so $P(X \geq t) = P(X > t)$.
\item Writing $E(X) = \lambda$: the mean is $\dfrac{1}{\lambda}$, the mean \emph{waiting time}.
\item Applying the memoryless property to a non-exponential lifetime without justification.
\end{itemize}
\end{attention}

\begin{exercice}[Practice]
The time $X$ (in minutes) between arrivals of customers at a petrol station in Buea is exponential with mean $5$ minutes.\\
(a) Write down $\lambda$ and the pdf of $X$.\\
(b) Find $P(X > 10)$ and $P(2 \leq X \leq 8)$.\\
(c) A customer has just arrived. Find the probability that exactly two customers arrive in the next $10$ minutes.\\
(d) Given that nobody has arrived for $6$ minutes, find the probability of waiting at least $4$ more minutes.
\end{exercice}

\begin{corrige}[Exercise 1]
(a) $E(X) = 1/\lambda = 5$, so $\lambda = 0.2$ per minute and $f(x) = 0.2e^{-0.2x}$, $x \geq 0$.\\
(b) $P(X > 10) = e^{-2} \approx 0.1353$; $P(2 \leq X \leq 8) = e^{-0.4} - e^{-1.6} \approx 0.6703 - 0.2019 = 0.4684$.\\
(c) Count in $10$ min: $N \sim \mathrm{Po}(2)$, so $P(N = 2) = \frac{e^{-2}2^2}{2!} = 2e^{-2} \approx 0.2707$.\\
(d) Memoryless: $P(X > 10 \mid X > 6) = P(X > 4) = e^{-0.8} \approx 0.4493$.
\end{corrige}

\begin{aretenir}[Key points]
\begin{itemize}
\item $X \sim \mathrm{Exp}(\lambda)$: $f(x) = \lambda e^{-\lambda x}$ ($x \geq 0$), $F(x) = 1 - e^{-\lambda x}$, $P(X > x) = e^{-\lambda x}$.
\item $E(X) = \dfrac{1}{\lambda}$, $\mathrm{Var}(X) = \dfrac{1}{\lambda^2}$ (derivation by parts examinable).
\item $\lambda$ is a rate per unit time; $1/\lambda$ is the mean waiting time. Convert units!
\item Memoryless: $P(X > s+t \mid X > s) = P(X > t)$.
\item Poisson process of rate $\lambda$: counts are $\mathrm{Po}(\lambda t)$, gaps are $\mathrm{Exp}(\lambda)$ — two faces of the same situation.
\end{itemize}
\end{aretenir}

\begin{savaistu}[Did you know?]
The exponential distribution models radioactive decay: the waiting time until a single unstable nucleus decays is exponential, which is why the \emph{number} of decays per second in a sample follows a Poisson law. The same mathematics that describes a Geiger counter in a physics laboratory describes the queue at an Express Union counter in town — a beautiful example of the unity of mathematics.
\end{savaistu}

\coursec{The Normal Distribution and Standardisation}

In the previous section we studied the exponential distribution, a model for \emph{waiting times} which is strongly skewed: small values are the most likely, and the density decreases steadily. Nature, however, offers us a very different and far more common pattern. Measure the heights of all the students in an Upper Sixth class in Bamenda, the masses of bags of cement leaving a factory in Douala, or the marks of thousands of candidates at the GCE: in each case the values cluster around a central average, and values far from the centre --- on \emph{both} sides --- become rarer and rarer, in a perfectly symmetric way. The histogram has the shape of a \cle{bell}. The mathematical model behind this bell shape is the \cle{Normal distribution}, the most important distribution in all of statistics, and the heart of the GCE Advanced Level statistics paper.

\courssub{Discovering the Bell Curve}

\begin{experience}[From a histogram to a curve]
Imagine weighing 1000 loaves of bread sold as ``\SI{500}{g} loaves''. Most loaves weigh close to \SI{500}{g}; a few weigh \SI{490}{g} or \SI{510}{g}; extremely few weigh \SI{470}{g} or \SI{530}{g}. If we draw a histogram with very narrow classes and then smooth its outline, we obtain a continuous curve with three striking features:
\begin{itemize}
\item it is \textbf{symmetric} about the mean $\mu$;
\item it has a \textbf{single peak} (mode) at $\mu$;
\item its two tails approach the horizontal axis but \textbf{never touch it}.
\end{itemize}
This is exactly the shape of the Normal density curve.
\end{experience}

\begin{definition}[Normal distribution]
A continuous random variable $X$ follows a \cle{Normal distribution} with mean $\mu$ and variance $\sigma^2$, written $X \sim N(\mu, \sigma^2)$, when its probability density function is
\[ f(x) = \frac{1}{\sigma\sqrt{2\pi}}\, e^{-\frac{(x-\mu)^2}{2\sigma^2}}, \qquad x \in \mathbb{R}. \]
The curve of $f$ is called the \emph{Normal curve} or \emph{bell curve}. It satisfies:
\begin{itemize}
\item symmetry about the vertical line $x = \mu$;
\item a maximum at $x = \mu$, of value $\dfrac{1}{\sigma\sqrt{2\pi}}$;
\item points of inflexion at $x = \mu - \sigma$ and $x = \mu + \sigma$;
\item $f(x) > 0$ for all $x$, and $f(x) \to 0$ as $x \to \pm\infty$ (the curve never touches the axis).
\end{itemize}
\end{definition}

\begin{attention}[Recognition only]
At the GCE Advanced Level you are \textbf{never} asked to integrate this density: the exponential of $-x^2$ has no elementary antiderivative. All probabilities are obtained from \textbf{tables}, after standardisation. You must, however, be able to recognise the parameters: in $N(\mu, \sigma^2)$, the second parameter is the \textbf{variance}, not the standard deviation. Thus $X \sim N(50, 9)$ means $\mu = 50$ and $\sigma = 3$.
\end{attention}

The two parameters play clearly separated roles, which you can see on the figure below:
\begin{itemize}
\item $\mu$ is a \textbf{location} parameter: changing $\mu$ translates the curve horizontally without changing its shape;
\item $\sigma$ is a \textbf{spread} parameter: a larger $\sigma$ gives a flatter, wider curve; a smaller $\sigma$ gives a tall, narrow curve.
\end{itemize}

\begin{popfigure}
\begin{center}
\begin{tikzpicture}
\begin{axis}[
  width=13cm, height=7cm,
  axis lines=middle,
  xlabel={$x$}, ylabel={$f(x)$},
  xmin=-7, xmax=8, ymin=0, ymax=0.45,
  xtick={-4,-2,0,2,4,6},
  legend style={at={(0.98,0.95)}, anchor=north east, font=\small},
  samples=100, domain=-7:8, thick
]
\addplot[popblue] {0.3989*exp(-x^2/2)};
\addlegendentry{$N(0,1)$}
\addplot[poporange] {0.3989*exp(-(x-2)^2/2)};
\addlegendentry{$N(2,1)$}
\addplot[poppurple] {0.1995*exp(-x^2/8)};
\addlegendentry{$N(0,4)$}
\end{axis}
\end{tikzpicture}
\end{center}
\legende{Effect of the parameters: $N(2,1)$ is $N(0,1)$ shifted $2$ units right (role of $\mu$); $N(0,4)$, with $\sigma=2$, is flatter and wider (role of $\sigma$).}
\end{popfigure}

\begin{propriete}[Fundamental characteristics of $N(\mu,\sigma^2)$]
If $X \sim N(\mu, \sigma^2)$ then
\[ E(X) = \mu, \qquad \mathrm{Var}(X) = \sigma^2, \qquad P(X < \mu) = P(X > \mu) = 0.5, \]
and the total area under the curve equals $1$. The symmetry about $\mu$ implies in particular that the mean, median and mode of a Normal distribution all coincide.
\end{propriete}

\courssub{The Standard Normal Variable}

There are infinitely many Normal distributions, one for each pair $(\mu, \sigma^2)$. Printing tables for all of them is impossible. The brilliant solution is to \textbf{reduce every Normal variable to one single reference variable}, by a simple change of origin and scale.

\begin{definition}[Standard Normal distribution]
The \cle{standard Normal distribution} is $N(0,1)$: mean $0$, variance $1$. A variable with this distribution is usually denoted $Z$, and its cumulative distribution function is denoted by the Greek letter Phi:
\[ \Phi(z) = P(Z < z). \]
\end{definition}

\begin{propriete}[Expectation and variance of $Z$ --- proof not examinable]
For $Z \sim N(0,1)$:
\[ E(Z) = 0 \qquad \text{and} \qquad \mathrm{Var}(Z) = 1. \]
\emph{Illustration.} The density of $Z$ is $\varphi(z) = \frac{1}{\sqrt{2\pi}}e^{-z^2/2}$, an \textbf{even} function (symmetric about $0$). Hence $z\,\varphi(z)$ is an odd function, and the integral defining $E(Z)$ over $\mathbb{R}$ is $0$. A computation (integration by parts) then gives $E(Z^2) = 1$, so $\mathrm{Var}(Z) = E(Z^2) - [E(Z)]^2 = 1 - 0 = 1$. This is exactly what ``standard'' means: centred (mean $0$) and reduced (standard deviation $1$).
\end{propriete}

\begin{propriete}[Standardisation --- examinable and essential]
If $X \sim N(\mu, \sigma^2)$, then the variable
\[ Z = \frac{X - \mu}{\sigma} \quad \text{follows} \quad N(0,1). \]
Consequently, for any real numbers $a < b$:
\[ P(a < X < b) = P\!\left(\frac{a-\mu}{\sigma} < Z < \frac{b-\mu}{\sigma}\right). \]
\end{propriete}

\emph{Justification.} The transformation $Z = \frac{X-\mu}{\sigma}$ is linear, so $Z$ remains Normally distributed. Moreover $E(Z) = \frac{E(X)-\mu}{\sigma} = 0$ and $\mathrm{Var}(Z) = \frac{\mathrm{Var}(X)}{\sigma^2} = \frac{\sigma^2}{\sigma^2} = 1$. Hence $Z \sim N(0,1)$. \hfill$\blacksquare$

\begin{methode}[Standardising $X \sim N(\mu, \sigma^2)$]
\begin{enumerate}
\item Identify $\mu$ and $\sigma$ (remember: the second parameter of $N(\mu,\sigma^2)$ is the \textbf{variance}, so take its square root).
\item Convert each boundary value $x$ into $z = \dfrac{x-\mu}{\sigma}$.
\item Sketch the standard Normal curve and \textbf{shade} the required region.
\item Read $\Phi$ from the tables and use the properties of $\Phi$ to obtain the answer.
\end{enumerate}
\end{methode}

\courssub{Reading the Standard Normal Tables}

The GCE tables give $\Phi(z) = P(Z < z)$ for $z \geq 0$, usually to two decimal places: the row gives $z$ to one decimal, the column gives the second decimal. For example, for $z = 1.23$ we read the row $1.2$ and the column $0.03$:

\begin{center}
\begin{tabularx}{\linewidth}{|l|X|X|X|X|}
\hline
\rowcolor{popblueL} $z$ & $0.00$ & $0.01$ & $0.02$ & $0.03$ \\
\hline
$1.1$ & $0.8643$ & $0.8665$ & $0.8686$ & $0.8708$ \\
\hline
$1.2$ & $0.8849$ & $0.8869$ & $0.8888$ & $\mathbf{0.8907}$ \\
\hline
$1.3$ & $0.9032$ & $0.9049$ & $0.9066$ & $0.9082$ \\
\hline
\end{tabularx}
\end{center}

So $\Phi(1.23) = 0.8907$. When the exam requires more precision, we \textbf{interpolate linearly} between two consecutive entries. For instance, for $z = 1.645$: $\Phi(1.64) = 0.9495$ and $\Phi(1.65) = 0.9505$, so
\[ \Phi(1.645) \approx \tfrac{1}{2}(0.9495 + 0.9505) = 0.9500. \]

\begin{attention}[The table gives a LEFT tail]
The table gives $\Phi(z) = P(Z < z)$, the area to the \textbf{left} of $z$. You can \emph{never} read a right-tail probability $P(Z > z)$ directly: you must first rewrite it as $1 - \Phi(z)$. Reading $\Phi(1.5)$ as the answer to $P(Z > 1.5)$ is one of the most frequent errors at the GCE.
\end{attention}

\courssub{Properties of \texorpdfstring{$\Phi$}{Phi}}

Because the standard Normal curve is symmetric about $0$, the area to the left of $-z$ equals the area to the right of $z$. This single geometric fact generates all the working formulas.

\begin{propriete}[Properties of $\Phi$ --- must be known perfectly]
For $Z \sim N(0,1)$ and any real numbers $z$, $a < b$:
\begin{align*}
\Phi(-z) &= 1 - \Phi(z) && \text{(symmetry)}\\
P(Z > z) &= 1 - \Phi(z) && \text{(upper tail)}\\
P(Z < -z) &= \Phi(-z) = 1 - \Phi(z)\\
P(a < Z < b) &= \Phi(b) - \Phi(a) && \text{(interval)}\\
P(|Z| < z) &= 2\Phi(z) - 1, \qquad z \geq 0 && \text{(central interval)}
\end{align*}
\end{propriete}

\begin{popfigure}
\begin{center}
\begin{tikzpicture}[scale=1.05]
\fill[popblueL] plot[domain=-3.5:1, samples=60] (\x, {1.6*exp(-\x*\x/2)}) -- (1,0) -- (-3.5,0) -- cycle;
\fill[popgreenL] plot[domain=1:3.5, samples=40] (\x, {1.6*exp(-\x*\x/2)}) -- (3.5,0) -- (1,0) -- cycle;
\draw[popdark, thick] plot[domain=-3.5:3.5, samples=80] (\x, {1.6*exp(-\x*\x/2)});
\draw[->] (-4,0) -- (4.2,0) node[below left] {$z$};
\draw[->] (0,-0.15) -- (0,2) node[left] {$\varphi(z)$};
\draw[dashed, popblue] (1,0) -- (1,{1.6*exp(-0.5)});
\node[below, popblue] at (1,0) {$z$};
\draw[dashed, poppurple] (-1,0) -- (-1,{1.6*exp(-0.5)});
\node[below, poppurple] at (-1,0) {$-z$};
\node[below] at (0,0) {$0$};
\node[popblue, align=center] at (-1.9,0.45) {$\Phi(z)$};
\node[poppurple, align=center, font=\small] at (2.6,0.45) {$1-\Phi(z)$};
\end{tikzpicture}
\end{center}
\legende{The shaded blue area is $\Phi(z)=P(Z<z)$. By symmetry, the area left of $-z$ equals the area right of $z$: $\Phi(-z)=1-\Phi(z)$.}
\end{popfigure}

\begin{exemplebox}[Warm-up with the tables]
Let $Z \sim N(0,1)$. Using $\Phi(1.20) = 0.8849$ and $\Phi(0.85) = 0.8023$:
\begin{itemize}
\item $P(Z < 1.20) = \Phi(1.20) = 0.8849$;
\item $P(Z > 1.20) = 1 - 0.8849 = 0.1151$;
\item $P(Z < -1.20) = \Phi(-1.20) = 1 - 0.8849 = 0.1151$;
\item $P(-1.20 < Z < 1.20) = 2(0.8849) - 1 = 0.7698$;
\item $P(0.85 < Z < 1.20) = \Phi(1.20) - \Phi(0.85) = 0.8849 - 0.8023 = 0.0826$.
\end{itemize}
\end{exemplebox}

\begin{exoresolu}[Standardising a real situation]
The masses of loaves produced by a bakery are Normally distributed with mean $\mu = \SI{500}{g}$ and standard deviation $\sigma = \SI{12}{g}$. A loaf is chosen at random. Find the probability that its mass is (a) less than $\SI{482}{g}$; (b) between $\SI{494}{g}$ and $\SI{518}{g}$; (c) more than $\SI{527}{g}$.
\tcblower
Let $X \sim N(500, 12^2)$ and set $Z = \dfrac{X - 500}{12} \sim N(0,1)$.

\textbf{(a)} $z = \dfrac{482 - 500}{12} = -1.5$, so
\[ P(X < 482) = P(Z < -1.5) = 1 - \Phi(1.5) = 1 - 0.9332 = 0.0668. \]

\textbf{(b)} $z_1 = \dfrac{494-500}{12} = -0.5$ and $z_2 = \dfrac{518-500}{12} = 1.5$, so
\begin{align*}
P(494 < X < 518) &= \Phi(1.5) - \Phi(-0.5) = \Phi(1.5) - \bigl[1 - \Phi(0.5)\bigr] \\
&= 0.9332 - 1 + 0.6915 = 0.6247.
\end{align*}

\textbf{(c)} $z = \dfrac{527-500}{12} = 2.25$, so
\[ P(X > 527) = 1 - \Phi(2.25) = 1 - 0.9878 = 0.0122. \]
About $1.2\%$ of loaves exceed $\SI{527}{g}$: heavy loaves are rare, exactly as the bell shape predicts.
\end{exoresolu}

\courssub{Using the Tables in Reverse: Percentage Points}

Very often the question is inverted: we know the \emph{probability} and we look for the \emph{value} of $z$. We then read the table \textbf{inside-out}: locate the probability in the body of the table and read off the corresponding $z$.

\begin{definition}[Percentage point]
For $0 < p < 1$, the value $z_p$ such that $\Phi(z_p) = p$ is called a \cle{percentage point} (or quantile) of the standard Normal distribution. By symmetry, $\Phi(z) = p$ with $p < 0.5$ gives $z = -z_{1-p} < 0$.
\end{definition}

\begin{exemplebox}[Reverse reading]
Find $z$ such that $\Phi(z) = 0.9750$. Searching the body of the table, $0.9750$ corresponds to $z = 1.96$. Hence $P(Z < 1.96) = 0.975$ --- the famous value behind the central $95\%$ interval. Similarly, $\Phi(z) = 0.05$ gives $z = -1.645$, since $\Phi(1.645) = 0.95$.
\end{exemplebox}

The following percentage points occur so often at the GCE that they are worth memorising:

\begin{center}
\begin{tabularx}{\linewidth}{|l|X|X|X|X|}
\hline
\rowcolor{popblueL} $p$ & $0.90$ & $0.95$ & $0.975$ & $0.995$ \\
\hline
$z$ with $\Phi(z)=p$ & $1.282$ & $1.645$ & $1.960$ & $2.576$ \\
\hline
\end{tabularx}
\end{center}

\begin{methode}[Exam tip: sketch before you search]
Before touching the tables, \textbf{always} draw a quick bell curve, mark $0$ (or $\mu$), place the boundary value(s), and shade the required area. The sketch tells you immediately whether your answer should be greater or smaller than $0.5$, whether you need $\Phi(z)$, $1-\Phi(z)$ or a difference, and whether a reverse reading should give a positive or negative $z$. A ten-second sketch eliminates most table errors.
\end{methode}

\begin{attention}[Common mistakes with $\Phi$]
\begin{itemize}
\item Reading $P(Z > z)$ directly from the table instead of computing $1 - \Phi(z)$.
\item Forgetting that the table only covers $z \geq 0$: for negative values use $\Phi(-z) = 1 - \Phi(z)$.
\item Writing $P(a<Z<b) = \Phi(a) - \Phi(b)$ instead of $\Phi(b) - \Phi(a)$ (the result would be negative!).
\item Confusing $\sigma$ and $\sigma^2$ when standardising: divide by $\sigma$, the \textbf{standard deviation}.
\item Believing the curve touches the axis: probabilities such as $P(Z > 4)$ are tiny but never exactly $0$.
\end{itemize}
\end{attention}

\begin{exercice}[Practice]
Let $Z \sim N(0,1)$. Using the tables, find:
\begin{enumerate}
\item $P(Z < 0.75)$, $P(Z > 1.34)$, $P(Z < -2.05)$;
\item $P(-1.28 < Z < 2.10)$ and $P(|Z| < 1.96)$;
\item the value of $z$ such that $\Phi(z) = 0.9918$;
\item the value of $z$ such that $P(Z > z) = 0.025$.
\end{enumerate}
\end{exercice}

\begin{corrige}[Exercise 1]
\begin{enumerate}
\item $P(Z<0.75) = \Phi(0.75) = 0.7734$; \quad $P(Z>1.34) = 1 - \Phi(1.34) = 1 - 0.9099 = 0.0901$; \quad $P(Z<-2.05) = 1 - \Phi(2.05) = 1 - 0.9798 = 0.0202$.
\item $P(-1.28<Z<2.10) = \Phi(2.10) - \bigl[1-\Phi(1.28)\bigr] = 0.9821 - 1 + 0.8997 = 0.8818$; \quad $P(|Z|<1.96) = 2\Phi(1.96) - 1 = 2(0.9750) - 1 = 0.9500$.
\item $\Phi(z) = 0.9918$ corresponds to $z = 2.40$.
\item $P(Z > z) = 0.025 \iff \Phi(z) = 0.975$, hence $z = 1.96$.
\end{enumerate}
\end{corrige}

\begin{aretenir}[Key points of this section]
\begin{itemize}
\item $X \sim N(\mu, \sigma^2)$: bell curve, symmetric about $x=\mu$, inflexion points at $\mu \pm \sigma$, total area $1$, $P(X<\mu) = 0.5$.
\item Standardisation: $Z = \dfrac{X-\mu}{\sigma} \sim N(0,1)$, with $E(Z)=0$ and $\mathrm{Var}(Z)=1$.
\item The tables give the \textbf{left tail} $\Phi(z) = P(Z<z)$ for $z \geq 0$.
\item Master formulas: $\Phi(-z) = 1-\Phi(z)$; $P(Z>z) = 1-\Phi(z)$; $P(a<Z<b) = \Phi(b)-\Phi(a)$; $P(|Z|<z) = 2\Phi(z)-1$.
\item Reverse reading: find the probability inside the table, read $z$ on the border. Key values: $\Phi(1.645)=0.95$, $\Phi(1.96)=0.975$, $\Phi(2.576)=0.995$.
\item Always sketch and shade before using the tables.
\end{itemize}
\end{aretenir}

\begin{savaistu}[Why ``Normal''?]
The distribution is also called the \emph{Gaussian} distribution, after Carl Friedrich Gauss, who used it to model errors in astronomical measurements. Its ubiquity is explained by the \emph{Central Limit Theorem}: the sum (or average) of many small independent random effects is approximately Normal, whatever the original distributions. That is why heights, measurement errors, examination marks and so many natural phenomena all follow the same bell curve --- and why the Normal distribution will reappear throughout the rest of this chapter, from central intervals to approximations of the binomial and Poisson distributions.
\end{savaistu}

\coursec{Central Intervals and Working with Any Normal Distribution}

In the previous section we learnt how to \emph{standardise}: any Normal variable $X\sim N(\mu,\sigma^{2})$ is transformed into the standard Normal variable $Z=\dfrac{X-\mu}{\sigma}\sim N(0,1)$, whose probabilities are read from the tables of $\Phi(z)=P(Z<z)$. We now put this tool to work in the two directions that the GCE examiner loves: \textbf{forwards} (find the probability of an event concerning $X$) and \textbf{backwards} (given a probability, find the corresponding value of $X$). Along the way we shall meet the famous \cle{central 95\%} and \cle{central 99\%} intervals, which every candidate must know by heart.

\courssub{The Central 95\% and 99\% Intervals}

\textbf{Intuition.} We already know that a Normal variable clusters around its mean $\mu$: values far from $\mu$ are possible but rare. A natural question is: \emph{within what interval, symmetric about $\mu$, does the variable fall with a given high probability?} For example, a bottling company in Douala may wish to state: ``95\% of our bottles contain between $a$ and $b$ centilitres.'' Because the Normal curve is symmetric, the fairest interval is the \emph{central} one: we leave the same probability in each tail.

\begin{experience}[Finding the central 95\% point]
We want the number $z$ such that $P(-z<Z<z)=0.95$. Since the total area under the curve is $1$, the two tails together carry $0.05$, and by symmetry each tail carries $0.025$. Hence
\[
\Phi(z)=P(Z<z)=1-0.025=0.975.
\]
Reading the standard Normal tables \emph{backwards} (looking for $0.975$ in the body of the table), we find $\Phi(1.96)=0.9750$. Therefore $z=1.96$, that is
\[
P(-1.96<Z<1.96)=0.95.
\]
The same computation with $0.99$ in the centre gives $\Phi(z)=0.995$, and the tables give $z=2.576$ (more precisely $2.5758$).
\end{experience}

\begin{propriete}[Central intervals of a Normal distribution — proof examinable in outline]
If $X\sim N(\mu,\sigma^{2})$ and $Z=\dfrac{X-\mu}{\sigma}$, then:
\begin{itemize}
\item $P(-1.96<Z<1.96)=0.95$, equivalently $P(\mu-1.96\sigma<X<\mu+1.96\sigma)=0.95$;
\item $P(-2.576<Z<2.576)=0.99$, equivalently $P(\mu-2.576\sigma<X<\mu+2.576\sigma)=0.99$;
\item $P(-1.645<Z<1.645)=0.90$, equivalently $P(\mu-1.645\sigma<X<\mu+1.645\sigma)=0.90$.
\end{itemize}
\emph{Justification.} $P(-z<Z<z)=\Phi(z)-\Phi(-z)=\Phi(z)-\bigl(1-\Phi(z)\bigr)=2\Phi(z)-1$. Setting $2\Phi(z)-1=0.95$ gives $\Phi(z)=0.975$, so $z=1.96$; setting it equal to $0.99$ gives $\Phi(z)=0.995$, so $z=2.576$; setting it equal to $0.90$ gives $\Phi(z)=0.95$, so $z=1.645$. The statement for $X$ follows because $-z<Z<z \iff \mu-z\sigma<X<\mu+z\sigma$.
\end{propriete}

\begin{aretenir}[Percentage points to memorise]
\begin{center}
\begin{tabularx}{\linewidth}{|l|X|X|X|}
\hline
\rowcolor{popblueL} Purpose & Central probability & Tail probability (each) & $z$-value \\
\hline
Central 90\% & $0.90$ & $0.05$ & $1.645$ \\
\hline
Central 95\% & $0.95$ & $0.025$ & $1.96$ \\
\hline
Central 99\% & $0.99$ & $0.005$ & $2.576$ \\
\hline
One-sided 5\% point & $P(Z<z)=0.95$ & $0.05$ (upper) & $1.645$ \\
\hline
One-sided 1\% point & $P(Z<z)=0.99$ & $0.01$ (upper) & $2.326$ \\
\hline
\end{tabularx}
\end{center}
Notice that $1.645$ appears twice: it is the \emph{central} $90\%$ point and the \emph{one-sided} $5\%$ point. The context decides which role it plays.
\end{aretenir}

\begin{attention}[Central 95\% leaves 2.5\% in EACH tail]
A very common error is to use $z=1.645$ for a ``95\% interval''. That is wrong: $1.645$ is the value for which $P(Z<1.645)=0.95$, i.e. it leaves the whole $5\%$ in the \emph{upper} tail only. A \emph{central} 95\% interval must leave $2.5\%$ in \emph{each} tail, so $\Phi(z)=0.975$ and $z=1.96$. Always ask yourself: \emph{central} (two tails) or \emph{one-sided} (one tail)?
\end{attention}

\begin{popfigure}
\begin{center}
\begin{tikzpicture}
\begin{axis}[
  width=13cm, height=6.2cm,
  axis lines=middle,
  xlabel={$z$}, ylabel={$\phi(z)$},
  xmin=-3.6, xmax=3.6, ymin=0, ymax=0.45,
  xtick={-1.96,0,1.96}, xticklabels={$-1.96$,$0$,$1.96$},
  ytick=\empty,
  every axis x label/.style={at={(ticklabel* cs:1.02)}, anchor=west},
  samples=100, domain=-3.6:3.6,
]
\addplot[draw=none, fill=popblueL, domain=-1.96:1.96] {exp(-x^2/2)/sqrt(max(0,2*pi))} \closedcycle;
\addplot[thick, popblue] {exp(-x^2/2)/sqrt(max(0,2*pi))};
\draw[dashed, popdark] (axis cs:-1.96,0) -- (axis cs:-1.96,0.0584);
\draw[dashed, popdark] (axis cs:1.96,0) -- (axis cs:1.96,0.0584);
\node[popdark] at (axis cs:0,0.18) {$0.95$};
\node[popdark, align=center] at (axis cs:-2.9,0.05) {$0.025$};
\node[popdark, align=center] at (axis cs:2.9,0.05) {$0.025$};
\end{axis}
\end{tikzpicture}
\legende{The standard Normal curve: the central 95\% region lies between $z=-1.96$ and $z=1.96$, leaving $2.5\%$ in each tail.}
\end{center}
\end{popfigure}

\begin{savaistu}[The 1.96 legacy]
The value $1.96$ became iconic largely through the work of Sir Ronald Fisher (1925), who recommended the $5\%$ significance level for hypothesis testing. In Cameroon's quality control laboratories -- from the CDC tea plantations in Tiko to the SONARA refinery in Limbe -- the ``$\mu \pm 1.96\sigma$'' interval remains the standard benchmark for setting tolerance limits on product specifications.
\end{savaistu}

\courssub{Using the Standard Normal Tables for Any Normal Distribution}

\begin{methode}[The four-step procedure for probabilities with $X\sim N(\mu,\sigma^{2})$]
\begin{enumerate}
\item \textbf{Standardise.} Replace $X$ by $Z=\dfrac{X-\mu}{\sigma}$ and rewrite the event in terms of $Z$.
\item \textbf{Sketch.} Draw a quick Normal curve, mark $0$ and the $z$-value(s), and shade the required region. The sketch tells you which combination of $\Phi$ values you need.
\item \textbf{Read.} Use the tables of $\Phi$, remembering $\Phi(-z)=1-\Phi(z)$ for negative values.
\item \textbf{Combine.} Use $P(Z<z)=\Phi(z)$, $P(Z>z)=1-\Phi(z)$, $P(a<Z<b)=\Phi(b)-\Phi(a)$, and give the answer to 4 decimal places or 3 significant figures.
\end{enumerate}
\end{methode}

\begin{exoresolu}[Masses of bags of cement]
The mass $X$ of a bag of cement produced by a factory in Figuil is Normally distributed with mean $\SI{50}{kg}$ and standard deviation $\SI{0.8}{kg}$. A bag is chosen at random. Find the probability that its mass is (a) less than $\SI{49.2}{kg}$; (b) more than $\SI{51.1}{kg}$; (c) between $\SI{49.2}{kg}$ and $\SI{51.1}{kg}$; (d) within one interval containing the central $95\%$ of all masses.
\tcblower
We have $X\sim N(50,\,0.8^{2})$, so $Z=\dfrac{X-50}{0.8}$.

\textbf{(a)} Standardise: $z=\dfrac{49.2-50}{0.8}=-1.00$. Hence
\[
P(X<49.2)=P(Z<-1)=\Phi(-1)=1-\Phi(1)=1-0.8413=0.1587.
\]

\textbf{(b)} $z=\dfrac{51.1-50}{0.8}=1.375\approx 1.38$. Hence
\[
P(X>51.1)=1-\Phi(1.375)\approx 1-0.9154=0.0846.
\]
(Using $z=1.38$: $1-0.9162=0.0838$; both are accepted if the rounding is stated.)

\textbf{(c)} Using the two previous results,
\[
P(49.2<X<51.1)=\Phi(1.375)-\Phi(-1)=0.9154-0.1587=0.7567.
\]

\textbf{(d)} The central $95\%$ interval is $\mu\pm 1.96\sigma$:
\[
50\pm 1.96\times 0.8 = 50\pm 1.568,
\]
i.e. between $\SI{48.43}{kg}$ and $\SI{51.57}{kg}$. So $95\%$ of the bags have a mass between about $\SI{48.4}{kg}$ and $\SI{51.6}{kg}$.
\end{exoresolu}

\courssub{De-standardising: From a Probability Back to a Value of $X$}

So far the value $x$ was given and we found the probability. The \emph{inverse} problem is just as important at the GCE: \textbf{the probability is given, and we must find the value of $x$.} The key is to read the tables \emph{backwards}: locate the probability $p$ in the body of the $\Phi$ table and read off the corresponding $z$.

\begin{methode}[De-standardising — from $\Phi(z)=p$ to $x=\mu+z\sigma$]
\begin{enumerate}
\item \textbf{Translate.} Write the given statement in the form $P(X<x)=p$ (if you are given $P(X>x)=p$, rewrite it as $P(X<x)=1-p$).
\item \textbf{Invert the tables.} Find $z$ such that $\Phi(z)=p$. If $p<0.5$, then $z$ is negative: find $z'$ with $\Phi(z')=1-p$ and take $z=-z'$.
\item \textbf{De-standardise.} Compute $x=\mu+z\sigma$.
\item \textbf{Check.} If $p>0.5$ you must get $x>\mu$; if $p<0.5$ you must get $x<\mu$. This one-second check catches most sign errors.
\end{enumerate}
\end{methode}

\begin{popfigure}
\begin{center}
\begin{tikzpicture}[scale=1.0]
% Top panel: X scale
\draw[->, thick] (0,0) -- (10.4,0) node[right] {$X$};
\foreach \x/\lab in {2/{$\mu-1.96\sigma$},5/{$\mu$},8/{$\mu+1.96\sigma$}}{
  \draw (\x,0.12) -- (\x,-0.12);
  \node[below, popdark] at (\x,-0.15) {\lab};
}
\draw[popblue, very thick] (2,0.35) -- (8,0.35);
\fill[popblue] (2,0.35) circle (1.6pt);
\fill[popblue] (8,0.35) circle (1.6pt);
\node[popblue, above] at (5,0.4) {central $95\%$ of the $X$-values};
% arrows
\draw[<->, poporange, thick] (2,-1.1) -- (2,-2.2);
\draw[<->, poporange, thick] (5,-1.1) -- (5,-2.2);
\draw[<->, poporange, thick] (8,-1.1) -- (8,-2.2);
\node[poporange, align=center] at (3.5,-1.65) {$z=\dfrac{x-\mu}{\sigma}$};
\node[poporange, align=center] at (6.9,-1.65) {$x=\mu+z\sigma$};
% Bottom panel: Z scale
\draw[->, thick] (0,-2.6) -- (10.4,-2.6) node[right] {$Z$};
\foreach \x/\lab in {2/{$-1.96$},5/{$0$},8/{$1.96$}}{
  \draw (\x,-2.48) -- (\x,-2.72);
  \node[below, popdark] at (\x,-2.75) {\lab};
}
\draw[popgreen, very thick] (2,-2.25) -- (8,-2.25);
\fill[popgreen] (2,-2.25) circle (1.6pt);
\fill[popgreen] (8,-2.25) circle (1.6pt);
\node[popgreen, below] at (5,-3.35) {central $95\%$ of the $Z$-values};
\end{tikzpicture}
\legende{Standardisation and de-standardisation: the map $z=\frac{x-\mu}{\sigma}$ sends the $X$-scale to the $Z$-scale; the inverse map $x=\mu+z\sigma$ brings a $z$-value back to the original units.}
\end{center}
\end{popfigure}

\begin{exoresolu}[Fill volumes of bottles]
A machine at a brewery in Yaound\'e fills bottles so that the volume $V$ delivered is Normally distributed with mean $\SI{33.0}{cl}$ and standard deviation $\SI{0.4}{cl}$.
(a) Find the volume $v$ such that $90\%$ of bottles contain \emph{more} than $v$.
(b) Find the volume $w$ exceeded by only $1\%$ of bottles.
\tcblower
\textbf{(a)} We want $v$ with $P(V>v)=0.90$, i.e. $P(V<v)=0.10$. Since $0.10<0.5$, the $z$-value is negative. From the tables, $\Phi(1.282)=0.90$, so $z=-1.282$. De-standardising:
\[
v=\mu+z\sigma=33.0+(-1.282)(0.4)=33.0-0.5128\approx \SI{32.49}{cl}.
\]
Check: $p=0.10<0.5$ and indeed $v<\mu=33.0$. \checkmark

\textbf{(b)} We want $w$ with $P(V>w)=0.01$, i.e. $\Phi(z)=0.99$, giving $z=2.326$. Hence
\[
w=33.0+(2.326)(0.4)=33.0+0.9304\approx \SI{33.93}{cl}.
\]
Only $1\%$ of bottles receive more than about $\SI{33.93}{cl}$.
\end{exoresolu}

\courssub{Inverse Problems in Terms of $\mu$ and $\sigma$}

Sometimes the mean and standard deviation are not substituted immediately: the answer is required \emph{in terms of $\mu$ and $\sigma$}, or the information is used to set up an equation (this prepares the next section, where $\mu$ or $\sigma$ is unknown).

\begin{exemplebox}[General inverse statement]
If $X\sim N(\mu,\sigma^{2})$ and $P(X<x)=0.975$, then $\Phi(z)=0.975$ gives $z=1.96$, so
\[
x=\mu+1.96\sigma.
\]
Similarly $P(X<x)=0.05$ gives $z=-1.645$, so $x=\mu-1.645\sigma$; and $P(X>x)=0.005$ gives $\Phi(z)=0.995$, $z=2.576$, so $x=\mu+2.576\sigma$.
\end{exemplebox}

\begin{exoresolu}[Marks of candidates — GCE type]
The marks $X$ obtained by candidates in a mock Advanced Level examination are Normally distributed with mean $52$ and standard deviation $14$.
(a) Find the central interval within which $99\%$ of the marks lie.
(b) The best $10\%$ of candidates are awarded grade A. Find the least mark required for grade A.
(c) Find the mark $m$ such that $P(X<m)=0.20$, and interpret it.
\tcblower
\textbf{(a)} Central $99\%$: $\mu\pm 2.576\sigma = 52\pm 2.576\times 14 = 52\pm 36.06$, i.e. approximately from $15.9$ to $88.1$ marks.

\textbf{(b)} We need $a$ with $P(X>a)=0.10$, i.e. $P(X<a)=0.90$, so $z=1.282$ and
\[
a=52+1.282\times 14=52+17.95\approx 69.9.
\]
The least mark for grade A is therefore about $70$.

\textbf{(c)} $\Phi(z)=0.20$ with $0.20<0.5$: $\Phi(0.842)=0.80$, so $z=-0.842$ and
\[
m=52-0.842\times 14=52-11.79\approx 40.2.
\]
Interpretation: $20\%$ of the candidates scored less than about $40$ marks.
\end{exoresolu}

\begin{attention}[Exam technique: precision of $z$-values]
GCE mark schemes expect you to quote $z$-values to at least 3 significant figures: write $1.96$, $2.576$, $1.645$, $2.326$, and not $1.9$ or $2.6$. When inverting the tables, choose the tabulated probability \emph{closest} to the required one (e.g. use $\Phi(1.96)=0.9750$ for $0.975$). Premature rounding of $z$ before computing $x=\mu+z\sigma$ is a frequent cause of lost accuracy marks.
\end{attention}

\begin{attention}[Signs and direction of the inequality]
Before de-standardising, always rewrite the statement as $P(X<x)=p$ with $p=\Phi(z)$. If you are given $P(X>x)=p$, do \emph{not} look up $p$ directly: look up $1-p$. And remember that a probability below $0.5$ on the left means a \emph{negative} $z$ and a value of $x$ \emph{below} the mean.
\end{attention}

\begin{exercice}[Consolidation]
The heights of students in an Upper Sixth class in Bamenda are Normally distributed with mean $\SI{168}{cm}$ and standard deviation $\SI{7}{cm}$.
\begin{enumerate}
\item Find the probability that a randomly chosen student is taller than $\SI{180}{cm}$.
\item Find the central interval containing the heights of $95\%$ of the students.
\item Find the height $h$ such that only $5\%$ of the students are taller than $h$.
\item Find the height $k$ such that $P(X<k)=0.75$.
\end{enumerate}
\end{exercice}

\begin{corrige}[Exercise 1]
1. $z=\dfrac{180-168}{7}=1.714\approx 1.71$; $P(X>180)=1-\Phi(1.714)\approx 1-0.9568=0.0432$.

2. $\mu\pm 1.96\sigma=168\pm 13.72$, i.e. from $\SI{154.3}{cm}$ to $\SI{181.7}{cm}$.

3. $P(X>h)=0.05\Rightarrow\Phi(z)=0.95\Rightarrow z=1.645$; $h=168+1.645\times 7=\SI{179.5}{cm}$.

4. $\Phi(z)=0.75\Rightarrow z=0.674$; $k=168+0.674\times 7=\SI{172.7}{cm}$.
\end{corrige}

\begin{aretenir}[What you must retain from this section]
\begin{itemize}
\item Central intervals: $90\%\to 1.645$, $95\%\to 1.96$, $99\%\to 2.576$; for $X$ these become $\mu\pm z\sigma$.
\item Forwards: standardise, sketch, read $\Phi$, combine.
\item Backwards (de-standardising): write $P(X<x)=p$, invert the tables to get $z$, then $x=\mu+z\sigma$; check the position of $x$ relative to $\mu$.
\item Central probability $c$ means $(1-c)/2$ in \emph{each} tail: never confuse the central $95\%$ point $1.96$ with the one-sided $5\%$ point $1.645$.
\end{itemize}
\end{aretenir}

These inverse skills — moving freely between a probability and a value of $X$ — are exactly what we shall need in the next section, where the unknown is no longer $x$ but the parameter $\mu$ or $\sigma$ itself.

\coursec{Determining Unknown $\mu$ and/or $\sigma$}

In the previous section we learnt to \emph{use} a Normal distribution whose parameters $\mu$ and $\sigma$ were given: we standardised with $Z=\dfrac{X-\mu}{\sigma}$, read the tables, and de-standardised to answer questions about $X$. In real life, however, the mean and the standard deviation are rarely handed to us on a plate. A quality-control engineer in a soap factory in Douala may know that ``$10\%$ of tablets weigh more than $62$ g'' and ``$5\%$ weigh less than $55$ g'', and from these two pieces of information she must \emph{reconstruct} $\mu$ and $\sigma$. This section is exactly about that reverse problem: we are given probability statements and we must find the unknown parameter(s). The key observation is simple: each probability statement, once converted into a $z$-value by \emph{reading the tables backwards}, produces one linear equation in $\mu$ and $\sigma$. One statement, one equation, one unknown; two statements, two equations, two unknowns.

\courssub{The General Principle: From a Probability Statement to an Equation}

Suppose $X\sim N(\mu,\sigma^2)$. Any statement of the form
\[ P(X<a)=p \qquad\text{or}\qquad P(X>a)=p \]
can be standardised. For instance, $P(X<a)=p$ becomes
\[ P\!\left(Z<\frac{a-\mu}{\sigma}\right)=p, \qquad\text{so}\qquad \Phi\!\left(\frac{a-\mu}{\sigma}\right)=p. \]
Reading the standard normal table \emph{in reverse} (looking for the probability $p$ in the body of the table and reading off the corresponding $z$) gives a number $z_0$ such that $\Phi(z_0)=p$, and hence the \cle{linear equation}
\[ \frac{a-\mu}{\sigma}=z_0 \qquad\Longleftrightarrow\qquad a=\mu+z_0\,\sigma. \]
For a ``greater than'' statement we first convert: $P(X>a)=p$ gives $P(X<a)=1-p$, so $\Phi\!\left(\dfrac{a-\mu}{\sigma}\right)=1-p$.

\begin{attention}[The sign of $z$]
A probability \emph{greater} than $0.5$ corresponds to a \emph{positive} $z$-value; a probability \emph{less} than $0.5$ corresponds to a \emph{negative} $z$-value, and the tables only tabulate $\Phi(z)$ for $z\geq 0$. If $p<0.5$, find $z_1$ with $\Phi(z_1)=1-p$ and take $z=-z_1$ (symmetry: $\Phi(-z_1)=1-\Phi(z_1)$). Getting this sign wrong is the single most common error in this topic. Always ask yourself: is the boundary $a$ above or below the mean?
\end{attention}

\begin{methode}[Setting up equations for unknown $\mu$ and/or $\sigma$]
\begin{enumerate}
\item Write the model: $X\sim N(\mu,\sigma^2)$, stating what is known and what is unknown.
\item Translate each worded probability statement into $P(X<a)=p$ or $P(X>a)=p$.
\item Standardise: $\Phi\!\left(\dfrac{a-\mu}{\sigma}\right)=p$ (use $1-p$ for a ``greater than'' statement).
\item Read the table backwards to find $z_0$ with $\Phi(z_0)=p$; fix the sign using the warning above.
\item Write the equation in the form $\dfrac{a-\mu}{\sigma}=z_0$, i.e. $a=\mu+z_0\sigma$.
\item Solve for the unknown(s); check that $\sigma>0$ and that the answers are sensible in context.
\end{enumerate}
\end{methode}

\begin{popfigure}
\begin{tikzpicture}[node distance=0.55cm, every node/.style={font=\small},
box/.style={draw=popblue, thick, rounded corners, fill=popblueL, text width=4.6cm, align=center, inner sep=5pt},
arr/.style={-{Stealth[length=2.5mm]}, thick, popdark}]
\node[box] (p1) {Probability statement(s): $P(X>a)=p$};
\node[box, below=of p1] (p2) {Standardise: $\Phi\!\left(\frac{a-\mu}{\sigma}\right)=p$ or $1-p$};
\node[box, below=of p2] (p3) {Inverse table reading: $z$-value (mind the sign!)};
\node[box, below=of p3] (p4) {Linear equation(s): $a=\mu+z\sigma$};
\node[box, below=of p4] (p5) {Solve: one unknown $\Rightarrow$ one equation; two unknowns $\Rightarrow$ two equations};
\node[box, below=of p5, fill=popgreenL, draw=popgreen] (p6) {Check: $\sigma>0$, values sensible in context};
\draw[arr] (p1) -- (p2);
\draw[arr] (p2) -- (p3);
\draw[arr] (p3) -- (p4);
\draw[arr] (p4) -- (p5);
\draw[arr] (p5) -- (p6);
\end{tikzpicture}
\legende{Flow-chart of the method: from a probability statement to the values of $\mu$ and $\sigma$.}
\end{popfigure}

\courssub{Type 1: $\mu$ Known, $\sigma$ Unknown}

Here one probability statement provides one equation in the single unknown $\sigma$.

\begin{exoresolu}[Type 1 — Finding $\sigma$]
The masses of loaves produced by a bakery are Normally distributed with mean $500$ g and standard deviation $\sigma$ g. Given that $20\%$ of loaves weigh more than $510$ g, find $\sigma$.
\tcblower
Let $X\sim N(500,\sigma^2)$. The statement ``$20\%$ weigh more than $510$ g'' gives
\[ P(X>510)=0.20 \;\Longrightarrow\; P(X<510)=0.80 \;\Longrightarrow\; \Phi\!\left(\frac{510-500}{\sigma}\right)=0.80. \]
Reading the table backwards, $\Phi(z)=0.80$ gives $z=0.8416$. Since $0.80>0.5$, $z$ is positive (and indeed $510>500$). Hence
\[ \frac{10}{\sigma}=0.8416 \;\Longrightarrow\; \sigma=\frac{10}{0.8416}\approx 11.9\ \text{g}. \]
Check: $\sigma>0$ and of plausible size for a spread of loaf masses. $\boxed{\sigma\approx 11.9\text{ g}}$
\end{exoresolu}

\courssub{Type 2: $\sigma$ Known, $\mu$ Unknown}

Symmetrically, one probability statement gives one equation in $\mu$.

\begin{exoresolu}[Type 2 — Finding $\mu$]
A machine fills sugar packets so that the mass delivered is Normally distributed with standard deviation $4$ g and unknown mean $\mu$ g. Regulations require that only $2.5\%$ of packets contain less than $500$ g. Find the value at which $\mu$ must be set.
\tcblower
Let $X\sim N(\mu,16)$. We need $P(X<500)=0.025$, so
\[ \Phi\!\left(\frac{500-\mu}{4}\right)=0.025. \]
Since $0.025<0.5$, the $z$-value is \emph{negative}. From the tables, $\Phi(1.96)=0.975$, so $\Phi(-1.96)=0.025$, giving
\[ \frac{500-\mu}{4}=-1.96 \;\Longrightarrow\; 500-\mu=-7.84 \;\Longrightarrow\; \mu=507.84\ \text{g}. \]
Check: the mean must exceed $500$ g (so that few packets fall below $500$ g), which it does. $\boxed{\mu\approx 507.8\text{ g}}$
\end{exoresolu}

\begin{attention}[Do not confuse the two directions]
In Type 2 problems the unknown $\mu$ often lies on the \emph{opposite side} of the boundary from what students guess. If only a small percentage may fall \emph{below} a limit, the mean must be \emph{above} that limit. Sketch a quick curve before solving — it will save you from sign errors.
\end{attention}

\courssub{Type 3: Both $\mu$ and $\sigma$ Unknown}

With two unknowns we need \emph{two} probability statements, giving two equations
\[ a_1=\mu+z_1\sigma, \qquad a_2=\mu+z_2\sigma, \]
which we solve simultaneously. The cleanest technique is \cle{elimination}: subtracting the two equations eliminates $\mu$ and yields $\sigma$ immediately; substitution then gives $\mu$.

\begin{methode}[Solving the two simultaneous equations]
\begin{enumerate}
\item Write both equations in the standard form $a_i=\mu+z_i\sigma$ \emph{before} any algebra.
\item Subtract: $a_2-a_1=(z_2-z_1)\sigma$, so $\sigma=\dfrac{a_2-a_1}{z_2-z_1}$.
\item Substitute $\sigma$ into either equation to find $\mu=a_1-z_1\sigma$.
\item Verify: $\sigma$ must be positive; substitute both values back into the \emph{other} equation as a check.
\end{enumerate}
\end{methode}

\begin{popfigure}
\begin{tikzpicture}
\begin{axis}[
width=12cm, height=6.2cm,
axis lines=middle, axis line style={-{Stealth}},
xlabel={$x$}, ylabel={density},
xlabel style={at={(ticklabel* cs:1)}, anchor=west},
ylabel style={at={(ticklabel* cs:1)}, anchor=south},
xtick={55,62}, xticklabels={$a_1=55$,$a_2=62$},
ytick=\empty,
domain=47:71, samples=200,
enlargelimits=false, clip=false]
\addplot[domain=47:55, poporange, fill=poporange, fill opacity=0.45, draw=none] {exp(-((x-59)^2)/(2*3.57^2))/(3.57*2.5066)} \closedcycle;
\addplot[domain=62:71, popgreen, fill=popgreen, fill opacity=0.45, draw=none] {exp(-((x-59)^2)/(2*3.57^2))/(3.57*2.5066)} \closedcycle;
\addplot[popblue, very thick] {exp(-((x-59)^2)/(2*3.57^2))/(3.57*2.5066)};
\draw[dashed, popdark] (axis cs:59,0) -- (axis cs:59,0.112) node[above, font=\small]{$\mu$};
\node[font=\small, popdark] at (axis cs:51.5,0.02) {$P(X<55)=0.1314$};
\node[font=\small, popdark] at (axis cs:66,0.03) {$P(X>62)=0.20$};
\end{axis}
\end{tikzpicture}
\legende{Two tail statements give two equations: $55=\mu-1.12\sigma$ and $62=\mu+0.8416\sigma$.}
\end{popfigure}

\begin{exoresolu}[Type 3 — Finding both $\mu$ and $\sigma$]
The masses of tablets produced by a machine are Normally distributed with mean $\mu$ g and standard deviation $\sigma$ g. It is found that $13.14\%$ of tablets weigh less than $55$ g and $20\%$ weigh more than $62$ g. Find $\mu$ and $\sigma$, and hence the percentage of tablets weighing between $56$ g and $63$ g.
\tcblower
\textbf{Step 1: the two equations.} From $P(X<55)=0.1314$: since $0.1314<0.5$, $z$ is negative; $\Phi(1.12)=0.8686$, so $z_1=-1.12$:
\[ \frac{55-\mu}{\sigma}=-1.12 \;\Longrightarrow\; 55=\mu-1.12\sigma. \tag{1} \]
From $P(X>62)=0.20$: $P(X<62)=0.80$, $\Phi(0.8416)=0.80$, $z_2=+0.8416$:
\[ \frac{62-\mu}{\sigma}=0.8416 \;\Longrightarrow\; 62=\mu+0.8416\sigma. \tag{2} \]
\textbf{Step 2: elimination.} Subtracting (1) from (2):
\[ 7=(0.8416+1.12)\sigma=1.9616\,\sigma \;\Longrightarrow\; \sigma=\frac{7}{1.9616}\approx 3.57\ \text{g}. \]
Substituting into (2): $\mu=62-0.8416(3.57)\approx 62-3.00=59.0$ g.
Check in (1): $59.0-1.12(3.57)=59.0-4.00=55.0$. \checkmark\ Also $\sigma>0$. \checkmark
\textbf{Step 3: further probability with the fitted model.} Now $X\sim N(59,\,3.57^2)$:
\[ P(56<X<63)=\Phi\!\left(\frac{63-59}{3.57}\right)-\Phi\!\left(\frac{56-59}{3.57}\right)=\Phi(1.12)-\Phi(-0.84). \]
Using the tables: $\Phi(1.12)=0.8686$ and $\Phi(-0.84)=1-\Phi(0.84)=1-0.7995=0.2005$. Hence
\[ P(56<X<63)=0.8686-0.2005=0.6681. \]
About $\boxed{66.8\%}$ of tablets weigh between $56$ g and $63$ g.
\end{exoresolu}

\begin{attention}[Exam technique]
State each equation explicitly in the form $\dfrac{x-\mu}{\sigma}=z$ \emph{before} clearing the fraction. Examiners award method marks for this line, and it prevents the classic slip of writing $x=\mu\sigma+z$ or inverting the sign of $z$. Quote $z$-values to 4 decimal places where the tables allow, and keep full precision until the final rounding.
\end{attention}

\courssub{GCE-Style Mixed Problems}

\begin{exoresolu}[Exam marks context]
In a GCE mock examination, the marks obtained by a large group of candidates may be modelled by $X\sim N(\mu,\sigma^2)$. It is known that $15\%$ of candidates scored more than $72$ marks and that $10\%$ scored less than $38$ marks.
\begin{enumerate}
\item Find $\mu$ and $\sigma$.
\item The pass mark is $50$. Estimate the percentage of candidates who passed.
\end{enumerate}
\tcblower
\textbf{(a)} $P(X>72)=0.15\Rightarrow P(X<72)=0.85$, and $\Phi(1.0364)=0.85$, so
\[ 72=\mu+1.0364\,\sigma. \tag{1} \]
$P(X<38)=0.10$: negative $z$; $\Phi(1.2816)=0.90$, so $z=-1.2816$:
\[ 38=\mu-1.2816\,\sigma. \tag{2} \]
Subtracting (2) from (1): $34=2.318\,\sigma$, so $\sigma\approx 14.67$. Then $\mu=72-1.0364(14.67)\approx 72-15.20=56.8$.
Check in (2): $56.8-1.2816(14.67)\approx 56.8-18.8=38.0$. \checkmark
\textbf{(b)} $P(X>50)=P\!\left(Z>\dfrac{50-56.8}{14.67}\right)=P(Z>-0.4635)=\Phi(0.4635)\approx 0.6785$.
About $\boxed{67.9\%}$ of candidates passed.
\end{exoresolu}

\begin{exoresolu}[Crop yield context]
The yield, in kg per plot, of a maize variety is Normally distributed. Experiments show that $5\%$ of plots yield less than $18$ kg and $25\%$ yield more than $30$ kg. Find the mean and standard deviation of the yield.
\tcblower
$P(X<18)=0.05$: $\Phi(1.6449)=0.95$, so $z_1=-1.6449$:
\[ 18=\mu-1.6449\,\sigma. \tag{1} \]
$P(X>30)=0.25\Rightarrow P(X<30)=0.75$, $\Phi(0.6745)=0.75$:
\[ 30=\mu+0.6745\,\sigma. \tag{2} \]
Subtracting: $12=2.3194\,\sigma$, so $\sigma\approx 5.17$ kg, and $\mu=30-0.6745(5.17)\approx 26.5$ kg.
Check in (1): $26.5-1.6449(5.17)\approx 26.5-8.5=18.0$. \checkmark\ $\boxed{\mu\approx 26.5\text{ kg},\ \sigma\approx 5.17\text{ kg}}$
\end{exoresolu}

\begin{exercice}[Practice]
\begin{enumerate}
\item The heights of students in a school are $N(165,\sigma^2)$ cm. If $30\%$ of students are taller than $170$ cm, find $\sigma$.
\item A machine produces bolts whose diameters are $N(\mu,0.04)$ mm (so $\sigma=0.2$ mm). If $1\%$ of bolts have diameter less than $24.5$ mm, find $\mu$.
\item The lifetime $X$ (hours) of a brand of battery is $N(\mu,\sigma^2)$. Given $P(X<40)=0.0668$ and $P(X>70)=0.0228$, find $\mu$ and $\sigma$, and hence $P(50<X<60)$.
\end{enumerate}
\end{exercice}

\begin{corrige}[Exercise 1]
\begin{enumerate}
\item $P(X>170)=0.30\Rightarrow\Phi\!\left(\frac{5}{\sigma}\right)=0.70$, $z=0.5244$, so $\sigma=\dfrac{5}{0.5244}\approx 9.53$ cm.
\item $\Phi\!\left(\frac{24.5-\mu}{0.2}\right)=0.01$, $z=-2.3263$, so $\mu=24.5+0.2(2.3263)\approx 24.97$ mm.
\item $\Phi(1.5)=0.9332\Rightarrow z_1=-1.5$: $40=\mu-1.5\sigma$. $\Phi(2.0)=0.9772\Rightarrow z_2=2.0$: $70=\mu+2\sigma$. Subtracting: $30=3.5\sigma$, $\sigma\approx 8.571$, $\mu=40+1.5(8.571)\approx 52.86$. Then $P(50<X<60)=\Phi(0.8333)-\Phi(-0.3333)=0.7977-0.3694\approx 0.428$.
\end{enumerate}
\end{corrige}

\begin{aretenir}[Key points]
\begin{itemize}
\item Each probability statement $P(X<a)=p$ or $P(X>a)=p$ converts, via standardisation and a \emph{backwards} table reading, into one linear equation $a=\mu+z\sigma$.
\item $p>0.5\Rightarrow z>0$; \ $p<0.5\Rightarrow z<0$ (use symmetry: $\Phi(-z)=1-\Phi(z)$).
\item One unknown requires one statement; two unknowns ($\mu$ and $\sigma$) require two statements, solved by elimination: $\sigma=\dfrac{a_2-a_1}{z_2-z_1}$, then back-substitute for $\mu$.
\item Always check $\sigma>0$ and verify the solution in the unused equation; then the fitted distribution can be used for any further probability question.
\end{itemize}
\end{aretenir}

\coursec{Continuity Correction and Normal Approximations}

In the previous section we learned how to handle any normal distribution by standardising it to $Z \sim N(0,1)$ and using the standard normal tables. We also solved problems where the mean $\mu$ or the standard deviation $\sigma$ (or both) were unknown. Now we turn to a powerful practical tool: using the normal distribution to \emph{approximate} probabilities for \emph{discrete} distributions — the binomial and the Poisson — when exact calculation becomes tedious or impossible without technology. The bridge between the discrete world and the continuous normal curve is the \cle{continuity correction}.

\courssub{Why a Continuity Correction is Needed}

The binomial distribution $B(n,p)$ and the Poisson distribution $\text{Po}(\lambda)$ are \emph{discrete}: the random variable takes only integer values $0,1,2,\dots$. The normal distribution is \emph{continuous}: its probability density function is defined for every real number, and $\Pr(X = x) = 0$ for any single value $x$.

When we approximate a discrete distribution by a continuous one, we must account for the fact that a discrete probability $\Pr(X = r)$ corresponds to an \emph{area} under the continuous curve. The natural way to do this is to associate the integer value $r$ with the interval $[r-0.5,\; r+0.5]$ on the real line. This is the essence of the continuity correction.

\begin{popfigure}
\centering
\begin{tikzpicture}[scale=0.9]
% Axes
\draw[->] (-0.5,0) -- (10.5,0) node[right] {$x$};
\draw[->] (0,-0.3) -- (0,2.2) node[above] {$f(x)$};

% Binomial bars (n=10, p=0.4)
\foreach \x/\h in {0/0.04, 1/0.12, 2/0.21, 3/0.25, 4/0.21, 5/0.12, 6/0.04, 7/0.01, 8/0.002} {
    \draw[fill=popblue!30, draw=popblue] (\x-0.4,0) rectangle (\x+0.4,\h);
    \node[below] at (\x,-0.15) {\small $\x$};
}

% Normal curve (mu=4, sigma=1.55)
\draw[poporange, thick, domain=0:10, samples=200] plot (\x, {1/(1.55*sqrt(max(0,2*pi)))*exp(-(\x-4)^2/(2*1.55^2))});

% Highlight bar for r=4
\draw[fill=popgreen!50, draw=popgreen, thick] (3.6,0) rectangle (4.4,0.21);
\node[above, popgreen] at (4,0.23) {$\Pr(X=4)$};

% Continuity correction interval
\draw[popred, thick, <->] (3.5, -0.5) -- (4.5, -0.5);
\node[popred, below] at (4, -0.65) {$[3.5,\;4.5]$};

% Legend
\node[anchor=north west, text width=5cm, align=left] at (0.5,2.0) {
    \textcolor{popblue}{\rule{4pt}{4pt}} Binomial $B(10,0.4)$ bars\\
    \textcolor{poporange}{\rule{10pt}{2pt}} Normal $N(4,\,2.4)$ curve\\
    \textcolor{popgreen}{\rule{4pt}{4pt}} $\Pr(X=4)$ bar\\
    \textcolor{popred}{\rule{10pt}{2pt}} Continuity correction interval
};
\end{tikzpicture}
\legende{Binomial histogram with superimposed normal curve. The bar at $r=4$ corresponds to the area under the normal curve between $3.5$ and $4.5$.}
\end{popfigure}

\begin{popfigure}
\centering
\begin{tikzpicture}[scale=1.2]
% Number line
\draw[thick, <->] (0,0) -- (10,0);
\foreach \x in {0,1,2,3,4,5,6,7,8,9,10} {
    \draw (\x,0.1) -- (\x,-0.1) node[below] {\small $\x$};
}

% Intervals for each integer
\foreach \r/\col in {0/popblue, 1/poporange, 2/popgreen, 3/poppurple, 4/poppink, 5/popgold, 6/popblue, 7/poporange, 8/popgreen, 9/poppurple, 10/poppink} {
    \draw[\col, thick] (\r-0.5, 0.3) -- (\r+0.5, 0.3);
    \draw[\col] (\r-0.5, 0.25) -- (\r-0.5, 0.35);
    \draw[\col] (\r+0.5, 0.25) -- (\r+0.5, 0.35);
    \node[\col, above] at (\r, 0.4) {\small $[\r-0.5,\,\r+0.5]$};
}

\node[above] at (5, 0.8) {Mapping each integer $r$ to the interval $[r-0.5,\; r+0.5]$};
\end{tikzpicture}
\legende{Number-line strip showing how each discrete value $r$ maps to a continuous interval of width 1 centred on $r$.}
\end{popfigure}

\courssub{The Continuity Correction Rules}

The following table summarises how to translate every discrete probability statement into a continuous one before standardising. \textbf{Memorise this table — it is the single most important tool in this section.}

\begin{methode}[Continuity Correction Table]
For a discrete integer-valued random variable $X$ approximated by a continuous normal variable $Y$:
\begin{center}
\begin{tabularx}{\linewidth}{|l|X|X|}\hline
\rowcolor{popblueL}
\textbf{Discrete probability} & \textbf{Continuous equivalent} & \textbf{Explanation} \\ \hline
$\Pr(X \le r)$ & $\Pr(Y < r + 0.5)$ & Include the whole bar at $r$ \\ \hline
$\Pr(X < r)$ = $\Pr(X \le r-1)$ & $\Pr(Y < r - 0.5)$ & Stop before the bar at $r$ \\ \hline
$\Pr(X = r)$ & $\Pr(r - 0.5 < Y < r + 0.5)$ & Area of the single bar at $r$ \\ \hline
$\Pr(X \ge r)$ & $\Pr(Y > r - 0.5)$ & Include the bar at $r$ and all to the right \\ \hline
$\Pr(X > r)$ = $\Pr(X \ge r+1)$ & $\Pr(Y > r + 0.5)$ & Start after the bar at $r$ \\ \hline
\end{tabularx}
\end{center}
\end{methode}

\begin{attention}[The Most Penalised Error]
\textbf{Forgetting the continuity correction} is the number one mistake in GCE examinations on this topic. Even if you correctly identify the normal approximation, correctly compute $\mu$ and $\sigma$, and correctly use the tables, \textbf{omitting the $\pm 0.5$ adjustment will lose you most of the marks}. Always write the corrected inequality explicitly before standardising.
\end{attention}

\courssub{Normal Approximation to the Binomial Distribution}

When $n$ is large, the binomial distribution $B(n,p)$ becomes symmetric and bell-shaped (provided $p$ is not too close to 0 or 1). The Central Limit Theorem justifies approximating it by a normal distribution with the same mean and variance.

\begin{propriete}[Normal Approximation to the Binomial]
If $X \sim B(n,p)$ with $n$ large, then
\[
X \;\approx\; Y \sim N(\mu,\sigma^2) \quad\text{where}\quad \mu = np,\quad \sigma^2 = npq,\quad q = 1-p.
\]
\textbf{Validity conditions (GCE requirement):} $np > 5$ \textbf{and} $nq > 5$.
\end{propriete}

\begin{savaistu}[Why $np>5$ and $nq>5$?]
These conditions ensure that the binomial distribution is sufficiently spread out and not too skewed. If $p$ is very small, the Poisson approximation (next subsection) is often better. The threshold $5$ is a conventional rule of thumb; some textbooks use $10$. For GCE, always check and \textbf{state} both $np>5$ and $nq>5$ explicitly to earn method marks.
\end{savaistu}

\begin{exemplebox}[Proportion of Defective Items]
A factory produces light bulbs. Historical data shows that $4\%$ of bulbs are defective. A quality control inspector randomly selects $200$ bulbs. Find the probability that \textbf{at most $10$ bulbs are defective}.
\end{exemplebox}

\textbf{Solution.}
Let $X$ be the number of defective bulbs in the sample. Then $X \sim B(200, 0.04)$.
Check conditions: $np = 200 \times 0.04 = 8 > 5$, $nq = 200 \times 0.96 = 192 > 5$. $\checkmark$ Normal approximation is valid.

Approximate by $Y \sim N(\mu,\sigma^2)$ with
\[
\mu = np = 8,\qquad \sigma^2 = npq = 200 \times 0.04 \times 0.96 = 7.68,\qquad \sigma = \sqrt{7.68} \approx 2.771.
\]

We need $\Pr(X \le 10)$. Apply continuity correction: $\Pr(X \le 10) \approx \Pr(Y < 10.5)$.
Standardise:
\[
z = \frac{10.5 - 8}{\sqrt{7.68}} = \frac{2.5}{2.771} \approx 0.902.
\]
Using standard normal tables: $\Phi(0.902) \approx 0.8164$.

\[
\boxed{\Pr(X \le 10) \approx 0.8164}
\]

\begin{exemplebox}[Number of Voters in a Sample]
In a large constituency, $55\%$ of voters support Candidate A. A random sample of $100$ voters is taken. Use a normal approximation to find the probability that \textbf{more than $60$ voters} support Candidate A.
\end{exemplebox}

\textbf{Solution.}
$X \sim B(100, 0.55)$. Check: $np = 55 > 5$, $nq = 45 > 5$. $\checkmark$

$Y \sim N(55,\; 100 \times 0.55 \times 0.45 = 24.75)$, $\sigma = \sqrt{24.75} \approx 4.975$.

$\Pr(X > 60) = \Pr(X \ge 61)$. Continuity correction: $\Pr(Y > 60.5)$.
\[
z = \frac{60.5 - 55}{\sqrt{24.75}} = \frac{5.5}{4.975} \approx 1.105.
\]
$\Pr(Z > 1.105) = 1 - \Phi(1.105) \approx 1 - 0.8654 = 0.1346$.

\[
\boxed{\Pr(X > 60) \approx 0.1346}
\]

\courssub{Normal Approximation to the Poisson Distribution}

For large $\lambda$, the Poisson distribution $\text{Po}(\lambda)$ also becomes approximately normal. This is useful because Poisson tables often stop at $\lambda = 10$ or $15$, and exact calculation of $\Pr(X \ge k)$ for large $\lambda$ involves summing many terms.

\begin{propriete}[Normal Approximation to the Poisson]
If $X \sim \text{Po}(\lambda)$ with $\lambda$ large, then
\[
X \;\approx\; Y \sim N(\mu,\sigma^2) \quad\text{where}\quad \mu = \lambda,\quad \sigma^2 = \lambda.
\]
\textbf{Guideline validity condition:} $\lambda > 10$ (some texts use $\lambda > 15$; for GCE, $\lambda > 10$ is accepted).
\end{propriete}

\begin{exemplebox}[Calls to a Call Centre]
A call centre receives an average of $18$ calls per hour. Assuming a Poisson model, use a normal approximation to find the probability that in a given hour \textbf{more than $25$ calls} are received.
\end{exemplebox}

\textbf{Solution.}
$X \sim \text{Po}(18)$. $\lambda = 18 > 10$, so normal approximation is valid.

$Y \sim N(18,\; 18)$, $\sigma = \sqrt{18} \approx 4.243$.

$\Pr(X > 25) = \Pr(X \ge 26)$. Continuity correction: $\Pr(Y > 25.5)$.
\[
z = \frac{25.5 - 18}{\sqrt{18}} = \frac{7.5}{4.243} \approx 1.768.
\]
$\Pr(Z > 1.768) = 1 - \Phi(1.768) \approx 1 - 0.9615 = 0.0385$.

\[
\boxed{\Pr(X > 25) \approx 0.0385}
\]

\courssub{Choosing the Correct Approximation and Applying the Correction}

In examination questions, you must decide which approximation to use. Here is a decision flowchart:

\begin{popfigure}
\centering
\begin{tikzpicture}[node distance=1.8cm, every node/.style={font=\small}]
\node[draw, rectangle, rounded corners, fill=popblueL, text width=2.5cm, align=center] (start) {Discrete distribution\\$X \sim B(n,p)$ or $\text{Po}(\lambda)$};
\node[draw, diamond, aspect=2, below of=start, fill=poporange!20, text width=2.5cm, align=center] (check) {Is $n$ large?\\($np>5,\;nq>5$)};
\node[draw, rectangle, rounded corners, right of=check, xshift=3.5cm, fill=popgreenL, text width=2.5cm, align=center] (bin) {Use Normal\\$N(np,\,npq)$\\with continuity correction};
\node[draw, diamond, aspect=2, below of=check, yshift=-1.5cm, fill=poporange!20, text width=2.5cm, align=center] (check2) {Is $\lambda$ large?\\($\lambda>10$)};
\node[draw, rectangle, rounded corners, right of=check2, xshift=3.5cm, fill=popgreenL, text width=2.5cm, align=center] (poi) {Use Normal\\$N(\lambda,\,\lambda)$\\with continuity correction};
\node[draw, rectangle, rounded corners, below of=check2, yshift=-1.5cm, fill=poppink!30, text width=2.5cm, align=center] (exact) {Use exact\\binomial/Poisson\\formula or tables};

\draw[->] (start) -- (check);
\draw[->] (check) -- node[left] {Yes} (bin);
\draw[->] (check) -- node[left] {No} (check2);
\draw[->] (check2) -- node[left] {Yes} (poi);
\draw[->] (check2) -- node[left] {No} (exact);
\end{tikzpicture}
\legende{Decision flowchart for choosing the normal approximation.}
\end{popfigure}

\begin{aretenir}[Key Steps for Every Approximation Question]
\begin{enumerate}
\item \textbf{Identify} the exact distribution ($B(n,p)$ or $\text{Po}(\lambda)$) and the required probability.
\item \textbf{Check and state} the validity conditions ($np>5,\;nq>5$ or $\lambda>10$).
\item \textbf{Write down} the approximating normal distribution $N(\mu,\sigma^2)$ with explicit $\mu$ and $\sigma^2$.
\item \textbf{Apply the continuity correction} to convert the discrete inequality to a continuous one. \textbf{Show this step clearly.}
\item \textbf{Standardise} to $z = \frac{x - \mu}{\sigma}$.
\item \textbf{Use the standard normal tables} to find the probability.
\item \textbf{Give the final answer} to 3 or 4 significant figures (or as required).
\end{enumerate}
\end{aretenir}

\begin{exoresolu}[Mixed Approximation Problem]
A machine produces components of which $2\%$ are defective. Components are packed in boxes of $500$.
\begin{enumerate}
\item Use a suitable approximation to find the probability that a randomly chosen box contains \textbf{fewer than $5$ defective components}.
\item The factory produces $2000$ boxes per day. Estimate the number of boxes per day that contain \textbf{at least $15$ defective components}.
\end{enumerate}
\end{exoresolu}

\textbf{Detailed solution.}

\textbf{Part (a).}
$X \sim B(500, 0.02)$. $np = 10 > 5$, $nq = 490 > 5$. Normal approximation valid.
Also $\lambda = np = 10$, so Poisson approximation $\text{Po}(10)$ is possible. Since $\lambda = 10$ is borderline for normal approximation to Poisson, and the question says ``suitable approximation'', the normal approximation to the binomial is the standard route.

$Y \sim N(\mu = 10,\; \sigma^2 = 500 \times 0.02 \times 0.98 = 9.8)$, $\sigma = \sqrt{9.8} \approx 3.1305$.

$\Pr(X < 5) = \Pr(X \le 4)$. Continuity correction: $\Pr(Y < 4.5)$.
\[
z = \frac{4.5 - 10}{\sqrt{9.8}} = \frac{-5.5}{3.1305} \approx -1.757.
\]
$\Pr(Z < -1.757) = 1 - \Phi(1.757) \approx 1 - 0.9605 = 0.0395$.

\[
\boxed{\Pr(X < 5) \approx 0.0395}
\]

\textbf{Part (b).}
We need $\Pr(X \ge 15)$ for one box, then multiply by $2000$.
Continuity correction: $\Pr(X \ge 15) \approx \Pr(Y > 14.5)$.
\[
z = \frac{14.5 - 10}{\sqrt{9.8}} = \frac{4.5}{3.1305} \approx 1.437.
\]
$\Pr(Z > 1.437) = 1 - \Phi(1.437) \approx 1 - 0.9247 = 0.0753$.

Expected number of boxes per day $= 2000 \times 0.0753 = 150.6 \approx \boxed{151 \text{ boxes}}$.

\begin{exoresolu}[Poisson Approximation with Continuity Correction]
The number of accidents per month at a certain intersection follows a Poisson distribution with mean $25$. Use a normal approximation to find the probability that in a given month there are \textbf{between $20$ and $30$ accidents inclusive}.
\end{exoresolu}

\textbf{Detailed solution.}
$X \sim \text{Po}(25)$. $\lambda = 25 > 10$, normal approximation valid.

$Y \sim N(25,\; 25)$, $\sigma = 5$.

Required: $\Pr(20 \le X \le 30)$.
Continuity correction: $\Pr(19.5 < Y < 30.5)$.

Standardise both ends:
\[
z_1 = \frac{19.5 - 25}{5} = -1.1,\qquad z_2 = \frac{30.5 - 25}{5} = 1.1.
\]
$\Pr(-1.1 < Z < 1.1) = \Phi(1.1) - \Phi(-1.1) = \Phi(1.1) - (1 - \Phi(1.1)) = 2\Phi(1.1) - 1$.
From tables: $\Phi(1.1) = 0.8643$.
Probability $= 2 \times 0.8643 - 1 = 0.7286$.

\[
\boxed{\Pr(20 \le X \le 30) \approx 0.7286}
\]

\courssub{Estimating Binomial Probabilities That Are Tedious Exactly}

A classic use of the normal approximation is to estimate tail probabilities like $\Pr(X \ge 40)$ for $X \sim B(100, 0.35)$ without summing $61$ binomial terms.

\begin{exemplebox}[Large Tail Probability]
$X \sim B(100, 0.35)$. Estimate $\Pr(X \ge 40)$ using a normal approximation.
\end{exemplebox}

\textbf{Solution.}
$np = 35 > 5$, $nq = 65 > 5$. Valid.

$Y \sim N(35,\; 100 \times 0.35 \times 0.65 = 22.75)$, $\sigma = \sqrt{22.75} \approx 4.770$.

$\Pr(X \ge 40) \approx \Pr(Y > 39.5)$.
\[
z = \frac{39.5 - 35}{\sqrt{22.75}} = \frac{4.5}{4.770} \approx 0.943.
\]
$\Pr(Z > 0.943) = 1 - \Phi(0.943) \approx 1 - 0.8272 = 0.1728$.

\[
\boxed{\Pr(X \ge 40) \approx 0.1728}
\]

\begin{attention}[Common Traps in Approximation Questions]
\begin{itemize}
\item \textbf{Using the wrong variance:} For binomial, $\sigma^2 = npq$ (not $np$). For Poisson, $\sigma^2 = \lambda$.
\item \textbf{Forgetting to take the square root} when computing $\sigma$ for standardisation.
\item \textbf{Applying continuity correction in the wrong direction:} $\Pr(X \ge r) \to \Pr(Y > r-0.5)$, not $r+0.5$.
\item \textbf{Not stating the validity conditions:} GCE mark schemes allocate specific marks for checking $np>5,\;nq>5$ or $\lambda>10$.
\item \textbf{Using the normal approximation when $p$ is very small or very large:} If $p < 0.1$ and $n$ is large, the Poisson approximation to the binomial ($B(n,p) \approx \text{Po}(np)$) may be more accurate, but the syllabus focuses on Normal approximations here.
\end{itemize}
\end{attention}

\begin{experience}[Comparing Exact and Approximate Probabilities]
Use a calculator or software to compute the exact binomial probability $\Pr(X \ge 40)$ for $X \sim B(100, 0.35)$ and compare with the normal approximation $0.1728$ found above. You will find the exact value is approximately $0.1795$. The error is about $0.0067$, which is typical for a normal approximation with continuity correction. Without the continuity correction ($z = \frac{40-35}{4.770} = 1.048$, probability $0.1474$), the error would be much larger ($0.032$). This demonstrates why the continuity correction is essential.
\end{experience}

\courssub{Examination-Style Practice}

\begin{exercice}[GCE Style Questions]
\begin{enumerate}
\item The probability that a student passes a certain test is $0.7$. In a random sample of $80$ students, use a normal approximation to find the probability that \textbf{more than $60$ students pass}. State the validity conditions clearly.
\item A radioactive source emits particles at an average rate of $12$ per minute. Assuming a Poisson distribution, use a normal approximation to find the probability that in a given minute \textbf{at most $15$ particles} are emitted.
\item $X \sim B(200, 0.12)$. Use a normal approximation to find $\Pr(20 \le X \le 30)$.
\item A factory produces metal rods. The number of defects per rod follows a Poisson distribution with mean $0.8$. A batch contains $500$ rods. Use a suitable approximation to find the probability that the total number of defects in the batch \textbf{exceeds $420$}.
\item In a large population, $30\%$ of people have blood type O. A random sample of $150$ people is taken. Find, using a normal approximation, the probability that the number with blood type O is \textbf{less than $40$}.
\end{enumerate}
\end{exercice}

\begin{corrige}[Exercise 1]
$X \sim B(80, 0.7)$. $np = 56 > 5$, $nq = 24 > 5$. Valid.
$Y \sim N(56,\; 80 \times 0.7 \times 0.3 = 16.8)$, $\sigma = \sqrt{16.8} \approx 4.099$.
$\Pr(X > 60) = \Pr(X \ge 61) \approx \Pr(Y > 60.5)$.
$z = \frac{60.5 - 56}{\sqrt{16.8}} = \frac{4.5}{4.099} \approx 1.098$.
$\Pr(Z > 1.098) = 1 - \Phi(1.098) \approx 1 - 0.8641 = 0.1359$.
\end{corrige}

\begin{corrige}[Exercise 2]
$X \sim \text{Po}(12)$. $\lambda = 12 > 10$. Valid.
$Y \sim N(12,\; 12)$, $\sigma = \sqrt{12} \approx 3.464$.
$\Pr(X \le 15) \approx \Pr(Y < 15.5)$.
$z = \frac{15.5 - 12}{\sqrt{12}} = \frac{3.5}{3.464} \approx 1.010$.
$\Phi(1.010) \approx 0.8438$.
\end{corrige}

\begin{corrige}[Exercise 3]
$X \sim B(200, 0.12)$. $np = 24 > 5$, $nq = 176 > 5$. Valid.
$Y \sim N(24,\; 200 \times 0.12 \times 0.88 = 21.12)$, $\sigma = \sqrt{21.12} \approx 4.596$.
$\Pr(20 \le X \le 30) \approx \Pr(19.5 < Y < 30.5)$.
$z_1 = \frac{19.5 - 24}{\sqrt{21.12}} = \frac{-4.5}{4.596} \approx -0.979$.
$z_2 = \frac{30.5 - 24}{\sqrt{21.12}} = \frac{6.5}{4.596} \approx 1.414$.
$\Pr(-0.979 < Z < 1.414) = \Phi(1.414) - \Phi(-0.979) = \Phi(1.414) - (1 - \Phi(0.979))$.
$\Phi(1.414) \approx 0.9214$, $\Phi(0.979) \approx 0.8363$.
Probability $\approx 0.9214 - 0.1637 = 0.7577$.
\end{corrige}

\begin{corrige}[Exercise 4]
Total defects in 500 rods: $X \sim \text{Po}(500 \times 0.8 = 400)$. $\lambda = 400 > 10$.
$Y \sim N(400,\; 400)$, $\sigma = 20$.
$\Pr(X > 420) \approx \Pr(Y > 420.5)$.
$z = \frac{420.5 - 400}{20} = 1.025$.
$\Pr(Z > 1.025) = 1 - \Phi(1.025) \approx 1 - 0.8473 = 0.1527$.
\end{corrige}

\begin{corrige}[Exercise 5]
$X \sim B(150, 0.3)$. $np = 45 > 5$, $nq = 105 > 5$. Valid.
$Y \sim N(45,\; 150 \times 0.3 \times 0.7 = 31.5)$, $\sigma = \sqrt{31.5} \approx 5.612$.
$\Pr(X < 40) = \Pr(X \le 39) \approx \Pr(Y < 39.5)$.
$z = \frac{39.5 - 45}{\sqrt{31.5}} = \frac{-5.5}{5.612} \approx -0.980$.
$\Pr(Z < -0.980) = 1 - \Phi(0.980) \approx 1 - 0.8365 = 0.1635$.
\end{corrige}

\begin{savaistu}[Historical Note: De Moivre and the Normal Approximation]
Abraham de Moivre (1667–1754), a French mathematician who fled to England, discovered the normal approximation to the binomial in 1733 while studying gambling problems. He showed that for large $n$, the binomial probabilities could be approximated by the curve $y = \frac{1}{\sqrt{2\pi npq}} e^{-\frac{(x-np)^2}{2npq}}$. This was the first appearance of the normal distribution, long before Gauss and Laplace. De Moivre's work laid the foundation for the Central Limit Theorem.
\end{savaistu}

\begin{aretenir}[Summary of Section 6]
\begin{itemize}
\item \textbf{Continuity correction} maps discrete integer $r$ to continuous interval $[r-0.5,\; r+0.5]$.
\item \textbf{Binomial approximation:} $B(n,p) \approx N(np,\,npq)$ if $np>5$ and $nq>5$.
\item \textbf{Poisson approximation:} $\text{Po}(\lambda) \approx N(\lambda,\,\lambda)$ if $\lambda>10$.
\item \textbf{Always:} State conditions $\to$ Write approximating $N(\mu,\sigma^2)$ $\to$ Apply continuity correction $\to$ Standardise $\to$ Use tables.
\item \textbf{Most common error:} Forgetting the $\pm 0.5$ adjustment.
\end{itemize}
\end{aretenir}

\coursec{Combinations of Independent Normal Variables}

In the previous section we mastered the art of \emph{standardising} any normal variable $X \sim N(\mu,\sigma^2)$ to $Z = \frac{X-\mu}{\sigma} \sim N(0,1)$ and using the standard normal tables (or the inverse process) to find probabilities and unknown parameters. We also learned how the normal distribution can approximate the binomial and Poisson distributions when a continuity correction is applied. 

Now we take a decisive step forward: \textbf{what happens when we combine two or more independent normal variables?} This question lies at the heart of many real-world problems — the total weight of a crate and its contents, the difference in performance between two athletes, the combined error of two independent measurements. The beautiful property of the normal family is that \emph{any linear combination of independent normal variables is itself normally distributed}. This section gives you the rigorous tools to handle such combinations, with a special emphasis on the critical distinction between \emph{scaling a single observation} ($nX$) and \emph{summing $n$ independent observations} ($X_1 + \dots + X_n$) — a distinction that examiners love to test.

\courssub{Sum and Difference of Two Independent Normal Variables}

Let us start with the fundamental theorem. Its proof uses moment generating functions (MGFs) or characteristic functions, which are beyond the GCE syllabus, but the statement and its application are examinable and must be known perfectly.

\begin{propriete}[Sum and Difference of Two Independent Normal Variables]
Let $X \sim N(\mu_1,\sigma_1^2)$ and $Y \sim N(\mu_2,\sigma_2^2)$ be \textbf{independent} random variables. Then:
\begin{enumerate}
    \item \textbf{Sum:} $X + Y \sim N(\mu_1 + \mu_2,\; \sigma_1^2 + \sigma_2^2)$.
    \item \textbf{Difference:} $X - Y \sim N(\mu_1 - \mu_2,\; \sigma_1^2 + \sigma_2^2)$.
\end{enumerate}
In words: \emph{means add (or subtract) as expected, but variances always add — even for a difference.}
\end{propriete}

\begin{attention}[Variances ADD for a Difference — Never Subtract!]
A very common error is to write $\mathrm{Var}(X - Y) = \sigma_1^2 - \sigma_2^2$. This is \textbf{wrong}. 
Recall the general rule for independent variables: $\mathrm{Var}(aX + bY) = a^2\sigma_1^2 + b^2\sigma_2^2$. 
For $X - Y$ we have $a=1,\; b=-1$, so the variances add: $(1)^2\sigma_1^2 + (-1)^2\sigma_2^2 = \sigma_1^2 + \sigma_2^2$.
\textbf{Variance measures spread; subtracting two independent variables increases uncertainty, it does not cancel it.}
\end{attention}

\begin{exemplebox}[Total Mass of a Bag and Its Contents]
A bag of rice has mass $R \sim N(5.02,\; 0.04^2)$ kg. The empty sack has mass $S \sim N(0.15,\; 0.01^2)$ kg. $R$ and $S$ are independent. Find the probability that the total mass of a filled sack exceeds $5.20$ kg.

\textbf{Solution:}
Let $T = R + S$ be the total mass. By the theorem:
\[
T \sim N(\mu_T,\sigma_T^2), \quad \mu_T = 5.02 + 0.15 = 5.17,\quad \sigma_T^2 = 0.04^2 + 0.01^2 = 0.0016 + 0.0001 = 0.0017.
\]
So $\sigma_T = \sqrt{0.0017} \approx 0.04123$ kg.
We need $P(T > 5.20)$. Standardise:
\[
z = \frac{5.20 - 5.17}{0.04123} = \frac{0.03}{0.04123} \approx 0.7276.
\]
Using tables: $P(Z > 0.7276) = 1 - \Phi(0.7276) \approx 1 - 0.7664 = 0.2336$.
\textbf{Answer:} $\approx 0.234$ (or $23.4\%$).
\end{exemplebox}

\begin{exemplebox}[Difference Between Two Athletes' Times]
Two sprinters, A and B, run the 100 m. Their times (in seconds) are modelled as independent normal variables:
$A \sim N(10.2,\; 0.15^2)$ and $B \sim N(10.5,\; 0.12^2)$. Find the probability that A beats B by more than $0.2$ seconds.

\textbf{Solution:}
We want $P(A < B - 0.2)$, i.e. $P(B - A > 0.2)$. Let $D = B - A$. 
$D \sim N(\mu_D,\sigma_D^2)$ with $\mu_D = 10.5 - 10.2 = 0.3$ and $\sigma_D^2 = 0.12^2 + 0.15^2 = 0.0144 + 0.0225 = 0.0369$.
$\sigma_D = \sqrt{0.0369} \approx 0.1921$.
\[
P(D > 0.2) = P\left(Z > \frac{0.2 - 0.3}{0.1921}\right) = P(Z > -0.5206) = \Phi(0.5206) \approx 0.6985.
\]
\textbf{Answer:} $\approx 0.699$.
\end{exemplebox}

\begin{methode}[Computing $P(X > Y)$ for Independent Normals]
To find $P(X > Y)$ where $X \sim N(\mu_1,\sigma_1^2)$, $Y \sim N(\mu_2,\sigma_2^2)$ independent:
\begin{enumerate}
    \item Define $D = X - Y$. State $D \sim N(\mu_1 - \mu_2,\; \sigma_1^2 + \sigma_2^2)$.
    \item The required probability is $P(D > 0)$.
    \item Standardise: $z = \frac{0 - (\mu_1 - \mu_2)}{\sqrt{\sigma_1^2 + \sigma_2^2}} = \frac{\mu_2 - \mu_1}{\sqrt{\sigma_1^2 + \sigma_2^2}}$.
    \item Use tables: $P(D > 0) = 1 - \Phi(z)$ if $z>0$, or $\Phi(-z)$ by symmetry.
\end{enumerate}
\end{methode}

\begin{exoresolu}[Probability One Component Outlasts Another]
The lifetimes (in hours) of two types of electronic component are modelled as independent normal variables:
Type X: $N(500,\; 50^2)$, Type Y: $N(480,\; 40^2)$.
A device uses one component of each type. Find the probability that the Type X component outlasts the Type Y component.

\textbf{Statement:} Find $P(X > Y)$.
\tcblower
\textbf{Detailed Solution:}
Let $D = X - Y$. Since $X$ and $Y$ are independent,
\[
D \sim N(\mu_D,\sigma_D^2), \quad \mu_D = 500 - 480 = 20,\quad \sigma_D^2 = 50^2 + 40^2 = 2500 + 1600 = 4100.
\]
$\sigma_D = \sqrt{4100} = 10\sqrt{41} \approx 64.03$.
We need $P(D > 0)$.
\[
z = \frac{0 - 20}{64.03} = -0.3123.
\]
$P(D > 0) = P(Z > -0.3123) = \Phi(0.3123) \approx 0.6225$.
\textbf{Answer:} $0.623$ (3 s.f.).
\end{exoresolu}

\begin{popfigure}
\begin{tikzpicture}
\begin{axis}[
    width=0.95\linewidth,
    height=7cm,
    axis lines=middle,
    xlabel={$x$}, ylabel={Density},
    xmin=-4, xmax=4,
    ymin=0, ymax=0.65,
    xtick={-3,-2,-1,0,1,2,3},
    ytick={0,0.2,0.4,0.6},
    legend style={at={(0.98,0.98)}, anchor=north east, font=\small},
    samples=200,
    domain=-4:4,
    smooth,
]
\addplot[popblue, thick] {1/sqrt(max(0,2*pi))*exp(-x^2/2)};
\addlegendentry{$X \sim N(0,1)$}
\addplot[poporange, thick] {1/sqrt(max(0,2*pi*1.5))*exp(-(x-1)^2/(2*1.5))};
\addlegendentry{$Y \sim N(1,1.5)$}
\addplot[popgreen, thick, dashed] {1/sqrt(max(0,2*pi*2.5))*exp(-(x-1)^2/(2*2.5))};
\addlegendentry{$X+Y \sim N(1,2.5)$}
\addplot[poppurple, thick, dotted] {1/sqrt(max(0,2*pi*2.5))*exp(-(x+1)^2/(2*2.5))};
\addlegendentry{$X-Y \sim N(-1,2.5)$}
\end{axis}
\end{tikzpicture}
\legende{Normal curves for $X\sim N(0,1)$, $Y\sim N(1,1.5)$, their sum $X+Y\sim N(1,2.5)$ and difference $X-Y\sim N(-1,2.5)$. Note how the sum/difference curves are wider (larger variance) and centred at the sum/difference of the means.}
\end{popfigure}

\courssub{Linear Combinations: $aX + bY$ and the Crucial Distinction}

The theorem extends immediately to any linear combination.

\begin{propriete}[Linear Combination of Two Independent Normal Variables]
If $X \sim N(\mu_1,\sigma_1^2)$ and $Y \sim N(\mu_2,\sigma_2^2)$ are independent, then for any constants $a,b \in \mathbb{R}$:
\[
aX + bY \sim N(a\mu_1 + b\mu_2,\; a^2\sigma_1^2 + b^2\sigma_2^2).
\]
\end{propriete}

This formula is the engine for all combination problems. But there is a \textbf{subtle and vital distinction} that separates A-grade candidates from the rest: the difference between \emph{multiplying a single observation by $n$} and \emph{summing $n$ independent observations}.

\begin{aretenir}[The $nX$ vs $X_1 + \dots + X_n$ Trap — Examiners' Favourite]
Let $X \sim N(\mu,\sigma^2)$.
\begin{itemize}
    \item \textbf{One observation scaled by $n$:} $nX \sim N(n\mu,\; n^2\sigma^2)$. \quad Variance is multiplied by $n^2$.
    \item \textbf{Sum of $n$ independent copies:} $X_1 + X_2 + \dots + X_n \sim N(n\mu,\; n\sigma^2)$. \quad Variance is multiplied by $n$.
\end{itemize}
\textbf{The means are the same, but the variances differ by a factor of $n$.} 
$nX$ is \emph{much more variable} than the sum of $n$ independent observations because all the ``eggs are in one basket'' — a single extreme value gets magnified $n$ times. In the sum, extremes tend to cancel out.
\end{aretenir}

\begin{attention}[$3X$ vs $X_1+X_2+X_3$ — Never Confuse Them!]
\begin{itemize}
    \item If a problem says ``the mass of a box is $X$ kg; find the probability that the total mass of \textbf{3 boxes} exceeds...'', you need $X_1+X_2+X_3$ (variance $3\sigma^2$).
    \item If it says ``a box is weighed and the reading is \textbf{multiplied by 3}...'' or ``a single box's mass is $X$; find $P(3X > ...)$'', you use $3X$ (variance $9\sigma^2$).
\end{itemize}
\textbf{Read the wording carefully:} ``three boxes'' $\rightarrow$ sum of three independent variables; ``three times the mass'' $\rightarrow$ multiple of one variable.
\end{attention}

\begin{popfigure}
\begin{tikzpicture}
\begin{axis}[
    width=0.95\linewidth,
    height=7cm,
    axis lines=middle,
    xlabel={$x$}, ylabel={Density},
    xmin=-10, xmax=10,
    ymin=0, ymax=0.5,
    xtick={-8,-4,0,4,8},
    ytick={0,0.1,0.2,0.3,0.4},
    legend style={at={(0.98,0.98)}, anchor=north east, font=\small},
    samples=200,
    domain=-10:10,
    smooth,
]
% X ~ N(0,1)
\addplot[popblue, thick] {1/sqrt(max(0,2*pi))*exp(-x^2/2)};
\addlegendentry{$X \sim N(0,1)$}
% 4X ~ N(0,16)
\addplot[poporange, thick] {1/sqrt(max(0,2*pi*16))*exp(-x^2/(2*16))};
\addlegendentry{$4X \sim N(0,16)$}
% X1+X2+X3+X4 ~ N(0,4)
\addplot[popgreen, thick, dashed] {1/sqrt(max(0,2*pi*4))*exp(-x^2/(2*4))};
\addlegendentry{$X_1+X_2+X_3+X_4 \sim N(0,4)$}
\end{axis}
\end{tikzpicture}
\legende{Comparison of distributions: $X \sim N(0,1)$ (blue), $4X \sim N(0,16)$ (orange, very wide), and $X_1+X_2+X_3+X_4 \sim N(0,4)$ (green dashed, narrower). The sum of four independent copies has variance $4$, while four times a single copy has variance $16$.}
\end{popfigure}

\begin{exemplebox}[Crates of Bottles — Sum vs Multiple]
The mass of a glass bottle is $B \sim N(0.35,\; 0.02^2)$ kg. The mass of an empty crate is $C \sim N(1.2,\; 0.05^2)$ kg. All masses are independent.
\begin{enumerate}
    \item A crate holds 12 bottles. Find the probability that a full crate (crate + 12 bottles) has mass between $5.4$ kg and $5.6$ kg.
    \item A quality control check weighs one bottle and multiplies the reading by 12 to estimate the total bottle mass. Find the probability that this estimate exceeds $4.5$ kg.
\end{enumerate}

\textbf{Solution:}
\begin{enumerate}
    \item Let $T = C + B_1 + \dots + B_{12}$ (sum of 12 independent bottles + crate).
    $\mu_T = 1.2 + 12 \times 0.35 = 1.2 + 4.2 = 5.4$ kg.
    $\sigma_T^2 = 0.05^2 + 12 \times 0.02^2 = 0.0025 + 12 \times 0.0004 = 0.0025 + 0.0048 = 0.0073$.
    $\sigma_T = \sqrt{0.0073} \approx 0.08544$ kg.
    $P(5.4 < T < 5.6) = P\left(\frac{5.4-5.4}{0.08544} < Z < \frac{5.6-5.4}{0.08544}\right) = P(0 < Z < 2.341) = \Phi(2.341) - 0.5 \approx 0.9904 - 0.5 = 0.4904$.
    \item Estimate $E = 12B$ (one bottle scaled by 12).
    $E \sim N(12 \times 0.35,\; 12^2 \times 0.02^2) = N(4.2,\; 144 \times 0.0004) = N(4.2,\; 0.0576)$.
    $\sigma_E = \sqrt{0.0576} = 0.24$ kg.
    $P(E > 4.5) = P\left(Z > \frac{4.5-4.2}{0.24}\right) = P(Z > 1.25) = 1 - \Phi(1.25) \approx 1 - 0.8944 = 0.1056$.
\end{enumerate}
Notice how the estimate $12B$ is far more variable (sd $0.24$ kg) than the actual total of 12 bottles (sd $\approx 0.085$ kg for the bottles alone).
\end{exemplebox}

\courssub{Sums of Several Independent Normal Variables}

The pattern generalises naturally.

\begin{propriete}[Sum of $n$ Independent Normal Variables]
If $X_1, X_2, \dots, X_n$ are independent with $X_i \sim N(\mu_i,\sigma_i^2)$, then
\[
S_n = X_1 + X_2 + \dots + X_n \sim N\left(\sum_{i=1}^n \mu_i,\; \sum_{i=1}^n \sigma_i^2\right).
\]
If they are \emph{identically distributed} (i.i.d.) with $X_i \sim N(\mu,\sigma^2)$, then
\[
S_n \sim N(n\mu,\; n\sigma^2).
\]
\end{propriete}

\begin{exemplebox}[Total Length of Rods]
Metal rods have lengths $L \sim N(1.50,\; 0.02^2)$ m. Five rods are laid end to end. Find the probability that the total length exceeds $7.55$ m.

\textbf{Solution:}
Let $T = L_1 + L_2 + L_3 + L_4 + L_5$ (five independent rods).
$T \sim N(5 \times 1.50,\; 5 \times 0.02^2) = N(7.50,\; 5 \times 0.0004) = N(7.50,\; 0.002)$.
$\sigma_T = \sqrt{0.002} \approx 0.04472$ m.
$P(T > 7.55) = P\left(Z > \frac{7.55 - 7.50}{0.04472}\right) = P(Z > 1.118) = 1 - \Phi(1.118) \approx 1 - 0.8686 = 0.1314$.
\end{exemplebox}

\courssub{Synthesis Problems: Combining with Standardisation and Inverse Tables}

GCE questions often blend this section with earlier skills: finding an unknown $\mu$ or $\sigma$ from a probability statement about a combination, or using the inverse normal tables (percentage points) to determine a threshold.

\begin{exemplebox}[Finding an Unknown Mean from a Combination Probability]
The mass of a standard loaf of bread is $N(\mu,\; 0.01^2)$ kg. A pack contains 6 loaves. The probability that a pack weighs less than $4.80$ kg is $0.025$. Find $\mu$.

\textbf{Solution:}
Let $P = L_1 + \dots + L_6$ be the pack mass. $P \sim N(6\mu,\; 6 \times 0.01^2) = N(6\mu,\; 0.0006)$.
$\sigma_P = \sqrt{0.0006} \approx 0.02449$.
Given $P(P < 4.80) = 0.025$.
From tables, $\Phi^{-1}(0.025) = -1.96$ (the $2.5\%$ lower tail).
Standardise:
\[
\frac{4.80 - 6\mu}{0.02449} = -1.96 \quad \Rightarrow \quad 4.80 - 6\mu = -1.96 \times 0.02449 \approx -0.0480.
\]
\[
6\mu = 4.80 + 0.0480 = 4.848 \quad \Rightarrow \quad \mu = \frac{4.848}{6} = 0.808 \text{ kg}.
\]
\end{exemplebox}

\begin{exemplebox}[Finding an Unknown Standard Deviation — Difference of Two Variables]
Two machines fill jars of jam. Machine A: $N(350,\; \sigma^2)$ g. Machine B: $N(355,\; \sigma^2)$ g. (Same variance $\sigma^2$). Jars from A and B are independent. The probability that a jar from A contains more jam than a jar from B is $0.15$. Find $\sigma$.

\textbf{Solution:}
Let $D = A - B$. $D \sim N(350-355,\; \sigma^2 + \sigma^2) = N(-5,\; 2\sigma^2)$.
We need $P(A > B) = P(D > 0) = 0.15$.
Standardise: $P\left(Z > \frac{0 - (-5)}{\sqrt{2}\sigma}\right) = 0.15 \;\Rightarrow\; P\left(Z > \frac{5}{\sqrt{2}\sigma}\right) = 0.15$.
From tables, the upper $15\%$ point is $z_{0.15} \approx 1.036$ (since $\Phi(1.036) \approx 0.85$).
\[
\frac{5}{\sqrt{2}\sigma} = 1.036 \quad \Rightarrow \quad \sqrt{2}\sigma = \frac{5}{1.036} \approx 4.826 \quad \Rightarrow \quad \sigma = \frac{4.826}{\sqrt{2}} \approx 3.413 \text{ g}.
\]
\end{exemplebox}

\begin{exoresolu}[GCE-Style Synthesis: Crate Weight Limit]
A factory packs apples into crates. The mass of an apple is $A \sim N(150,\; 20^2)$ g. The mass of an empty crate is $C \sim N(500,\; 30^2)$ g. Each crate holds 30 apples. All masses are independent.
\begin{enumerate}
    \item Find the mean and standard deviation of the mass of a full crate.
    \item The factory requires that no more than $2.5\%$ of full crates exceed a weight limit $W$ g. Find $W$.
    \item A quality inspector picks one apple and multiplies its mass by 30 to estimate the total apple mass in a crate. Find the probability that this estimate exceeds the actual total apple mass by more than $200$ g.
\end{enumerate}

\textbf{Statement:} Full crate = crate + 30 apples. Estimate = $30 \times$ (one apple).
\tcblower
\textbf{Detailed Solution:}
\begin{enumerate}
    \item Let $F = C + A_1 + \dots + A_{30}$.
    $\mu_F = 500 + 30 \times 150 = 500 + 4500 = 5000$ g.
    $\sigma_F^2 = 30^2 + 30 \times 20^2 = 900 + 30 \times 400 = 900 + 12000 = 12900$.
    $\sigma_F = \sqrt{12900} \approx 113.58$ g.
    \item We need $P(F > W) = 0.025 \;\Rightarrow\; P(F \le W) = 0.975$.
    The $97.5\%$ point of $N(0,1)$ is $z = 1.96$.
    $\frac{W - 5000}{113.58} = 1.96 \;\Rightarrow\; W = 5000 + 1.96 \times 113.58 \approx 5000 + 222.6 = 5222.6$ g.
    \item Let $E = 30A$ (estimate from one apple). Actual total apple mass $T = A_1 + \dots + A_{30}$.
    $E \sim N(30 \times 150,\; 30^2 \times 20^2) = N(4500,\; 360000)$.
    $T \sim N(30 \times 150,\; 30 \times 20^2) = N(4500,\; 12000)$.
    We need $P(E - T > 200)$. Let $D = E - T$.
    $E$ and $T$ are independent (the single apple is different from the 30 in the crate).
    $D \sim N(4500 - 4500,\; 360000 + 12000) = N(0,\; 372000)$.
    $\sigma_D = \sqrt{372000} \approx 609.9$ g.
    $P(D > 200) = P\left(Z > \frac{200}{609.9}\right) = P(Z > 0.3279) = 1 - \Phi(0.3279) \approx 1 - 0.6285 = 0.3715$.
\end{enumerate}
\textbf{Answers:} (i) Mean $5000$ g, sd $\approx 114$ g. (ii) $W \approx 5223$ g. (iii) $\approx 0.372$.
\end{exoresolu}

\begin{exercice}[Practice: Combinations of Normals]
\begin{enumerate}
    \item $X \sim N(10, 4)$, $Y \sim N(6, 9)$, independent. Find $P(X + Y > 20)$ and $P(X - Y < 2)$.
    \item The height of a man is $M \sim N(175, 7^2)$ cm, a woman $W \sim N(162, 6^2)$ cm, independent. Find the probability that a randomly chosen man is at least 15 cm taller than a randomly chosen woman.
    \item A box contains 20 screws. The mass of a screw is $N(5.0, 0.2^2)$ g. The empty box is $N(50, 1^2)$ g. Find the probability that a full box weighs between $149$ g and $151$ g.
    \item $X \sim N(\mu, 25)$. A sample of 4 independent observations gives $\bar{X} = \frac{X_1+X_2+X_3+X_4}{4}$. Given $P(\bar{X} > 52) = 0.05$, find $\mu$. (Hint: $\bar{X}$ is a linear combination.)
    \item Two independent measurements of a distance are $D_1 \sim N(d, \sigma^2)$ and $D_2 \sim N(d, \sigma^2)$. The average $\frac{D_1+D_2}{2}$ is used as the final estimate. Show that the variance of the estimate is $\sigma^2/2$. If instead a single measurement is multiplied by $1$ (i.e. just $D_1$), compare the variances. Explain why averaging reduces uncertainty.
\end{enumerate}
\end{exercice}

\begin{corrige}[Exercise 1]
$X+Y \sim N(16, 13)$. $P(X+Y>20) = P(Z > \frac{4}{\sqrt{13}}) = P(Z > 1.109) \approx 0.1336$.
$X-Y \sim N(4, 13)$. $P(X-Y<2) = P(Z < \frac{-2}{\sqrt{13}}) = P(Z < -0.555) = \Phi(-0.555) \approx 0.2895$.
\end{corrige}

\begin{corrige}[Exercise 2]
$D = M - W \sim N(13, 7^2+6^2) = N(13, 85)$. $\sigma_D = \sqrt{85} \approx 9.22$.
$P(D \ge 15) = P(Z \ge \frac{2}{9.22}) = P(Z \ge 0.217) \approx 1 - 0.586 = 0.414$.
\end{corrige}

\begin{corrige}[Exercise 3]
Full box $F = B + S_1+\dots+S_{20} \sim N(50+100, 1^2 + 20\times 0.2^2) = N(150, 1+0.8) = N(150, 1.8)$.
$\sigma_F = \sqrt{1.8} \approx 1.3416$.
$P(149 < F < 151) = P(\frac{-1}{1.3416} < Z < \frac{1}{1.3416}) = P(-0.745 < Z < 0.745) = 2\Phi(0.745)-1 \approx 2\times 0.772 - 1 = 0.544$.
\end{corrige}

\begin{corrige}[Exercise 4]
$\bar{X} = \frac{1}{4}(X_1+X_2+X_3+X_4) \sim N(\mu, \frac{25}{4}) = N(\mu, 6.25)$.
$P(\bar{X} > 52) = 0.05 \Rightarrow \frac{52-\mu}{2.5} = 1.645 \Rightarrow 52-\mu = 4.1125 \Rightarrow \mu = 47.8875$.
\end{corrige}

\begin{corrige}[Exercise 5]
Average $A = \frac{D_1+D_2}{2} \sim N(d, \frac{\sigma^2+\sigma^2}{4}) = N(d, \sigma^2/2)$. Variance $= \sigma^2/2$.
Single measurement $D_1 \sim N(d, \sigma^2)$. Variance $= \sigma^2$.
Averaging halves the variance because independent errors tend to cancel out ($\mathrm{Var}(D_1+D_2) = 2\sigma^2$, then dividing by $2$ gives factor $1/4$ on the sum, so $\frac{1}{4}\times 2\sigma^2 = \sigma^2/2$). Using one measurement keeps the full variance $\sigma^2$. Averaging is always better for precision.
\end{corrige}

\begin{savaistu}[Why the Normal Family is Closed Under Linear Combinations]
The normal distribution is a \emph{stable distribution}: any linear combination of independent normals is normal. This is not true for most other distributions (e.g., the sum of two independent exponentials is Gamma, not exponential; the sum of two independent uniforms is triangular, not uniform). This remarkable property makes the normal distribution uniquely tractable for modelling aggregate phenomena — total demand, combined errors, portfolio returns — and is a key reason for its central role in statistics (Central Limit Theorem).
\end{savaistu}

\begin{aretenir}[Key Points for the Exam]
\begin{itemize}
    \item \textbf{Independence is essential.} The theorems only hold if the variables are independent. Always state ``$X$ and $Y$ are independent''.
    \item \textbf{Define the new variable explicitly.} Write ``Let $T = X + Y \sim N(\mu_X+\mu_Y,\; \sigma_X^2+\sigma_Y^2)$'' before calculating. This gains method marks.
    \item \textbf{Variances always add.} For $X-Y$, variance is $\sigma_X^2 + \sigma_Y^2$, never $\sigma_X^2 - \sigma_Y^2$.
    \item \textbf{Distinguish $nX$ from $X_1+\dots+X_n$.} $nX$: variance $n^2\sigma^2$. Sum of $n$ i.i.d.: variance $n\sigma^2$. Read the wording: ``$n$ items'' $\rightarrow$ sum; ``$n$ times the value'' $\rightarrow$ multiple.
    \item \textbf{Standardise correctly.} $z = \frac{x - \mu_{\text{comb}}}{\sigma_{\text{comb}}}$. Use $\sqrt{\text{variance}}$ for $\sigma$.
    \item \textbf{Inverse problems.} If given a probability, find the $z$-value from tables (percentage points), then solve for the unknown $\mu$ or $\sigma$.
\end{itemize}
\end{aretenir}

\begin{experience}[Guided Investigation: Portfolio Risk]
In finance, a portfolio combines assets. Suppose Asset A has return $R_A \sim N(8\%, 15^2)$ and Asset B has return $R_B \sim N(12\%, 20^2)$ (percentages). An investor puts fraction $w$ in A and $(1-w)$ in B. The portfolio return is $R_P = w R_A + (1-w) R_B$.
\begin{enumerate}
    \item Assuming independence, write the distribution of $R_P$ in terms of $w$.
    \item Find the value of $w$ that minimises the variance (risk) of the portfolio.
    \item If the correlation between A and B is actually $\rho = 0.5$, how does the variance formula change? (Formula: $\mathrm{Var}(wX + (1-w)Y) = w^2\sigma_X^2 + (1-w)^2\sigma_Y^2 + 2w(1-w)\rho\sigma_X\sigma_Y$.)
    \item Discuss: why does diversification (mixing assets) reduce risk even when returns are positively correlated?
\end{enumerate}
This illustrates the real-world power of combining normal variables — the foundation of Modern Portfolio Theory (Markowitz, Nobel Prize 1990).
\end{experience}

\coursec{Integration, Assessment and Remediation}

Building on our work with combinations of independent Normal variables, we now consolidate the entire chapter. This section is designed to help you synthesise the three major continuous distributions — Uniform, Exponential and Normal — recognise which one applies to a given problem, practise under exam conditions, identify and correct your most frequent errors, and finally assess your readiness for the end-of-term test.

\courssub{Synthesis Map of the Chapter}

\begin{popfigure}
\begin{tikzpicture}[
  node distance=1.8cm and 2.2cm,
  every node/.style={font=\small},
  dist/.style={draw=popblue, thick, rounded corners, fill=popblueL, text width=2.8cm, align=center, minimum height=1.2cm},
  prop/.style={draw=popgreen, thick, rounded corners, fill=popgreenL, text width=2.5cm, align=center, minimum height=1cm},
  link/.style={draw=popdark, -stealth, thick},
  decision/.style={draw=poporange, thick, diamond, aspect=2, fill=poporange!20, text width=2.5cm, align=center, inner sep=2pt},
  approx/.style={draw=poppurple, thick, rounded corners, fill=poppurple!15, text width=2.5cm, align=center, minimum height=1cm},
]
% Distributions
\node[dist, align=center] (uniform){Uniform (Rectangular)\\ $X\sim U(a,b)$};
\node[dist, right=of uniform, align=center] (expo){Exponential\\ $X\sim Exp(\lambda)$};
\node[dist, right=of expo, align=center] (normal){Normal\\ $X\sim N(\mu,\sigma^2)$};

% Key properties
\node[prop, below=of uniform, align=center] (unifprop){Constant pdf on $[a,b]$\\ $\mathbb{E}(X)=\frac{a+b}{2}$\\ $\mathrm{Var}(X)=\frac{(b-a)^2}{12}$};
\node[prop, below=of expo, align=center] (expoprop){Memoryless property\\ $\mathbb{E}(X)=\frac{1}{\lambda}$\\ $\mathrm{Var}(X)=\frac{1}{\lambda^2}$};
\node[prop, below=of normal, align=center] (normprop){Bell-shaped, symmetric\\ $\mathbb{E}(X)=\mu$\\ $\mathrm{Var}(X)=\sigma^2$};

% Standardisation & Tables
\node[prop, below=of normprop, xshift=1.5cm, align=center] (std){Standard Normal\\ $Z\sim N(0,1)$\\ Tables: $\Phi(z)=P(Z\le z)$};
\node[prop, left=of std, align=center] (destd){De-standardising\\ $X=\mu+\sigma Z$};

% Approximations
\node[approx, below=of destd, xshift=-1.5cm, align=center] (approxbin){Normal approx.\ to\\ Binomial $B(n,p)$\\ if $np>5, nq>5$};
\node[approx, right=of approxbin, align=center] (approxpoi){Normal approx.\ to\\ Poisson $Po(\lambda)$\\ if $\lambda>15$};
\node[approx, right=of approxpoi, align=center] (contcorr){Continuity Correction\\ $P(X\le k)\approx P(Z\le \frac{k+0.5-\mu}{\sigma})$};

% Combinations
\node[prop, below=of approxpoi, yshift=-0.5cm, align=center] (comb){Independent Normals\\ $aX+bY\sim N(a\mu_X+b\mu_Y,\,a^2\sigma_X^2+b^2\sigma_Y^2)$};

% Decision guide
\node[decision, above=of uniform, yshift=1.2cm] (decide) {Which distribution?};

% Links
\draw[link] (decide) -- (uniform);
\draw[link] (decide) -- (expo);
\draw[link] (decide) -- (normal);
\draw[link] (uniform) -- (unifprop);
\draw[link] (expo) -- (expoprop);
\draw[link] (normal) -- (normprop);
\draw[link] (normprop) -- (std);
\draw[link] (std) -- (destd);
\draw[link] (destd) -- (approxbin);
\draw[link] (destd) -- (approxpoi);
\draw[link] (approxbin) -- (contcorr);
\draw[link] (approxpoi) -- (contcorr);
\draw[link] (normprop) -- (comb);
\draw[link] (std) -- (comb);
\end{tikzpicture}
\legende{Concept map linking the three distributions, their properties, standardisation, approximations and combinations.}
\end{popfigure}

The map above summarises the logical flow of the chapter. Use it as a mental checklist: when you read a problem, first decide which distribution fits the context (decision diamond), then follow the relevant branch to the required technique (expectation/variance, standardisation, approximation, or combination).

\begin{methode}[Decision Guide — Recognising the Correct Distribution]
\textbf{Step 1: Identify the random variable and its support.}
\begin{itemize}
\item If the variable is \emph{equally likely} over a finite interval $[a,b]$ (e.g. waiting time for a bus that arrives exactly every 20 minutes, rounding errors), think \cle{Uniform}.
\item If the variable measures \emph{time between random events} occurring at a constant average rate (e.g. time between customer arrivals, radioactive decay, component lifetime), think \cle{Exponential}. Check for the memoryless property: $P(X>s+t\mid X>s)=P(X>t)$.
\item If the variable is a \emph{naturally occurring measurement} that clusters around a mean (e.g. heights, weights, exam marks, measurement errors), think \cle{Normal}.
\end{itemize}
\textbf{Step 2: Check for approximations.}
\begin{itemize}
\item If the problem gives a Binomial $B(n,p)$ with $n$ large and $np>5$, $nq>5$, use \cle{Normal approximation} with continuity correction.
\item If the problem gives a Poisson $Po(\lambda)$ with $\lambda>15$, use \cle{Normal approximation} with continuity correction.
\end{itemize}
\textbf{Step 3: Look for combinations.}
If the problem involves sums or differences of \emph{independent} Normal variables (e.g. total weight of several items, difference in heights), use the linear combination rule: $aX+bY\sim N(a\mu_X+b\mu_Y,\,a^2\sigma_X^2+b^2\sigma_Y^2)$.
\end{methode}

\courssub{Mixed Revision Exercises}

The following exercises cover Lessons 166–182 in increasing order of difficulty. Attempt them without notes first; then check your work against the solutions.

\begin{exercice}[Uniform Distribution — Lesson 166, 167]
A machine produces metal rods whose lengths are uniformly distributed between $9.8$ cm and $10.2$ cm.
\begin{enumerate}
\item Write down the probability density function $f(x)$.
\item Find the probability that a randomly chosen rod has length greater than $10.1$ cm.
\item Calculate the mean and variance of the rod lengths.
\item If $100$ rods are selected independently, estimate the probability that their total length exceeds $1005$ cm. (Hint: use the Central Limit Theorem.)
\end{enumerate}
\end{exercice}

\begin{corrige}[Exercise 1]
\begin{enumerate}
\item $f(x)=\frac{1}{0.4}=2.5$ for $9.8\le x\le 10.2$, $0$ otherwise.
\item $P(X>10.1)=\frac{10.2-10.1}{0.4}=0.25$.
\item $\mathbb{E}(X)=\frac{9.8+10.2}{2}=10.0$ cm. $\mathrm{Var}(X)=\frac{(0.4)^2}{12}=\frac{0.16}{12}\approx 0.01333$ cm$^2$.
\item Let $S=\sum_{i=1}^{100} X_i$. $\mathbb{E}(S)=100\times10.0=1000$ cm. $\mathrm{Var}(S)=100\times\frac{0.16}{12}=\frac{16}{12}\approx 1.333$ cm$^2$. $\sigma_S\approx 1.155$ cm. $P(S>1005)\approx P\left(Z>\frac{1005-1000}{1.155}\right)=P(Z>4.33)\approx 0$. (Virtually impossible.)
\end{enumerate}
\end{corrige}

\begin{exercice}[Exponential Distribution — Lesson 168, 169, 170, 171]
The time between successive arrivals at a hospital emergency unit follows an exponential distribution with mean $6$ minutes.
\begin{enumerate}
\item Find the rate parameter $\lambda$ and write the pdf.
\item What is the probability that the next arrival occurs after $10$ minutes?
\item Given that no arrival has occurred in the first $5$ minutes, what is the probability that the next arrival occurs after $10$ minutes total? (Use the memoryless property.)
\item The number of arrivals in a one-hour period follows a Poisson distribution. State its parameter and find the probability of exactly $12$ arrivals in an hour.
\end{enumerate}
\end{exercice}

\begin{corrige}[Exercise 2]
\begin{enumerate}
\item $\lambda=\frac{1}{6}$ per minute. $f(t)=\frac{1}{6}e^{-t/6},\ t\ge0$.
\item $P(T>10)=e^{-10/6}=e^{-5/3}\approx 0.1889$.
\item By memoryless property, $P(T>10\mid T>5)=P(T>5)=e^{-5/6}\approx 0.4346$.
\item Arrivals per hour $\sim Po(\lambda t)=Po(10)$. $P(X=12)=\frac{e^{-10}10^{12}}{12!}\approx 0.0948$.
\end{enumerate}
\end{corrige}

\begin{exercice}[Normal Distribution — Lessons 172–179]
The masses of bags of cement produced by a factory are normally distributed with mean $50.0$ kg and standard deviation $0.5$ kg.
\begin{enumerate}
\item Find the probability that a randomly chosen bag has mass between $49.5$ kg and $50.5$ kg.
\item Determine the mass exceeded by the heaviest $2.5\%$ of bags.
\item A sample of $25$ bags is selected. Find the probability that the sample mean mass is less than $49.9$ kg.
\item The factory adjusts the machine so that only $1\%$ of bags are under $49.5$ kg, while the standard deviation remains $0.5$ kg. Find the new mean setting.
\end{enumerate}
\end{exercice}

\begin{corrige}[Exercise 3]
\begin{enumerate}
\item $Z_1=\frac{49.5-50}{0.5}=-1$, $Z_2=\frac{50.5-50}{0.5}=1$. $P(-1<Z<1)=\Phi(1)-\Phi(-1)=0.8413-0.1587=0.6826$.
\item Heaviest $2.5\% \Rightarrow P(X>x)=0.025 \Rightarrow P(Z>z)=0.025 \Rightarrow z=1.96$. $x=\mu+z\sigma=50+1.96\times0.5=50.98$ kg.
\item $\bar{X}\sim N(50, \frac{0.5^2}{25}=0.01)$. $P(\bar{X}<49.9)=P\left(Z<\frac{49.9-50}{0.1}\right)=P(Z<-1)=0.1587$.
\item $P(X<49.5)=0.01 \Rightarrow z=-2.3263$ (from tables). $49.5=\mu+(-2.3263)\times0.5 \Rightarrow \mu=49.5+1.16315=50.66315\approx 50.66$ kg.
\end{enumerate}
\end{corrige}

\begin{exercice}[Continuity Correction \& Approximations — Lessons 180, 181]
A fair coin is tossed $200$ times. Use a Normal approximation to find the probability of obtaining:
\begin{enumerate}
\item between $90$ and $110$ heads inclusive;
\item more than $115$ heads.
\end{enumerate}
State the conditions that justify the approximation.
\end{exercice}

\begin{corrige}[Exercise 4]
$X\sim B(200,0.5)$. $np=100>5$, $nq=100>5$ $\Rightarrow$ Normal approximation valid.
$\mu=100$, $\sigma=\sqrt{200\times0.5\times0.5}=\sqrt{50}\approx 7.071$.
\begin{enumerate}
\item $P(90\le X\le 110)\approx P(89.5<X<110.5)=P\left(\frac{89.5-100}{7.071}<Z<\frac{110.5-100}{7.071}\right)=P(-1.485<Z<1.485)=\Phi(1.485)-\Phi(-1.485)\approx 0.9312-0.0688=0.8624$.
\item $P(X>115)\approx P(X>115.5)=P\left(Z>\frac{115.5-100}{7.071}\right)=P(Z>2.192)=1-\Phi(2.192)\approx 1-0.9857=0.0143$.
\end{enumerate}
\end{corrige}

\begin{exercice}[Combinations of Independent Normals — Lesson 182]
The heights of adult males in a certain population are $N(175, 7^2)$ cm, and adult females are $N(162, 6^2)$ cm. A male and a female are chosen independently.
\begin{enumerate}
\item Find the probability that the male is taller than the female.
\item Find the probability that the male is at least $20$ cm taller than the female.
\item If two males and one female are chosen independently, find the probability that the sum of the two males' heights exceeds twice the female's height by more than $30$ cm.
\end{enumerate}
\end{exercice}

\begin{corrige}[Exercise 5]
Let $M\sim N(175,49)$, $F\sim N(162,36)$, independent.
\begin{enumerate}
\item $D=M-F\sim N(175-162, 49+36)=N(13, 85)$. $P(M>F)=P(D>0)=P\left(Z>\frac{0-13}{\sqrt{85}}\right)=P(Z>-1.409)=\Phi(1.409)\approx 0.9207$.
\item $P(M-F\ge 20)=P(D\ge 20)=P\left(Z\ge\frac{20-13}{\sqrt{85}}\right)=P(Z\ge 0.759)=1-\Phi(0.759)\approx 1-0.7764=0.2236$.
\item Let $M_1,M_2$ be the two males. $S=M_1+M_2-2F$. $\mathbb{E}(S)=2\times175-2\times162=26$. $\mathrm{Var}(S)=2\times49+4\times36=98+144=242$. $\sigma_S=\sqrt{242}\approx 15.556$. $P(S>30)=P\left(Z>\frac{30-26}{15.556}\right)=P(Z>0.257)=1-\Phi(0.257)\approx 1-0.6015=0.3985$.
\end{enumerate}
\end{corrige}

\courssub{Timed GCE Past-Paper Questions}

The following questions are typical of GCE Advanced Level Pure Mathematics with Statistics Paper 2. Allow \textbf{15 minutes} for Question 1 (12 marks) and \textbf{12 minutes} for Question 2 (10 marks). Write your answers as you would in the exam: show all steps, state distributions clearly, and use tables where required.

\begin{exoresolu}[GCE Style Question 1 (12 marks)]
The lifetime $T$ (in hours) of a certain electronic component follows an exponential distribution with mean $500$ hours.
\begin{enumerate}
\item Write down the probability density function of $T$. \hfill [1]
\item Find the probability that a component lasts more than $600$ hours. \hfill [2]
\item A system uses $4$ such components connected in series (the system fails if any component fails). Assuming independence, find the probability that the system operates for at least $600$ hours. \hfill [3]
\item The number of components failing in a $1000$-hour period follows a Poisson distribution. State its mean and find the probability that exactly $5$ components fail in this period. \hfill [3]
\item A batch contains $200$ components. Use a Normal approximation to estimate the probability that at least $180$ components are still functioning after $100$ hours. \hfill [3]
\end{enumerate}
\tcblower
\textbf{Detailed Solution}
\begin{enumerate}
\item $\lambda=\frac{1}{500}=0.002$. $f(t)=0.002e^{-0.002t},\ t\ge 0$.
\item $P(T>600)=e^{-0.002\times600}=e^{-1.2}\approx 0.3012$.
\item System lifetime $= \min(T_1,T_2,T_3,T_4)$. For independent exponentials, the minimum is exponential with rate $4\lambda=0.008$. $P(\text{system}>600)=e^{-0.008\times600}=e^{-4.8}\approx 0.00823$. Alternatively, $P(\text{all 4}>600)=[P(T>600)]^4=(0.3012)^4\approx 0.00823$.
\item Failures in $1000$ hours: rate $\lambda t = 0.002\times1000=2$. So $X\sim Po(2)$. $P(X=5)=\frac{e^{-2}2^5}{5!}=\frac{32e^{-2}}{120}\approx 0.0361$.
\item For one component, $P(T>100)=e^{-0.002\times100}=e^{-0.2}\approx 0.8187$. Let $Y$ = number functioning after $100$ hours out of $200$. $Y\sim B(200, 0.8187)$. $np=163.74>5$, $nq=36.26>5$ $\Rightarrow$ Normal approximation. $\mu=163.74$, $\sigma=\sqrt{200\times0.8187\times0.1813}\approx 5.45$. $P(Y\ge 180)\approx P(Y>179.5)=P\left(Z>\frac{179.5-163.74}{5.45}\right)=P(Z>2.89)\approx 1-0.9981=0.0019$.
\end{enumerate}
\end{exoresolu}

\begin{exoresolu}[GCE Style Question 2 (10 marks)]
The random variable $X\sim N(\mu, \sigma^2)$. It is known that $P(X<40)=0.025$ and $P(X>60)=0.10$.
\begin{enumerate}
\item Find the values of $\mu$ and $\sigma$. \hfill [6]
\item Hence find $P(45<X<55)$. \hfill [2]
\item A sample of $36$ observations is taken from this distribution. Find the probability that the sample mean lies between $48$ and $52$. \hfill [2]
\end{enumerate}
\tcblower
\textbf{Detailed Solution}
\begin{enumerate}
\item From tables: $P(Z<z)=0.025 \Rightarrow z=-1.96$. $P(Z>z)=0.10 \Rightarrow z=1.2816$.
Equations:
\[
\frac{40-\mu}{\sigma}=-1.96 \quad\Rightarrow\quad \mu-1.96\sigma=40 \tag{1}
\]
\[
\frac{60-\mu}{\sigma}=1.2816 \quad\Rightarrow\quad \mu+1.2816\sigma=60 \tag{2}
\]
Subtract (1) from (2): $3.2416\sigma=20 \Rightarrow \sigma=\frac{20}{3.2416}\approx 6.169$.
Substitute into (1): $\mu=40+1.96\times6.169\approx 40+12.09=52.09$.
\item $P(45<X<55)=P\left(\frac{45-52.09}{6.169}<Z<\frac{55-52.09}{6.169}\right)=P(-1.15<Z<0.472)=\Phi(0.472)-\Phi(-1.15)\approx 0.6814-0.1251=0.5563$.
\item $\bar{X}\sim N(52.09, \frac{6.169^2}{36}=1.056)$. $\sigma_{\bar{X}}\approx 1.028$.
$P(48<\bar{X}<52)=P\left(\frac{48-52.09}{1.028}<Z<\frac{52-52.09}{1.028}\right)=P(-3.98<Z<-0.0875)\approx \Phi(-0.0875)-\Phi(-3.98)\approx 0.4651-0=0.4651$.
\end{enumerate}
\end{exoresolu}

\courssub{Error-Analysis Workshop}

The following \textbf{five errors} appear repeatedly in GCE scripts. Study each wrong answer, identify the mistake, and write the correct version.

\begin{attention}[Checklist of the Five Most Frequent GCE Errors in This Chapter]
\begin{enumerate}
\item \textbf{Sign error in $z$-score:} Writing $z=\frac{x-\mu}{\sigma}$ but then using $z=\frac{\mu-x}{\sigma}$ or misreading the table for negative $z$. \emph{Always sketch the Normal curve and shade the required area.}
\item \textbf{Missing continuity correction:} Using $P(X\le k)$ directly as $P(Z\le \frac{k-\mu}{\sigma})$ when approximating Binomial/Poisson. \emph{Always add/subtract $0.5$: $P(X\le k)\approx P(Z\le \frac{k+0.5-\mu}{\sigma})$.}
\item \textbf{Variance of a difference:} For independent $X,Y$, writing $\mathrm{Var}(X-Y)=\sigma_X^2-\sigma_Y^2$ instead of $\sigma_X^2+\sigma_Y^2$. \emph{Variances always add for independent variables, even when subtracting.}
\item \textbf{Confusing $\sigma$ and $\sigma^2$ in standardisation:} Using $\sigma^2$ in the denominator: $z=\frac{x-\mu}{\sigma^2}$. \emph{Denominator is standard deviation $\sigma$, not variance.}
\item \textbf{Forgetting to adjust parameters for sample means:} Using $\sigma$ instead of $\frac{\sigma}{\sqrt{n}}$ when dealing with $\bar{X}$. \emph{Standard error $=\frac{\sigma}{\sqrt{n}}$.}
\end{enumerate}
\end{attention}

\begin{experience}[Scripted Wrong Answers — Correct Them]
\textbf{Scenario A:} $X\sim N(50, 25)$. Find $P(X>55)$.
\emph{Wrong script:} $z=\frac{55-50}{25}=0.2$. $P(Z>0.2)=1-0.5793=0.4207$.
\emph{Error:} Used variance $25$ instead of standard deviation $5$.
\emph{Correction:} $z=\frac{55-50}{5}=1$. $P(Z>1)=1-0.8413=0.1587$.

\textbf{Scenario B:} $X\sim B(100, 0.3)$. Approximate $P(X\le 25)$.
\emph{Wrong script:} $\mu=30$, $\sigma=\sqrt{21}\approx 4.58$. $P(X\le 25)\approx P\left(Z\le\frac{25-30}{4.58}\right)=P(Z\le -1.09)=0.1379$.
\emph{Error:} No continuity correction.
\emph{Correction:} $P(X\le 25)\approx P\left(Z\le\frac{25.5-30}{4.58}\right)=P(Z\le -0.982)=0.1635$.

\textbf{Scenario C:} $X\sim N(10,4)$, $Y\sim N(8,9)$ independent. Find $P(X-Y>5)$.
\emph{Wrong script:} $D=X-Y\sim N(2, 4-9=-5)$ — impossible variance.
\emph{Error:} Subtracted variances.
\emph{Correction:} $D\sim N(2, 4+9=13)$. $P(D>5)=P\left(Z>\frac{5-2}{\sqrt{13}}\right)=P(Z>0.832)=0.2027$.
\end{experience}

\courssub{Remediation Tasks}

Based on the error analysis above, the three weakest skills are:
\begin{enumerate}
\item \textbf{Standardisation and correct use of Normal tables (including negative $z$).}
\item \textbf{Applying continuity correction correctly in Normal approximations.}
\item \textbf{Handling combinations of independent Normal variables (especially variance of differences).}
\end{enumerate}

Complete the following targeted tasks. Check your answers against the solutions provided.

\begin{methode}[Remediation Task 1: Standardisation Drill]
For each of the following, find the required probability using Normal tables. Show the $z$-calculation and the table lookup.
\begin{enumerate}
\item $X\sim N(100, 16)$. Find $P(X<95)$.
\item $X\sim N(50, 9)$. Find $P(X>56)$.
\item $X\sim N(20, 4)$. Find $P(18<X<23)$.
\item $X\sim N(\mu, 25)$. Given $P(X<30)=0.9772$, find $\mu$.
\end{enumerate}
\textbf{Answers:} (1) $z=-1.25$, $P=0.1056$. (2) $z=2$, $P=0.0228$. (3) $z_1=-1$, $z_2=1.5$, $P=0.7745$. (4) $z=2$, $\frac{30-\mu}{5}=2\Rightarrow\mu=20$.
\end{methode}

\begin{methode}[Remediation Task 2: Continuity Correction Practice]
For each Binomial/Poisson situation, write the correct Normal approximation statement with continuity correction, then compute the approximate probability.
\begin{enumerate}
\item $X\sim B(50, 0.4)$. Approximate $P(X\ge 22)$.
\item $X\sim Po(20)$. Approximate $P(X<15)$.
\item $X\sim B(200, 0.15)$. Approximate $P(25\le X\le 35)$.
\end{enumerate}
\textbf{Answers:} (1) $\mu=20$, $\sigma=\sqrt{12}\approx 3.464$. $P(X\ge 22)\approx P(Z\ge\frac{21.5-20}{3.464})=P(Z\ge 0.433)=0.3325$. (2) $\mu=20$, $\sigma=\sqrt{20}\approx 4.472$. $P(X<15)\approx P(Z<\frac{14.5-20}{4.472})=P(Z<-1.23)=0.1093$. (3) $\mu=30$, $\sigma=\sqrt{25.5}\approx 5.05$. $P(25\le X\le 35)\approx P(\frac{24.5-30}{5.05}\le Z\le\frac{35.5-30}{5.05})=P(-1.09\le Z\le 1.09)=0.7242$.
\end{methode}

\begin{methode}[Remediation Task 3: Combinations of Normals]
Let $X\sim N(30, 9)$, $Y\sim N(20, 16)$, independent.
\begin{enumerate}
\item Find $P(X+Y > 55)$.
\item Find $P(X-Y < 5)$.
\item Find $P(2X+3Y < 130)$.
\end{enumerate}
\textbf{Answers:} (1) $X+Y\sim N(50, 25)$. $P(Z>\frac{55-50}{5})=P(Z>1)=0.1587$. (2) $X-Y\sim N(10, 25)$. $P(Z<\frac{5-10}{5})=P(Z<-1)=0.1587$. (3) $2X+3Y\sim N(120, 4\times9+9\times16=36+144=180)$. $P(Z<\frac{130-120}{\sqrt{180}})=P(Z<0.745)=0.7722$.
\end{methode}

\courssub{Self-Assessment Checklist}

Use this checklist to evaluate your mastery of the chapter objectives. Tick each item only if you can do it confidently \textbf{without notes} and \textbf{within the time limits} suggested.

\begin{center}
\begin{tabularx}{\linewidth}{|l|X|c|}
\hline
\rowcolor{popblueL}
\textbf{Objective} & \textbf{Description} & \textbf{Mastered?} \\
\hline
1 & Define Uniform distribution, write its pdf, cdf, mean, variance. & $\square$ \\
\hline
2 & Solve problems involving Uniform distribution (probabilities, percentiles). & $\square$ \\
\hline
3 & Define Exponential distribution, state its memoryless property, link to Poisson. & $\square$ \\
\hline
4 & Compute probabilities, mean, variance for Exponential; use memoryless property. & $\square$ \\
\hline
5 & State properties of Normal distribution; standardise to $Z\sim N(0,1)$. & $\square$ \\
\hline
6 & Use Normal tables to find $\Phi(z)$ and $z$-values for given probabilities. & $\square$ \\
\hline
7 & Solve problems for any $N(\mu,\sigma^2)$: probabilities, percentiles, find $\mu$ or $\sigma$. & $\square$ \\
\hline
8 & Apply continuity correction for Normal approximation to Binomial and Poisson. & $\square$ \\
\hline
9 & Check conditions for approximations ($np>5$, $nq>5$, $\lambda>15$). & $\square$ \\
\hline
10 & Handle linear combinations of independent Normal variables (sums, differences, multiples). & $\square$ \\
\hline
11 & Solve problems involving sample means (standard error $\sigma/\sqrt{n}$). & $\square$ \\
\hline
12 & Interpret and solve multi-part GCE questions under timed conditions. & $\square$ \\
\hline
\end{tabularx}
\end{center}

\begin{savaistu}[Preparation for the End-of-Term Test]
\begin{itemize}
\item \textbf{Review the concept map} at the start of this section — redraw it from memory.
\item \textbf{Re-do the two timed GCE questions} (Section 8.3) under strict time limits; mark them using the detailed solutions.
\item \textbf{Complete all three Remediation Tasks} if you ticked fewer than 10 items above.
\item \textbf{Create a one-page "cheat sheet"} (for revision only, not for the exam) containing: key formulas for each distribution, standardisation steps, continuity correction rules, combination rules, and the five common errors from the Warning box.
\item \textbf{Past papers:} Work through at least two full Paper 2 statistics sections from recent GCE sessions (available from your teacher or the GCE Board).
\end{itemize}
\end{savaistu}

\begin{aretenir}[Key Takeaways for Chapter PMS14]
\begin{itemize}
\item \textbf{Uniform:} Constant pdf on $[a,b]$; mean $(a+b)/2$, variance $(b-a)^2/12$.
\item \textbf{Exponential:} Memoryless; mean $1/\lambda$, variance $1/\lambda^2$; inter-arrival times of Poisson process.
\item \textbf{Normal:} Bell-shaped; standardise via $Z=\frac{X-\mu}{\sigma}$; use tables for $\Phi(z)$.
\item \textbf{Approximations:} Binomial $\to$ Normal if $np>5, nq>5$; Poisson $\to$ Normal if $\lambda>15$; \emph{always} use continuity correction.
\item \textbf{Combinations:} If $X,Y$ independent Normals, $aX+bY\sim N(a\mu_X+b\mu_Y,\,a^2\sigma_X^2+b^2\sigma_Y^2)$.
\item \textbf{Exam technique:} Sketch the Normal curve, show $z$-calculations, state approximations used, and check that variances add (never subtract).
\end{itemize}
\end{aretenir}

\newpage

\coursec{Integration activity}

\courssub{Aims of the activity}

This integration activity brings together, in a single authentic logistics case study, all the special continuous distributions studied in this chapter. By the end of the activity, you should be able to:

\begin{itemize}
\item compute the \cle{expectation} and \cle{variance} of a \cle{uniform (rectangular) distribution} and use it to find probabilities as ratios of lengths;
\item exploit the \cle{Poisson--exponential link}: Poisson arrivals in time $t$ correspond to exponential inter-arrival times with the same rate;
\item standardise a Normal variable, use the \cle{standard normal tables} in both directions, and \cle{de-standardise} to find real values;
\item construct the \cle{central 95\% interval} of a Normal distribution;
\item apply the \cle{Normal approximation to the binomial} with a \cle{continuity correction};
\item use the distribution of a \cle{sum of independent Normal variables}.
\end{itemize}

\begin{attention}[Working rules for the activity]
Work in groups of four. Use the standard normal tables provided. Give all probabilities to 4 decimal places and all times to 1 decimal place. State clearly every distribution you use, with its parameters, before computing. Remember: for a continuous variable, $P(X \le a) = P(X < a)$, but for the binomial approximation the continuity correction depends on whether the inequality is strict or not.
\end{attention}

\courssub{Situation}

\textbf{The Buea--Douala Express Agency.}\par
A transport agency operates express buses on the Buea--Douala highway. Over one term, the operations manager collected data on five aspects of the service, and the following models were found to fit the data well.

\begin{enumerate}
\item \textbf{Loading bay.} The time $T$ (in minutes) that a parcel spends on the loading bay before being placed on a bus is uniformly distributed between 5 and 25 minutes: $T \sim U(5, 25)$.
\item \textbf{Ticket office.} Customers arrive at the ticket office according to a Poisson process with rate $3$ customers per 10 minutes. The waiting time $W$ (in minutes) between two consecutive arrivals is therefore exponential.
\item \textbf{Journey time.} The journey time $J$ (in minutes) of a bus from Buea to Douala is Normally distributed with mean $75$ and standard deviation $8$: $J \sim N(75, 8^2)$.
\item \textbf{Passenger satisfaction.} Based on surveys, the probability that any passenger is satisfied with the service is $p = 0.3$, independently of other passengers. In a given week, $300$ passengers are surveyed; let $X$ be the number satisfied.
\item \textbf{Luggage.} The mass of luggage per passenger is Normally distributed with mean $18$ kg and variance $9\ \mathrm{kg}^2$ (standard deviation $3$ kg), independently from passenger to passenger. A bus carries $40$ passengers; let $S$ be the total luggage mass.
\end{enumerate}

The manager wants answers to several operational questions, and your group has been hired as consultants.

\courssub{Tasks}

\textbf{Part A — The loading bay (uniform distribution).}
\begin{enumerate}
\item Write down the probability density function of $T$ and sketch it.
\item Calculate $E(T)$ and $\operatorname{Var}(T)$.
\item Find the probability that a parcel spends more than 20 minutes on the loading bay.
\end{enumerate}

\textbf{Part B — The ticket office (Poisson and exponential distributions).}
\begin{enumerate}
\setcounter{enumi}{3}
\item Show that the rate per minute is $\lambda = 0.3$, and write down the probability density function of the inter-arrival time $W$, together with $E(W)$.
\item Find the probability that the next customer arrives within 2 minutes of the previous one.
\item Using the Poisson distribution, find the probability of exactly 4 arrivals in a 10-minute period.
\item State, in one sentence, the link between the exponential and Poisson distributions illustrated by this situation.
\end{enumerate}

\textbf{Part C — The journey time (Normal distribution).}
\begin{enumerate}
\setcounter{enumi}{7}
\item Find the probability that a journey lasts more than 90 minutes.
\item Determine the central 95\% interval for journey times, i.e.\ the interval within which 95\% of all journeys fall.
\item Find the journey time exceeded by only 5\% of journeys (de-standardising).
\end{enumerate}

\textbf{Part D — Passenger satisfaction (Normal approximation to the binomial).}
\begin{enumerate}
\setcounter{enumi}{10}
\item Explain why $X \sim B(300, 0.3)$ may be approximated by a Normal distribution, and give the parameters of the approximating distribution.
\item Using a continuity correction, estimate the probability that more than 100 of the 300 passengers are satisfied.
\end{enumerate}

\textbf{Part E — Total luggage (sum of independent Normal variables).}
\begin{enumerate}
\setcounter{enumi}{12}
\item State the distribution of the total luggage mass $S$ of the 40 passengers, justifying your answer.
\item Find the probability that the total luggage mass exceeds 750 kg.
\end{enumerate}

\textbf{Part F — Presentation.} Prepare a flipchart with your full solutions. Each step will be validated by the class against the methods of this chapter: correct identification of the model, correct use of tables, correct continuity correction, and correct combination of Normal variables.

\courssub{Model answer}

\textbf{Part A — The loading bay.}\par
1. Since $T \sim U(5,25)$, the probability density function is constant on $[5,25]$ and the total area under it must equal $1$, so its height is $\dfrac{1}{25-5}$:
\[
f(t) = \begin{cases} \dfrac{1}{20} = 0.05, & 5 \le t \le 25, \\[2mm] 0, & \text{otherwise}. \end{cases}
\]

\begin{popfigure}
\begin{tikzpicture}
\draw[->, popmuted] (0,0) -- (7.2,0) node[below left, popink] {$t$ (min)};
\draw[->, popmuted] (0,0) -- (0,2.4) node[below left, popink] {$f(t)$};
\draw[very thick, popblue] (1,1.5) -- (6,1.5);
\draw[dashed, popblue] (1,0) -- (1,1.5);
\draw[dashed, popblue] (6,0) -- (6,1.5);
\draw[dashed, popmuted] (0,1.5) node[left, popink] {$0.05$} -- (1,1.5);
\node[below, popink] at (1,0) {$5$};
\node[below, popink] at (6,0) {$25$};
\fill[popblue, opacity=0.25] (4.6,0) rectangle (6,1.5);
\node[popdark] at (5.3,0.75) {$P(T>20)$};
\end{tikzpicture}
\legende{Density of $T \sim U(5,25)$; the shaded area is $P(T>20) = 5 \times 0.05 = 0.25$.}
\end{popfigure}

2. For $U(a,b)$, $E(T) = \dfrac{a+b}{2}$ and $\operatorname{Var}(T) = \dfrac{(b-a)^2}{12}$:
\[
E(T) = \frac{5+25}{2} = 15 \text{ min}, \qquad \operatorname{Var}(T) = \frac{(25-5)^2}{12} = \frac{400}{12} = \frac{100}{3} \approx 33.3\ \mathrm{min}^2.
\]

3. For a uniform distribution, probabilities are ratios of lengths (areas of rectangles):
\[
P(T > 20) = \frac{25-20}{25-5} = \frac{5}{20} = 0.25.
\]

\textbf{Part B — The ticket office.}\par
4. The rate is 3 customers per 10 minutes, so per minute $\lambda = \dfrac{3}{10} = 0.3$ customers per minute. Hence $W \sim \mathrm{Exp}(0.3)$ with density
\[
f(w) = 0.3\,e^{-0.3 w}, \quad w \ge 0, \qquad E(W) = \frac{1}{\lambda} = \frac{1}{0.3} = \frac{10}{3} \approx 3.3 \text{ min}.
\]
This mean is consistent with the situation: at 3 customers per 10 minutes, one expects on average one customer every $\frac{10}{3}$ minutes.

5. Using the cumulative distribution function $P(W \le w) = 1 - e^{-\lambda w}$:
\[
P(W \le 2) = 1 - e^{-0.3 \times 2} = 1 - e^{-0.6} = 1 - 0.5488 = 0.4512.
\]

6. The number of arrivals in a 10-minute period is $N \sim \mathrm{Po}(\lambda t) = \mathrm{Po}(0.3 \times 10) = \mathrm{Po}(3)$:
\[
P(N = 4) = e^{-3}\,\frac{3^4}{4!} = e^{-3} \times \frac{81}{24} = 0.0498 \times 3.375 = 0.1680.
\]

7. \emph{Link:} if arrivals form a Poisson process with rate $\lambda$ per unit time, then the number of arrivals in a time interval of length $t$ is $\mathrm{Po}(\lambda t)$, and the waiting time between two consecutive arrivals is exponential with parameter $\lambda$ (mean $1/\lambda$).

\textbf{Part C — The journey time.}\par
8. $J \sim N(75, 8^2)$. Standardising with $Z = \dfrac{J - 75}{8} \sim N(0,1)$:
\[
P(J > 90) = P\!\left(Z > \frac{90-75}{8}\right) = P(Z > 1.875) = 1 - \Phi(1.875).
\]
Interpolating between $\Phi(1.87) = 0.9693$ and $\Phi(1.88) = 0.9699$, we get $\Phi(1.875) \approx 0.9696$, so
\[
P(J > 90) = 1 - 0.9696 = 0.0304.
\]

9. The central 95\% interval for $Z$ is $[-1.96,\, 1.96]$, since $\Phi(1.96) = 0.975$ leaves $2.5\%$ in each tail. De-standardising with $j = \mu + z\sigma$:
\[
75 - 1.96 \times 8 = 59.32, \qquad 75 + 1.96 \times 8 = 90.68.
\]
So 95\% of journeys last between approximately $59.3$ and $90.7$ minutes.

10. We need $j$ with $P(J > j) = 0.05$, i.e.\ $P(Z \le z) = 0.95$, giving $z = 1.645$ from the tables. De-standardising:
\[
j = 75 + 1.645 \times 8 = 75 + 13.16 = 88.2 \text{ minutes (1 d.p.).}
\]

\textbf{Part D — Passenger satisfaction.}\par
11. $X \sim B(300, 0.3)$ with $np = 300 \times 0.3 = 90$ and $n(1-p) = 300 \times 0.7 = 210$, both well above $5$, so the Normal approximation is valid:
\[
X \approx Y \sim N\big(np,\, np(1-p)\big) = N(90,\, 63), \qquad \sigma = \sqrt{63} \approx 7.937.
\]

12. ``More than 100'' means $X \ge 101$; the continuity correction replaces the discrete value $101$ by the boundary $100.5$:
\[
P(X > 100) \approx P(Y \ge 100.5) = P\!\left(Z \ge \frac{100.5 - 90}{7.937}\right) = P(Z \ge 1.323).
\]
From the tables, $\Phi(1.32) = 0.9066$ and $\Phi(1.33) = 0.9082$; interpolating, $\Phi(1.323) \approx 0.9071$. Hence
\[
P(X > 100) \approx 1 - 0.9071 = 0.0929.
\]

\textbf{Part E — Total luggage.}\par
13. Each luggage mass $M_i \sim N(18, 9)$ independently. A sum of independent Normal variables is Normal, with the means added and the variances added:
\[
E(S) = 40 \times 18 = 720 \text{ kg}, \qquad \operatorname{Var}(S) = 40 \times 9 = 360, \qquad \sigma_S = \sqrt{360} \approx 18.974 \text{ kg}.
\]
Thus $S \sim N(720, 360)$.

14. Standardising:
\[
P(S > 750) = P\!\left(Z > \frac{750 - 720}{18.974}\right) = P(Z > 1.581).
\]
From the tables, $\Phi(1.58) = 0.9429$, so interpolating $\Phi(1.581) \approx 0.9431$, and
\[
P(S > 750) = 1 - 0.9431 = 0.0569.
\]
There is about a $5.7\%$ chance that the luggage exceeds 750 kg — useful information for the agency's weight policy.

\begin{attention}[Pitfalls to avoid in this type of problem]
\begin{itemize}
\item Do not confuse the \emph{variance} with the \emph{standard deviation}: $N(75, 8^2)$ has $\sigma = 8$, but $N(720, 360)$ has $\sigma = \sqrt{360}$, not $360$.
\item For sums of independent variables, add the \emph{variances}, never the standard deviations.
\item The continuity correction moves the boundary by $0.5$ towards the included region: $X \ge 101$ becomes $Y \ge 100.5$, not $Y \ge 101.5$.
\item In the Poisson--exponential link, the rate $\lambda$ must be expressed in the same time unit as the waiting time.
\end{itemize}
\end{attention}

\begin{aretenir}[What this activity consolidates]
\begin{itemize}
\item Uniform: $E = \frac{a+b}{2}$, $\operatorname{Var} = \frac{(b-a)^2}{12}$, probabilities as length ratios.
\item Exponential: $P(W \le w) = 1 - e^{-\lambda w}$, $E(W) = \frac{1}{\lambda}$; it is the inter-arrival time of a Poisson process of rate $\lambda$.
\item Normal: standardise, read $\Phi$ from tables, de-standardise with $x = \mu + z\sigma$; the central 95\% interval uses $z = \pm 1.96$ and the upper 5\% point uses $z = 1.645$.
\item Binomial $\to$ Normal when $np$ and $n(1-p)$ both exceed 5, always with a continuity correction.
\item Sums of independent Normal variables are Normal: add the means and add the variances.
\end{itemize}
\end{aretenir}

\newpage

\coursec{Methods and worked examples}

\courssub{Method 1: Computing probabilities, mean and variance for a uniform distribution}

\begin{methode}[Working with the rectangular distribution $U(a,b)$]
\begin{enumerate}
\item Identify the interval $[a,b]$ and write the density $f(x)=\dfrac{1}{b-a}$ for $a\le x\le b$, $0$ otherwise.
\item For any sub-interval $[c,d]\subseteq[a,b]$, compute $P(c\le X\le d)=\dfrac{d-c}{b-a}$ (ratio of lengths). A sketch of the rectangle always helps.
\item Write down directly $E(X)=\dfrac{a+b}{2}$ and $\mathrm{Var}(X)=\dfrac{(b-a)^2}{12}$, hence $\sigma=\dfrac{b-a}{2\sqrt{3}}$.
\item For "given that" questions, restrict the interval and recompute the ratio of lengths (conditional probability): the conditional distribution is again uniform on the restricted interval.
\end{enumerate}
\textbf{Reflex:} probabilities for a uniform variable are \emph{areas of rectangles}, i.e. simple ratios of lengths — never integrate unless the question demands a proof.
\end{methode}

\begin{exoresolu}[Waiting for a bus in Bamenda]
A bus arrives at a park in Bamenda at a random time between 8:00 and 8:20 each morning. The waiting time $X$ (in minutes after 8:00) of a passenger who arrives at 8:00 is modelled by $X\sim U(0,20)$.
\begin{enumerate}
\item Write down the probability density function of $X$.
\item Find the probability that the passenger waits more than 12 minutes.
\item Find $E(X)$ and $\mathrm{Var}(X)$.
\item Given that the passenger has already waited 5 minutes, find the probability that he waits at least 10 minutes in total.
\end{enumerate}
\tcblower
\textbf{Solution.}
\begin{enumerate}
\item Here $a=0$, $b=20$, so
\[ f(x)=\begin{cases} \dfrac{1}{20} & 0\le x\le 20,\\[2pt] 0 & \text{otherwise.}\end{cases} \]
\item The favourable interval is $[12,20]$, of length $8$:
\[ P(X>12)=\frac{20-12}{20}=\frac{8}{20}=0.4. \]
\item By the standard formulae:
\[ E(X)=\frac{0+20}{2}=10\ \text{minutes},\qquad \mathrm{Var}(X)=\frac{(20-0)^2}{12}=\frac{400}{12}=\frac{100}{3}\approx 33.33\ \text{min}^2. \]
Hence $\sigma=\sqrt{100/3}\approx 5.77$ minutes.
\item We need $P(X\ge 10\mid X\ge 5)$. Since $X\mid X\ge 5$ is uniform on $[5,20]$ (length $15$):
\[ P(X\ge 10\mid X\ge 5)=\frac{20-10}{20-5}=\frac{10}{15}=\frac{2}{3}\approx 0.667. \]
\end{enumerate}
\end{exoresolu}

\courssub{Method 2: Using the exponential distribution}

\begin{methode}[Probabilities with $X\sim \mathrm{Exp}(\lambda)$]
\begin{enumerate}
\item Write the density $f(x)=\lambda e^{-\lambda x}$, $x\ge 0$, and the cumulative distribution function $F(x)=1-e^{-\lambda x}$.
\item Use the survival formula $P(X>t)=e^{-\lambda t}$ and $P(X\le t)=1-e^{-\lambda t}$.
\item For an interval: $P(a<X<b)=e^{-\lambda a}-e^{-\lambda b}$.
\item Quote $E(X)=\dfrac{1}{\lambda}$, $\mathrm{Var}(X)=\dfrac{1}{\lambda^2}$.
\item Use the \cle{memoryless property}: $P(X>s+t\mid X>s)=P(X>t)=e^{-\lambda t}$.
\item If $\lambda$ is unknown, find it from a given probability (e.g. solve $e^{-\lambda t_0}=p_0$ by taking natural logarithms, $\lambda=-\dfrac{\ln p_0}{t_0}$) or from a given mean ($\lambda=1/E(X)$).
\end{enumerate}
\textbf{Reflex:} check units — if $X$ is in minutes, $\lambda$ is "per minute".
\end{methode}

\begin{exoresolu}[Lifetime of a phone battery component]
The lifetime $T$, in hours, of an electronic component follows an exponential distribution with mean 500 hours.
\begin{enumerate}
\item Write down the value of $\lambda$ and the probability density function of $T$.
\item Find the probability that a component lasts more than 300 hours.
\item Find the probability that a component lasts between 200 and 600 hours.
\item Given that a component has already worked for 300 hours, find the probability that it works for at least 200 more hours. Comment.
\end{enumerate}
\tcblower
\textbf{Solution.}
\begin{enumerate}
\item $E(T)=\dfrac{1}{\lambda}=500 \Rightarrow \lambda=\dfrac{1}{500}=0.002$ per hour. Hence
\[ f(t)=0.002\,e^{-0.002t},\quad t\ge 0. \]
\item Using the survival formula:
\[ P(T>300)=e^{-0.002\times 300}=e^{-0.6}\approx 0.5488. \]
\item 
\[ P(200<T<600)=e^{-0.002(200)}-e^{-0.002(600)}=e^{-0.4}-e^{-1.2}\approx 0.6703-0.3012=0.3691. \]
\item By the memoryless property:
\[ P(T>500\mid T>300)=P(T>200)=e^{-0.002\times 200}=e^{-0.4}\approx 0.6703. \]
\textbf{Comment:} the component "forgets" its age — the probability of surviving 200 extra hours is the same whether it is new or already 300 hours old. This is the defining feature of the exponential model.
\end{enumerate}
\end{exoresolu}

\courssub{Method 3: Linking the exponential and Poisson distributions}

\begin{methode}[From Poisson events to exponential waiting times]
\begin{enumerate}
\item If events occur randomly in time at rate $\lambda$ per unit time, the \emph{number} of events in an interval of length $t$ is $N\sim \mathrm{Po}(\lambda t)$.
\item The \emph{waiting time} $T$ between two consecutive events is $T\sim \mathrm{Exp}(\lambda)$.
\item Key link: $\{T>t\}$ is exactly the event $\{N=0\ \text{in}\ [0,t]\}$, so
\[ P(T>t)=P(N=0)=e^{-\lambda t}. \]
\item To answer questions, decide whether the variable asked about is a \emph{count} (use Poisson) or a \emph{time} (use exponential).
\end{enumerate}
\end{methode}

\begin{exoresolu}[Calls at a call centre in Douala]
Calls arrive at a customer-care centre at an average rate of 4 calls per hour, according to a Poisson process.
\begin{enumerate}
\item Find the probability that exactly 3 calls arrive in one hour.
\item Find the probability that no call arrives in a 30-minute period.
\item Let $T$ be the waiting time, in hours, between two consecutive calls. State the distribution of $T$ and find $P(T>0.5)$.
\item Find the mean and variance of $T$.
\end{enumerate}
\tcblower
\textbf{Solution.}
\begin{enumerate}
\item $N\sim \mathrm{Po}(4)$: $P(N=3)=\dfrac{e^{-4}4^3}{3!}=\dfrac{64e^{-4}}{6}\approx 0.1954$.
\item In 30 minutes, the mean number of calls is $\lambda t=4\times 0.5=2$, so $N\sim\mathrm{Po}(2)$:
\[ P(N=0)=e^{-2}\approx 0.1353. \]
\item $T\sim \mathrm{Exp}(4)$ (rate 4 per hour). Then
\[ P(T>0.5)=e^{-4\times 0.5}=e^{-2}\approx 0.1353. \]
Note this equals the answer to (b): "waiting more than 30 min for the next call" is the same event as "no call in the next 30 min" — the exponential--Poisson link in action.
\item $E(T)=\dfrac{1}{\lambda}=\dfrac{1}{4}$ hour $=15$ minutes; $\mathrm{Var}(T)=\dfrac{1}{\lambda^2}=\dfrac{1}{16}=0.0625$ hour$^2$.
\end{enumerate}
\end{exoresolu}

\courssub{Method 4: Standardising and using the normal tables}

\begin{methode}[Probabilities for $X\sim N(\mu,\sigma^2)$]
\begin{enumerate}
\item Standardise: $Z=\dfrac{X-\mu}{\sigma}$, where $Z\sim N(0,1)$.
\item Rewrite the required probability in terms of $Z$ and express it using $\Phi(z)=P(Z\le z)$, with the properties
\[ \Phi(-z)=1-\Phi(z),\qquad P(Z>z)=1-\Phi(z),\qquad P(a<Z<b)=\Phi(b)-\Phi(a). \]
\item Read $\Phi$ from the standard normal tables (4 decimal places).
\item For a "reverse" question (probability given, value wanted), read the table \emph{backwards}: find the probability in the body of the table and read off $z$, then de-standardise with $x=\mu+\sigma z$.
\end{enumerate}
\textbf{Reflex:} always draw a quick sketch of the bell curve and shade the required region before touching the tables.
\end{methode}

\begin{aretenir}[Special values of $\Phi$ to know]
For the \cle{central intervals}: $P(-1.96<Z<1.96)=0.95$ and $P(-2.576<Z<2.576)=0.99$. Thus the central $95\%$ of any normal distribution lies in $\mu\pm1.96\sigma$, and the central $99\%$ in $\mu\pm2.576\sigma$.
\end{aretenir}

\begin{exoresolu}[Masses of bags of rice]
The mass $X$ (in kg) of bags of rice sold in a market is modelled by $X\sim N(50,\,1.44)$.
\begin{enumerate}
\item Find $P(X<51.2)$.
\item Find $P(X>48.5)$.
\item Find $P(48.5<X<51.2)$.
\item Find the mass $m$ such that $90\%$ of bags are heavier than $m$.
\end{enumerate}
\tcblower
\textbf{Solution.} Here $\mu=50$, $\sigma=\sqrt{1.44}=1.2$.
\begin{enumerate}
\item $Z=\dfrac{51.2-50}{1.2}=1.00$, so $P(X<51.2)=\Phi(1.00)=0.8413$.
\item $Z=\dfrac{48.5-50}{1.2}=-1.25$, so
\[ P(X>48.5)=P(Z>-1.25)=\Phi(1.25)=0.8944. \]
\item 
\[ P(48.5<X<51.2)=\Phi(1.00)-\Phi(-1.25)=0.8413-(1-0.8944)=0.7357. \]
\item We need $P(X>m)=0.90$, i.e. $P(X<m)=0.10$. From the tables, $\Phi(z)=0.10$ gives $z=-1.2816$. De-standardising:
\[ m=\mu+\sigma z=50+1.2(-1.2816)\approx 48.46\ \text{kg}. \]
\end{enumerate}
\end{exoresolu}

\courssub{Method 5: Finding unknown $\mu$ and/or $\sigma$}

\begin{methode}[Determining parameters from given probabilities]
\begin{enumerate}
\item Translate each given probability statement into an equation of the form $\Phi\!\left(\dfrac{x-\mu}{\sigma}\right)=p$.
\item Read the corresponding $z$-value from the tables (take care with signs: if $p<0.5$, then $z$ is negative).
\item Obtain linear equations in $\mu$ and $\sigma$: $x=\mu+z\sigma$.
\item With two statements, solve the simultaneous equations for $\mu$ and $\sigma$; with one statement and one known parameter, solve directly for the other.
\end{enumerate}
\textbf{Reflex:} write each equation as $x-\mu=z\sigma$ \emph{before} substituting numbers — this avoids sign errors.
\end{methode}

\begin{exoresolu}[Filling bottles of palm oil]
A machine fills bottles so that the volume dispensed is normally distributed. It is known that $10\%$ of bottles contain less than $498$ ml and $5\%$ contain more than $506$ ml.
\begin{enumerate}
\item Find the mean $\mu$ and standard deviation $\sigma$ of the volume dispensed.
\item The firm wants only $1\%$ of bottles to contain less than $500$ ml. Keeping $\sigma$ unchanged, to what value should the mean be adjusted?
\end{enumerate}
\tcblower
\textbf{Solution.}
\begin{enumerate}
\item From the tables: $\Phi(z)=0.10 \Rightarrow z=-1.2816$; $\Phi(z)=0.95 \Rightarrow z=1.6449$.
The two statements give:
\[ 498=\mu-1.2816\sigma \quad (1)\qquad 506=\mu+1.6449\sigma \quad (2) \]
Subtracting (1) from (2): $8=2.9265\sigma \Rightarrow \sigma\approx 2.734$ ml.
Then from (1): $\mu=498+1.2816(2.734)\approx 501.50$ ml.
So $\mu\approx 501.5$ ml and $\sigma\approx 2.73$ ml.
\item We require $P(X<500)=0.01$, i.e. $\Phi\!\left(\dfrac{500-\mu}{2.734}\right)=0.01$, so $\dfrac{500-\mu}{2.734}=-2.3263$.
\[ \mu=500+2.3263\times 2.734\approx 506.36\ \text{ml}. \]
The mean should be set to about $506.4$ ml.
\end{enumerate}
\end{exoresolu}

\courssub{Method 6: Continuity correction and normal approximations}

\begin{methode}[Approximating binomial and Poisson distributions]
\begin{enumerate}
\item \textbf{Check validity.} Binomial $B(n,p)$: use $N(np,npq)$ when $n$ is large and $np>5$, $nq>5$. Poisson $\mathrm{Po}(\lambda)$: use $N(\lambda,\lambda)$ when $\lambda$ is large (say $\lambda>10$).
\item \textbf{Apply the continuity correction} (the discrete variable $X$ takes integer values):
\begin{itemize}
\item $P(X\le k)\to P(Y<k+0.5)$; \quad $P(X<k)\to P(Y<k-0.5)$;
\item $P(X\ge k)\to P(Y>k-0.5)$; \quad $P(X>k)\to P(Y>k+0.5)$;
\item $P(X=k)\to P(k-0.5<Y<k+0.5)$.
\end{itemize}
\item Standardise with $\mu=np$ (or $\lambda$) and $\sigma=\sqrt{npq}$ (or $\sqrt{\lambda}$), then use the normal tables.
\end{enumerate}
\end{methode}

\begin{attention}[Forgetting the continuity correction]
Omitting the $\pm0.5$ is the most common error at the GCE. Without it, $P(X=k)$ would be $0$ for a continuous variable! Always ask: "which integers are included?" and place the boundary halfway between included and excluded values.
\end{attention}

\begin{exoresolu}[Defective items in a factory]
A machine produces items of which $8\%$ are defective. A random sample of 200 items is taken. Using a suitable approximation, find the probability that
\begin{enumerate}
\item fewer than 12 items are defective;
\item exactly 16 items are defective.
\end{enumerate}
\tcblower
\textbf{Solution.} Let $X\sim B(200,0.08)$. Then $np=16>5$ and $nq=184>5$, so the normal approximation is valid:
\[ X\approx Y\sim N(16,\;200\times0.08\times0.92)=N(16,\;14.72),\qquad \sigma=\sqrt{14.72}\approx 3.837. \]
\begin{enumerate}
\item "Fewer than 12" means $X\le 11$, corrected to $Y<11.5$:
\[ Z=\frac{11.5-16}{3.837}=-1.173,\qquad P(Y<11.5)=\Phi(-1.173)=1-\Phi(1.173)\approx 1-0.8796=0.1204. \]
\item $P(X=16)$ becomes $P(15.5<Y<16.5)$:
\[ z_1=\frac{15.5-16}{3.837}=-0.130,\qquad z_2=\frac{16.5-16}{3.837}=0.130. \]
\[ P=\Phi(0.130)-\Phi(-0.130)=2\Phi(0.130)-1\approx 2(0.5517)-1=0.1034. \]
\end{enumerate}
\end{exoresolu}

\courssub{Method 7: Combinations of independent normal variables}

\begin{methode}[Sums and differences of independent normal variables]
\begin{enumerate}
\item If $X\sim N(\mu_1,\sigma_1^2)$ and $Y\sim N(\mu_2,\sigma_2^2)$ are \cle{independent}, then
\[ aX+bY\sim N\!\left(a\mu_1+b\mu_2,\;a^2\sigma_1^2+b^2\sigma_2^2\right). \]
In particular $X-Y\sim N(\mu_1-\mu_2,\;\sigma_1^2+\sigma_2^2)$ — \textbf{variances add even for a difference}.
\item For a sum of $n$ independent copies of $X$: $S_n\sim N(n\mu,\,n\sigma^2)$.
\item Compute the new mean and variance, then standardise and use the tables as usual.
\end{enumerate}
\end{methode}

\begin{attention}[Variance of a difference]
$\mathrm{Var}(X-Y)=\sigma_1^2+\sigma_2^2$, \emph{not} $\sigma_1^2-\sigma_2^2$. Subtracting variances is a classic GCE trap.
\end{attention}

\begin{exoresolu}[Heavier boy than girl?]
At a school in Yaound\'e, the masses of boys are $B\sim N(62,25)$ kg and the masses of girls are $G\sim N(54,16)$ kg, independently. A boy and a girl are chosen at random.
\begin{enumerate}
\item Find the probability that the boy is heavier than the girl.
\item Find the probability that their combined mass exceeds $125$ kg.
\end{enumerate}
\tcblower
\textbf{Solution.}
\begin{enumerate}
\item Let $D=B-G$. Then
\[ D\sim N(62-54,\;25+16)=N(8,41),\qquad \sigma_D=\sqrt{41}\approx 6.403. \]
We want $P(D>0)$:
\[ Z=\frac{0-8}{6.403}=-1.249,\qquad P(D>0)=P(Z>-1.249)=\Phi(1.249)\approx 0.8942. \]
\item Let $S=B+G$. Then
\[ S\sim N(62+54,\;25+16)=N(116,41). \]
\[ Z=\frac{125-116}{6.403}=1.406,\qquad P(S>125)=1-\Phi(1.406)\approx 1-0.9201=0.0799. \]
\end{enumerate}
So the boy is heavier with probability about $0.894$, and the combined mass exceeds $125$ kg with probability about $0.080$.
\end{exoresolu}

\newpage

\coursec{Exercises}

\courssub{I check my knowledge}

\begin{exercice}[Uniform distribution basics]
Let $X \sim U(2, 10)$.  
\begin{enumerate}
\item Write down the probability density function $f(x)$.
\item Find $\mathbb{E}(X)$ and $\operatorname{Var}(X)$.
\item Calculate $\mathbb{P}(X > 6)$.
\item Determine the value of $k$ such that $\mathbb{P}(X < k) = 0.75$.
\end{enumerate}
\end{exercice}

\begin{exercice}[Exponential distribution: true/false with justification]
The lifetime $T$ (in months) of a phone battery follows an exponential distribution with mean $30$.  
State whether each of the following is true or false, and justify your answer.
\begin{enumerate}
\item $\mathbb{P}(T > 36) = e^{-1.2}$.
\item Given that the battery has already lasted $24$ months, the probability it lasts at least another $12$ months is $\mathbb{P}(T > 12)$.
\item The median lifetime is $30 \ln 2$ months.
\item The variance of $T$ is $900$.
\end{enumerate}
\end{exercice}

\begin{exercice}[Poisson–Exponential link]
Customers arrive at an MTN service point according to a Poisson process with an average rate of $2$ per minute.  
Let $X$ be the number of arrivals in a $2$-minute interval and $Y$ be the time (in minutes) between two successive arrivals.
\begin{enumerate}
\item Write down the distribution of $X$ and its parameter.
\item Write down the distribution of $Y$ and its parameter.
\item Find $\mathbb{P}(X = 5)$.
\item Find $\mathbb{P}(Y > 1)$.
\end{enumerate}
\end{exercice}

\begin{exercice}[Normal distribution: standardisation]
The marks in a mock GCE examination are normally distributed with mean $52$ and variance $144$.  
Let $X$ be the mark of a randomly chosen candidate.
\begin{enumerate}
\item Find the $z$-score corresponding to $X = 40$.
\item Use standard normal tables to find $\mathbb{P}(X < 40)$.
\item Find $\mathbb{P}(45 < X < 65)$.
\item Determine the mark exceeded by the top $10\%$ of candidates.
\end{enumerate}
\end{exercice}

\courssub{I practise}

\begin{exercice}[Finding parameters of a Normal distribution]
A random variable $X \sim N(\mu, \sigma^2)$ satisfies  
$\mathbb{P}(X < 30) = 0.1587$ and $\mathbb{P}(X > 55) = 0.0228$.  
Determine the values of $\mu$ and $\sigma$.
\end{exercice}

\begin{exercice}[Quality control: adjusting the standard deviation]
A machine fills sugar packets. The masses (in grams) are normally distributed with mean $505$ and standard deviation $4$. Packets with mass below $500$ g are rejected.
\begin{enumerate}
\item Find the percentage of packets rejected.
\item The manufacturer wants to reduce the rejection rate to $1\%$ by improving the machine (keeping $\mu = 505$). Find the required new standard deviation $\sigma$.
\end{enumerate}
\end{exercice}

\begin{exercice}[Normal approximation to the Binomial with continuity correction]
In a large batch of seeds, each seed germinates independently with probability $0.85$. A random sample of $200$ seeds is taken.
\begin{enumerate}
\item State the exact distribution of the number $G$ of seeds that germinate.
\item Explain why a Normal approximation is appropriate.
\item Using a Normal approximation with continuity correction, estimate $\mathbb{P}(G < 160)$.
\end{enumerate}
\end{exercice}

\begin{exercice}[Normal approximation to the Poisson distribution]
The number of accidents per week on a certain highway follows a Poisson distribution with mean $12$.
\begin{enumerate}
\item State the conditions under which a Normal approximation to the Poisson distribution is valid. Are they satisfied here?
\item Using a Normal approximation with continuity correction, estimate the probability that there are at most $8$ accidents in a given week.
\end{enumerate}
\end{exercice}

\courssub{I prepare for the GCE}

\begin{exercice}[{Combining independent Normal variables: filled sacks] [12 marks}]
The mass of a coffee bag (in kg) is normally distributed with mean $50.2$ and variance $0.16$. The mass of an empty sack (in kg) is normally distributed with mean $0.8$ and variance $0.04$. The masses of bags and sacks are independent. A filled sack consists of one bag placed inside one sack.
\begin{enumerate}
\item Find the mean and variance of the mass of a filled sack. [3]
\item Find the probability that a randomly chosen filled sack weighs more than $51.5$ kg. [4]
\item A quality control check selects $4$ filled sacks at random. Find the probability that exactly $2$ of them weigh more than $51.5$ kg. [5]
\end{enumerate}
\end{exercice}

\begin{exercice}[{Two athletes: difference of Normal variables] [10 marks}]
Two athletes, A and B, run the 100 m. Their times (in seconds) are independent and normally distributed as follows:
\[
T_A \sim N(11.2,\; 0.09), \qquad T_B \sim N(11.5,\; 0.16).
\]
In a given race, the athlete with the smaller time wins.
\begin{enumerate}
\item Find the distribution of $D = T_A - T_B$. [3]
\item Find the probability that athlete A beats athlete B. [4]
\item In a series of $5$ independent races, find the probability that A wins at least $3$ races. [3]
\end{enumerate}
\end{exercice}

\begin{exercice}[{Determining unknown parameters and using approximations] [14 marks}]
A factory produces metal rods. The length $L$ (in cm) of a rod is normally distributed with mean $\mu$ and standard deviation $\sigma$. It is known that $2.5\%$ of rods are shorter than $98.0$ cm and $5\%$ are longer than $102.0$ cm.
\begin{enumerate}
\item Find the values of $\mu$ and $\sigma$. [6]
\item Rods are packed in boxes of $50$. A box is accepted if at least $48$ rods have lengths between $98.0$ cm and $102.0$ cm. Using a suitable approximation, estimate the probability that a randomly chosen box is accepted. [8]
\end{enumerate}
\end{exercice}

\newpage

\coursec{Solutions to the exercises}

\coursec{Integration Exercises and Corrections: Special Continuous Distributions}

\begin{corrige}[Exercise 1 — Uniform distribution basics]
We have $X \sim U(2,10)$, so $a = 2$ and $b = 10$, with $b - a = 8$.
\begin{enumerate}
\item The probability density function is
\[
f(x) = \begin{cases} \dfrac{1}{b-a} = \dfrac{1}{8} & \text{if } 2 \le x \le 10,\\[4pt] 0 & \text{otherwise.} \end{cases}
\]
\item Using the standard results $\mathbb{E}(X) = \dfrac{a+b}{2}$ and $\operatorname{Var}(X) = \dfrac{(b-a)^2}{12}$:
\[
\mathbb{E}(X) = \frac{2+10}{2} = 6, \qquad \operatorname{Var}(X) = \frac{(10-2)^2}{12} = \frac{64}{12} = \frac{16}{3} \approx 5.33.
\]
\item For a uniform distribution, probabilities are ratios of lengths:
\[
\mathbb{P}(X > 6) = \frac{10 - 6}{10 - 2} = \frac{4}{8} = \frac{1}{2}.
\]
(This is also immediate from symmetry, since $6$ is the mean.)
\item We require $\mathbb{P}(X < k) = \dfrac{k-2}{8} = 0.75$, so
\[
k - 2 = 8 \times 0.75 = 6 \quad\Longrightarrow\quad k = 8.
\]
\end{enumerate}
\end{corrige}

\begin{corrige}[Exercise 2 — Exponential distribution: true/false with justification]
The mean is $30$ months, so the rate parameter is $\lambda = \dfrac{1}{30}$ and the cumulative distribution function is $\mathbb{P}(T \le t) = 1 - e^{-t/30}$, hence $\mathbb{P}(T > t) = e^{-t/30}$.
\begin{enumerate}
\item \textbf{True.} $\mathbb{P}(T > 36) = e^{-36/30} = e^{-1.2}$.
\item \textbf{True.} By the \cle{memoryless property} of the exponential distribution,
\[
\mathbb{P}(T > 24 + 12 \mid T > 24) = \frac{\mathbb{P}(T > 36)}{\mathbb{P}(T > 24)} = \frac{e^{-36/30}}{e^{-24/30}} = e^{-12/30} = \mathbb{P}(T > 12).
\]
\item \textbf{True.} The median $m$ satisfies $\mathbb{P}(T > m) = \tfrac{1}{2}$, i.e. $e^{-m/30} = \tfrac{1}{2}$, so $\dfrac{m}{30} = \ln 2$ and $m = 30\ln 2 \approx 20.8$ months.
\item \textbf{True.} For an exponential distribution, $\operatorname{Var}(T) = \dfrac{1}{\lambda^2} = 30^2 = 900\ \mathrm{months}^2$.
\end{enumerate}
All four statements are true.
\end{corrige}

\begin{corrige}[Exercise 3 — Poisson–Exponential link]
\begin{enumerate}
\item The number of arrivals in a $2$-minute interval follows a Poisson distribution with mean $2 \times 2 = 4$:
\[
X \sim \mathrm{Po}(4).
\]
\item The waiting time between two successive arrivals in a Poisson process of rate $\lambda = 2$ per minute is exponential with the same rate:
\[
Y \sim \mathrm{Exp}(2), \qquad \mathbb{E}(Y) = \tfrac{1}{2}\ \mathrm{minute}.
\]
\item $\mathbb{P}(X = 5) = \dfrac{e^{-4}\,4^5}{5!} = \dfrac{e^{-4} \times 1024}{120} \approx 0.1563$.
\item $\mathbb{P}(Y > 1) = e^{-\lambda \times 1} = e^{-2} \approx 0.1353$.
\end{enumerate}
Note the link: $\mathbb{P}(Y > 1)$ (no arrival during $1$ minute) equals $\mathbb{P}(N = 0)$ where $N \sim \mathrm{Po}(2)$, and indeed $\mathbb{P}(N=0) = e^{-2}$.
\end{corrige}

\begin{corrige}[Exercise 4 — Normal distribution: standardisation]
Here $X \sim N(52, 144)$, so $\mu = 52$ and $\sigma = \sqrt{144} = 12$.
\begin{enumerate}
\item $z = \dfrac{40 - 52}{12} = \dfrac{-12}{12} = -1$.
\item $\mathbb{P}(X < 40) = \mathbb{P}(Z < -1) = 1 - \Phi(1) = 1 - 0.8413 = 0.1587$.
\item Standardising both bounds:
\[
z_1 = \frac{45-52}{12} = -0.5833\ldots, \qquad z_2 = \frac{65-52}{12} = 1.0833\ldots.
\]
Hence
\begin{align*}
\mathbb{P}(45 < X < 65) &= \Phi(1.0833\ldots) - \Phi(-0.5833\ldots)\\
&= \Phi(1.0833\ldots) - \bigl[1 - \Phi(0.5833\ldots)\bigr]\\
&\approx 0.8607 - (1 - 0.7202) = 0.8607 - 0.2798 = 0.5809.
\end{align*}
(Using $\Phi(1.08) = 0.8599$, $\Phi(0.58) = 0.7190$ gives $\approx 0.579$; answers around $0.58$ are accepted.)
\item We need the mark $m$ with $\mathbb{P}(X > m) = 0.10$, i.e. $\mathbb{P}(X < m) = 0.90$. From the tables, $\Phi(z) = 0.90$ gives $z = 1.2816$. De-standardising:
\[
m = \mu + z\sigma = 52 + 1.2816 \times 12 = 52 + 15.3792 \approx 67.4.
\]
The top $10\%$ of candidates score above approximately $67.4$ marks.
\end{enumerate}
\end{corrige}

\begin{corrige}[Exercise 5 — Finding parameters of a Normal distribution]
From $\mathbb{P}(X < 30) = 0.1587$: since $\Phi(1) = 0.8413$, we have $1 - \Phi(1) = 0.1587$, so the standardised value is $-1$:
\[
\frac{30 - \mu}{\sigma} = -1 \quad\Longrightarrow\quad \mu - \sigma = 30. \qquad (1)
\]
From $\mathbb{P}(X > 55) = 0.0228$: since $\Phi(2) = 0.9772$, we have $1 - \Phi(2) = 0.0228$, so the standardised value is $2$:
\[
\frac{55 - \mu}{\sigma} = 2 \quad\Longrightarrow\quad \mu + 2\sigma = 55. \qquad (2)
\]
Subtracting $(1)$ from $(2)$: $3\sigma = 25$, so $\sigma = \dfrac{25}{3} \approx 8.33$. Then from $(1)$:
\[
\mu = 30 + \sigma = 30 + \frac{25}{3} = \frac{115}{3} \approx 38.33.
\]
\[
\boxed{\mu = \tfrac{115}{3} \approx 38.3, \qquad \sigma = \tfrac{25}{3} \approx 8.33.}
\]
\textbf{Check:} $\dfrac{30 - 38.33}{8.33} = -1$ \checkmark\ and $\dfrac{55 - 38.33}{8.33} = 2$ \checkmark.
\end{corrige}

\begin{corrige}[Exercise 6 — Quality control: adjusting the standard deviation]
Let $M \sim N(505, 4^2)$.
\begin{enumerate}
\item Standardising the rejection limit $500$:
\[
z = \frac{500 - 505}{4} = -1.25.
\]
\[
\mathbb{P}(M < 500) = \mathbb{P}(Z < -1.25) = 1 - \Phi(1.25) = 1 - 0.8944 = 0.1056.
\]
So about $\mathbf{10.56\%}$ of packets are rejected.
\item We now require $\mathbb{P}(M < 500) = 0.01$ with $\mu = 505$ and unknown $\sigma$. From the tables, $\Phi(z) = 0.99$ gives $z = 2.3263$, so the standardised value of $500$ must be $-2.3263$:
\[
\frac{500 - 505}{\sigma} = -2.3263 \quad\Longrightarrow\quad \sigma = \frac{5}{2.3263} \approx 2.15\ \mathrm{g}.
\]
The machine must be improved so that the standard deviation is reduced from $4\ \mathrm{g}$ to about $2.15\ \mathrm{g}$.
\end{enumerate}
\end{corrige}

\begin{corrige}[Exercise 7 — Normal approximation to the Binomial with continuity correction]
\begin{enumerate}
\item $G \sim B(200,\; 0.85)$.
\item The conditions for a Normal approximation are $np > 5$ and $n(1-p) > 5$ (some texts use $npq > 5$). Here
\[
np = 200 \times 0.85 = 170 > 5, \qquad n(1-p) = 200 \times 0.15 = 30 > 5,
\]
so the approximation $G \approx N(np, npq)$ is appropriate, with
\[
\mu = 170, \qquad \sigma^2 = 200 \times 0.85 \times 0.15 = 25.5, \qquad \sigma = \sqrt{25.5} \approx 5.05.
\]
\item Applying the \cle{continuity correction}, $\mathbb{P}(G < 160) = \mathbb{P}(G \le 159)$ becomes $\mathbb{P}(Y < 159.5)$ where $Y \sim N(170, 25.5)$:
\[
z = \frac{159.5 - 170}{5.05} = \frac{-10.5}{5.05} \approx -2.08.
\]
\[
\mathbb{P}(G < 160) \approx \mathbb{P}(Z < -2.08) = 1 - \Phi(2.08) = 1 - 0.9812 = 0.0188.
\]
So $\mathbb{P}(G < 160) \approx 0.019$ (about $1.9\%$).
\end{enumerate}
\end{corrige}

\begin{corrige}[Exercise 8 — Normal approximation to the Poisson distribution]
Let $X \sim \mathrm{Po}(12)$.
\begin{enumerate}
\item A Poisson distribution $\mathrm{Po}(\lambda)$ can be approximated by $N(\lambda, \lambda)$ when $\lambda$ is sufficiently large — the usual rule of thumb is $\lambda > 10$ (some texts say $\lambda \ge 15$). Here $\lambda = 12 > 10$, so the approximation is acceptable (though only moderately accurate).
\item With the approximation $X \approx N(12, 12)$, we have $\sigma = \sqrt{12} \approx 3.464$. "At most $8$ accidents" means $\mathbb{P}(X \le 8)$; with the continuity correction this becomes $\mathbb{P}(Y < 8.5)$:
\[
z = \frac{8.5 - 12}{\sqrt{12}} = \frac{-3.5}{3.464} \approx -1.01.
\]
\[
\mathbb{P}(X \le 8) \approx \mathbb{P}(Z < -1.01) = 1 - \Phi(1.01) = 1 - 0.8438 = 0.1562.
\]
So the probability of at most $8$ accidents in a week is approximately $0.156$.
\end{enumerate}
\end{corrige}

\begin{corrige}[Exercise 9 — Combining independent Normal variables: filled sacks]
Let $B \sim N(50.2, 0.16)$ (bag) and $S \sim N(0.8, 0.04)$ (sack), independent, and $F = B + S$.

\textbf{(a)} For independent Normal variables, means add and variances add:
\[
\mathbb{E}(F) = 50.2 + 0.8 = 51.0\ \mathrm{kg}, \qquad \operatorname{Var}(F) = 0.16 + 0.04 = 0.20\ \mathrm{kg}^2,
\]
so $F \sim N(51.0, 0.20)$, with $\sigma_F = \sqrt{0.20} \approx 0.4472$.

\emph{Mark scheme:} M1 for stating that means add; M1 for variances add (independence quoted); A1 for both values correct. \textbf{[3 marks]}

\textbf{(b)}
\[
z = \frac{51.5 - 51.0}{\sqrt{0.20}} = \frac{0.5}{0.4472} \approx 1.12.
\]
\[
\mathbb{P}(F > 51.5) = 1 - \Phi(1.12) = 1 - 0.8686 = 0.1314.
\]

\emph{Mark scheme:} M1 standardising with their $\sigma_F$; A1 for $z = 1.12$; M1 for using $1 - \Phi$; A1 for $0.1314$ (accept $0.131$–$0.132$). \textbf{[4 marks]}

\textbf{(c)} Let $Y$ be the number of the $4$ sacks weighing more than $51.5$ kg. Then $Y \sim B(4, p)$ with $p = 0.1314$:
\[
\mathbb{P}(Y = 2) = \binom{4}{2} p^2 (1-p)^2 = 6 \times (0.1314)^2 \times (0.8686)^2.
\]
\[
= 6 \times 0.017266 \times 0.754466 \approx 0.0781.
\]

\emph{Mark scheme:} B1 for identifying the binomial model $B(4, 0.1314)$; M1 for the binomial formula with $n=4$, $k=2$; A1 for $\binom{4}{2} = 6$ used; M1 correct substitution; A1 for $\approx 0.078$. \textbf{[5 marks]}

\emph{Examiner's expectations:} candidates must state the distribution of $F$ explicitly, quote independence when adding variances, and carry their answer from (b) into (c) (error carried forward applies).
\end{corrige}

\begin{corrige}[Exercise 10 — Two athletes: difference of Normal variables]
\textbf{(a)} For independent Normal variables, $D = T_A - T_B$ is Normal with
\[
\mathbb{E}(D) = 11.2 - 11.5 = -0.3, \qquad \operatorname{Var}(D) = 0.09 + 0.16 = 0.25,
\]
(variances \emph{add} even for a difference), so
\[
D \sim N(-0.3,\; 0.25), \qquad \sigma_D = 0.5.
\]

\emph{Mark scheme:} M1 for subtracting means; M1 for adding variances; A1 for $N(-0.3, 0.25)$. \textbf{[3 marks]}

\textbf{(b)} A beats B when $T_A < T_B$, i.e. $D < 0$:
\[
z = \frac{0 - (-0.3)}{0.5} = 0.6.
\]
\[
\mathbb{P}(D < 0) = \Phi(0.6) = 0.7257.
\]
Athlete A wins with probability about $0.726$.

\emph{Mark scheme:} B1 for translating "A wins" into $D < 0$; M1 standardising; A1 for $z = 0.6$; A1 for $0.7257$. \textbf{[4 marks]}

\textbf{(c)} Let $W$ be the number of races won by A out of $5$: $W \sim B(5, 0.7257)$.
\begin{align*}
\mathbb{P}(W \ge 3) &= \mathbb{P}(W=3) + \mathbb{P}(W=4) + \mathbb{P}(W=5)\\
&= \binom{5}{3}p^3q^2 + \binom{5}{4}p^4q + p^5, \quad p = 0.7257,\ q = 0.2743\\
&= 10(0.3822)(0.07524) + 5(0.2773)(0.2743) + 0.2012\\
&\approx 0.2876 + 0.3803 + 0.2012 = 0.8691.
\end{align*}
So $\mathbb{P}(W \ge 3) \approx 0.869$.

\emph{Mark scheme:} B1 for the binomial model; M1 for summing the three correct terms; A1 for $\approx 0.869$. \textbf{[3 marks]}

\emph{Common trap:} subtracting the variances in (a) — the variance of a difference of independent variables is always the \emph{sum} of the variances.
\end{corrige}

\begin{corrige}[Exercise 11 — Determining unknown parameters and using approximations]
\textbf{(a)} We have $L \sim N(\mu, \sigma^2)$.

$\mathbb{P}(L < 98.0) = 0.025$: from the tables, $\Phi(z) = 0.975$ gives $z = 1.96$, so the standardised value of $98.0$ is $-1.96$:
\[
\frac{98.0 - \mu}{\sigma} = -1.96 \quad\Longrightarrow\quad \mu - 1.96\sigma = 98.0. \qquad (1)
\]
$\mathbb{P}(L > 102.0) = 0.05$: from the tables, $\Phi(z) = 0.95$ gives $z = 1.6449$:
\[
\frac{102.0 - \mu}{\sigma} = 1.6449 \quad\Longrightarrow\quad \mu + 1.6449\sigma = 102.0. \qquad (2)
\]
Subtracting $(1)$ from $(2)$:
\[
(1.6449 + 1.96)\sigma = 4.0 \quad\Longrightarrow\quad 3.6049\,\sigma = 4.0 \quad\Longrightarrow\quad \sigma \approx 1.110\ \mathrm{cm}.
\]
Substituting in $(1)$:
\[
\mu = 98.0 + 1.96 \times 1.110 = 98.0 + 2.175 \approx 100.17\ \mathrm{cm}.
\]
\[
\boxed{\mu \approx 100.2\ \mathrm{cm}, \qquad \sigma \approx 1.11\ \mathrm{cm}.}
\]

\emph{Mark scheme:} B1 for $z = \pm 1.96$; B1 for $z = 1.6449$ (accept $1.645$); M1 for setting up both standardised equations; M1 for solving the simultaneous system; A1 for $\sigma$; A1 for $\mu$. \textbf{[6 marks]}

\textbf{(b)} First find the probability $p$ that one rod has length between $98.0$ and $102.0$ cm:
\[
p = 1 - 0.025 - 0.05 = 0.925.
\]
Let $X$ be the number of rods (out of $50$) with acceptable length: $X \sim B(50, 0.925)$. The box is accepted if $X \ge 48$.

Since $np = 46.25 > 5$ and $n(1-p) = 3.75 < 5$, the direct Normal approximation to $X$ is poor. Instead, consider the number of \emph{defective} rods $Y = 50 - X \sim B(50, 0.075)$. Here $np = 3.75$ is small, $p$ is small and $n$ large, so a \textbf{Poisson approximation} is suitable:
\[
Y \approx \mathrm{Po}(3.75).
\]
The box is accepted when $X \ge 48$, i.e. $Y \le 2$:
\begin{align*}
\mathbb{P}(Y \le 2) &\approx e^{-3.75}\left(1 + 3.75 + \frac{3.75^2}{2}\right)\\
&= e^{-3.75}\left(1 + 3.75 + 7.03125\right)\\
&= 0.02352 \times 11.78125 \approx 0.2771.
\end{align*}
So the probability that a box is accepted is approximately $0.277$.

\emph{Mark scheme:} B1 for $p = 0.925$; B1 for identifying $X \sim B(50, 0.925)$; M1 for checking the Normal-approximation conditions and rejecting it ($nq < 5$); M1 for switching to defective rods $Y \sim B(50, 0.075)$; B1 for the Poisson approximation $\mathrm{Po}(3.75)$ with justification; M1 for translating acceptance into $Y \le 2$; M1 for the Poisson cumulative sum; A1 for $\approx 0.277$. \textbf{[8 marks]}

\emph{Examiner's expectations:} the key insight is recognising that the Normal approximation fails here ($nq = 3.75 < 5$) and choosing the Poisson approximation to the binomial instead. Candidates who apply the Normal approximation blindly (even with continuity correction) score only the first two marks. Note also that no continuity correction is needed for the Poisson calculation since it remains a discrete distribution.
\end{corrige}

\newpage

\coursec{Summary sheet}

\begin{popfigure}
\begin{tikzpicture}[
  every node/.style={draw, rounded corners, fill=popblueL, text width=2.2cm, align=center, font=\small},
  central/.style={fill=popblue, text width=3cm, font=\large\bfseries, text=white},
  branch1/.style={fill=poporange},
  branch2/.style={fill=popgreen},
  branch3/.style={fill=poppurple},
  branch4/.style={fill=poppink},
  branch5/.style={fill=popgold},
  branch6/.style={fill=popblue},
  branch7/.style={fill=popmuted},
]
\node[central] (center) {Special Continuous Distributions};
\node[branch1] (u) at (3.5,2.5) {Uniform (Rectangular)};
\node[branch2] (e) at (3.5,1) {Exponential};
\node[branch3] (n) at (3.5,-0.5) {Normal};
\node[branch4] (s) at (3.5,-2) {Standardisation \& Tables};
\node[branch5] (a) at (0,-3.5) {Approximations \& Combinations};
\node[branch6] (f) at (-3.5,-2) {Finding $\mu$ \& $\sigma$};
\node[branch7] (c) at (-3.5,1) {Continuity Correction};
\foreach \n in {u,e,n,s,a,f,c}
  \draw[popline, thick] (center) -- (\n);
\end{tikzpicture}
\legende{Mind map of Special Continuous Distributions}
\end{popfigure}

\begin{bilan}
\textbf{Key Definitions}
\begin{itemize}
\item \cle{Uniform (Rectangular) Distribution}: Continuous random variable $X$ on $[a,b]$ with pdf $f(x)=\frac{1}{b-a}$ for $a\le x\le b$, zero elsewhere. $E(X)=\frac{a+b}{2}$, $\mathrm{Var}(X)=\frac{(b-a)^2}{12}$.
\item \cle{Exponential Distribution}: Continuous random variable $X\ge0$ with pdf $f(x)=\lambda e^{-\lambda x}$ ($x\ge0$), $\lambda>0$. $E(X)=\frac{1}{\lambda}$, $\mathrm{Var}(X)=\frac{1}{\lambda^2}$. Memoryless property: $P(X>s+t\mid X>s)=P(X>t)$.
\item \cle{Normal Distribution}: $X\sim N(\mu,\sigma^2)$ with pdf $f(x)=\frac{1}{\sigma\sqrt{2\pi}}e^{-\frac{(x-\mu)^2}{2\sigma^2}}$. Standard normal $Z\sim N(0,1)$ with cdf $\Phi(z)$.
\item \cle{Standardisation}: $Z=\frac{X-\mu}{\sigma}$ transforms any $N(\mu,\sigma^2)$ to $N(0,1)$.
\item \cle{Continuity Correction}: When approximating discrete (Binomial/Poisson) by Normal, add/subtract $0.5$ to the discrete bound.
\end{itemize}

\textbf{Laws, Formulas \& Theorems to Know}
\begin{itemize}
\item Uniform: $P(c\le X\le d)=\frac{d-c}{b-a}$ for $a\le c<d\le b$.
\item Exponential: $P(X>x)=e^{-\lambda x}$, $P(X\le x)=1-e^{-\lambda x}$.
\item Link Exponential–Poisson: If events occur as Poisson process with rate $\lambda$, inter-arrival times are Exponential($\lambda$).
\item Normal: $\Phi(-z)=1-\Phi(z)$; $\Phi(0)=0.5$; $\Phi(1.96)\approx0.975$, $\Phi(2.576)\approx0.995$.
\item Central intervals: $P(\mu-1.96\sigma<X<\mu+1.96\sigma)=0.95$; $P(\mu-2.576\sigma<X<\mu+2.576\sigma)=0.99$.
\item Linear combination: If $X_i\sim N(\mu_i,\sigma_i^2)$ independent, then $\sum a_i X_i \sim N(\sum a_i\mu_i,\sum a_i^2\sigma_i^2)$.
\item Normal approximation to Binomial: $X\sim B(n,p) \approx N(np,npq)$ if $np>5$ and $nq>5$.
\item Normal approximation to Poisson: $X\sim Po(\lambda) \approx N(\lambda,\lambda)$ if $\lambda>10$ (or $>15$).
\end{itemize}

\textbf{Methods}
\begin{enumerate}
\item \textbf{Uniform}: Identify $a,b$; compute probabilities as lengths; use $E(X),\mathrm{Var}(X)$ formulas.
\item \textbf{Exponential}: Identify $\lambda$ (rate) or mean $1/\lambda$; use cdf $1-e^{-\lambda x}$ or survival $e^{-\lambda x}$; apply memoryless property when needed.
\item \textbf{Normal (standard)}: Use tables to find $\Phi(z)$ or $z$ given probability; exploit symmetry $\Phi(-z)=1-\Phi(z)$.
\item \textbf{Any Normal}: Standardise $z=\frac{x-\mu}{\sigma}$; use tables; de-standardise $x=\mu+z\sigma$.
\item \textbf{Finding $\mu$ and/or $\sigma$}: Set up equations from given probabilities (e.g. $P(X<x_1)=p_1$, $P(X>x_2)=p_2$); standardise to get $z$-values; solve simultaneous equations.
\item \textbf{Continuity correction}: For $P(X\le k)$ use $P(X<k+0.5)$; for $P(X\ge k)$ use $P(X>k-0.5)$; for $P(X=k)$ use $P(k-0.5<X<k+0.5)$.
\item \textbf{Independent Normals}: For $X,Y$ independent normals, $X\pm Y$ is normal with mean $\mu_X\pm\mu_Y$ and variance $\sigma_X^2+\sigma_Y^2$.
\end{enumerate}

\textbf{Common Mistakes}
\begin{itemize}
\item Forgetting that the uniform pdf is zero outside $[a,b]$; integrating over wrong limits.
\item Confusing rate $\lambda$ and mean $1/\lambda$ in exponential; using $\lambda$ as mean.
\item Misreading normal tables: using $\Phi(z)$ as tail probability instead of cumulative.
\item Forgetting symmetry: $\Phi(-z)=1-\Phi(z)$.
\item Omitting continuity correction when approximating discrete distributions.
\item Using normal approximation when conditions ($np>5$, $nq>5$ or $\lambda>10$) are not met.
\item Incorrectly adding variances for independent normals: variances add, standard deviations do not.
\item Mixing up $z$-values for central 95\% (1.96) and 99\% (2.576).
\end{itemize}

\textbf{GCE Tips}
\begin{itemize}
\item Always define the distribution and its parameters clearly at the start of a solution.
\item Show the standardisation step explicitly: $z=\frac{x-\mu}{\sigma}$.
\item When using tables, state the value you look up (e.g. $\Phi(1.645)=0.95$).
\item For inverse problems (find $\mu,\sigma$), write the equations in terms of $z$-scores before solving.
\item In hypothesis testing context (though not in this chapter), remember the critical values 1.645, 1.96, 2.576.
\item Check that probabilities sum to 1 or lie in $[0,1]$ as a sanity check.
\item For exponential, the memoryless property is a favourite exam question; recognise wording like ``given that it has lasted $s$ units, probability it lasts $t$ more''.
\item When approximating binomial/poisson, state the conditions and apply continuity correction; lose marks if omitted.
\item For combinations of independent normals, write the resulting mean and variance clearly.
\end{itemize}
\end{bilan}

\findecours

\end{document}
